\chapter{The Griess algebra of the moonshine module and AxiomMath's Monster}\label{ch:21}

\providecommand{\chXXIB}{\mathcal{B}}
\providecommand{\chXXIC}{\mathcal{C}}
\providecommand{\chXXIS}{\mathcal{S}}
\providecommand{\chXXIL}{\bar{\Lambda}}
\providecommand{\chXXIF}{\mathbb{F}_2}
\providecommand{\chXXISym}{\mathrm{Sym}}
\providecommand{\chXXIT}{\mathcal{T}}

This chapter documents the release of the library named in the front matter; every Lean quotation is checked by
\texttt{scripts/textbook/sync.py} to occur verbatim (up to whitespace and comments) in the source files
of that release.  It documents $65$ files:
the $44$ files of \leanpath{VOA/\allowbreak Moonshine/\allowbreak MonsterGriess/}, which are a port of
part of the AxiomMath CFSG library; the $15$ files
\leanpath{VOA/\allowbreak Moonshine/\allowbreak MoonshineCNGriess*.lean}; the $5$ files
\leanpath{VOA/\allowbreak Moonshine/\allowbreak MoonshineCNMonsterAction*.lean}; and
\leanpath{VOA/\allowbreak Moonshine/\allowbreak MoonshineLRGriessTransfer.lean}.

\section{Overview}\label{ch21:sec:intro}

\subsection{The central theorem}

The chapter has one central theorem; everything else prepares it or uses it.  Write $\chXXIB$ for the
$196884$-dimensional \emph{rational Griess algebra} of AxiomMath's Lean library: a commutative,
non-associative algebra over $\Q$ with product $u\cdot v$ and a positive definite bilinear form $b$
(Definitions~\ref{ch21:def:spaceB} and~\ref{ch21:def:gmul}).  Write $\Vnat_2$ for the weight-two space of
the moonshine module constructed in Chapters~\ref{ch:02}--\ref{ch:14}, with the product
$u\times v=u_{(1)}v$ and the invariant bilinear form $\langle\cdot,\cdot\rangle$ of $\Vnat$, normalised by
$\langle\mathbb1,\mathbb1\rangle=1$.

\begin{theorem}[Griess isomorphism; proved as Theorem~\ref{ch21:thm:griessIso}]\label{ch21:thm:intro}
There is a $\C$-linear isomorphism $\Theta:\C\otimes_\Q\chXXIB\to\Vnat_2$ such that
$\Theta(x\cdot y)=\Theta(x)\times\Theta(y)$ and $\langle\Theta x,\Theta y\rangle=2\,b(x,y)$ for all
$x,y\in\C\otimes_\Q\chXXIB$ (product and form extended $\C$-bilinearly).
\end{theorem}
\begin{leancode}
noncomputable def griessIso : ℂ ⊗[ℚ] SpaceB ≃ₗ[ℂ] VnatVOA.grading 2 :=
  (Algebra.TensorProduct.basis ℂ griessBasisQ).equiv griessBasisC (Equiv.refl _)
\end{leancode}
\begin{leancode}
theorem griessIso_mul (x y : ℂ ⊗[ℚ] SpaceB) :
    griessIso (LinearMap.BilinMap.baseChange ℂ gmul x y)
      = VnatVOA.griessMul (griessIso x) (griessIso y) := by
\end{leancode}
\begin{leancode}
theorem griessIso_form (x y : ℂ ⊗[ℚ] SpaceB) :
    SimpleNat.natForm (griessIso x : Vnat) (griessIso y : Vnat)
      = 2 * LinearMap.BilinForm.baseChange ℂ bform x y := by
\end{leancode}
(\leanpath{VOA/\allowbreak Moonshine/\allowbreak MoonshineCNGriessIso.lean}, namespace
\lean{VOA.Moonshine.GriessDictionary}.)

\paragraph{Symbol key.}  \lean{SpaceB} is $\chXXIB$, \lean{gmul} its product and \lean{bform} its form
(\S\ref{ch21:sec:spaceB}, verbatim from AxiomMath).  \lean{ℂ ⊗[ℚ] SpaceB} is $\C\otimes_\Q\chXXIB$ and
\lean{≃ₗ[ℂ]} a $\C$-linear bijection.  \lean{VnatVOA} is the moonshine module of Chapter~\ref{ch:14}
(named in Chapter~\ref{ch:20}), \lean{VnatVOA.grading 2} its weight-two space $\Vnat_2$,
\lean{VnatVOA.griessMul} the product $u_{(1)}v$ on it, and \lean{SimpleNat.natForm} its invariant form
(Chapter~\ref{ch:10}).  \lean{LinearMap.BilinMap.baseChange ℂ gmul} and
\lean{LinearMap.BilinForm.baseChange ℂ bform} are the $\C$-bilinear extensions of product and form.
\lean{griessBasisQ}, \lean{griessBasisC} are bases used to build $\Theta$; the theorem
\lean{griessIso_tmul} says $\Theta(1\otimes u)$ is an explicit vertex-algebra state
(\S\ref{ch21:sec:iso}).  None of the three statements has a hypothesis.

\paragraph{Why the factor $2$.}  The two forms are normalised independently: $\Vnat$'s by
$\langle\mathbb1,\mathbb1\rangle=1$, which forces $\langle\omega,\omega\rangle=c/2=12$ for the conformal
vector $\omega$; AxiomMath's $b$ is the sum of the standard dot products of the coordinates.  The
identity element of $\chXXIB$ is $e=(\tfrac14I,0,0)$, with $b(e,e)=\tr(I/16)=\tfrac32$; it corresponds to
the identity $\omega/2$ of $\Vnat_2$ (both facts are Lean theorems, Remark~\ref{ch21:rem:unit}), and
$\langle\omega/2,\omega/2\rangle=3=2\cdot\tfrac32$.  Since a linear map preserves $b$ iff it preserves $2b$,
the factor is harmless for everything that follows.

\paragraph{Why this is the central fidelity check of the Monster part.}  The Monster is not defined
anywhere in LeanVOA.  The group that later statements call ``the Monster'' is AxiomMath's \lean{AutO}:
the group of $\Q$-linear automorphisms of $\chXXIB$ preserving product and form, transcribed verbatim
from AxiomMath's Lean source.  ``Every element of the Monster acts on $\Vnat$'' (Chapter~\ref{ch:22})
means something only if the object on which \lean{AutO} acts is visibly an object inside $\Vnat$, and
$\Theta$ is that identification.  A referee who accepts (i) that \lean{SpaceB}, \lean{gmul},
\lean{bform} are literally AxiomMath's definitions (\S\ref{ch21:sec:fidelityaudit}), (ii) that
\lean{VnatVOA} is the moonshine module of Chapter~\ref{ch:14}, and (iii) the three statements above, has
accepted that ``\lean{AutO} acts on the Griess algebra of \lean{VnatVOA}'' is a precise statement.

\paragraph{What this chapter does \emph{not} contain.}  The library does not prove that \lean{AutO}
is finite, that it has the order $2^{46}\cdot3^{20}\cdots71$ of the Monster, or that it is simple;
AxiomMath's library proves these facts in Lean~v4.35, and they were not ported
(\S\ref{ch21:sec:notported}).  Nor does it prove AxiomMath's \emph{generation theorem}, which says that
\lean{AutO} is generated by an explicit subgroup $G_{\mathrm{vis}}$ (classically
$\cong2^{1+24}.\mathrm{Co}_1$; this identification is not formalised) and one explicit element $\tau_B$
of order $3$; that theorem enters Chapter~\ref{ch:22} as the single named hypothesis \lean{GmatEqAutO}.
The lifting of $G_{\mathrm{vis}}$ and $\tau_B$ to automorphisms of $\Vnat$ and the converse inclusion
$\Aut(\Vnat)\hookrightarrow$ \lean{AutO} are Chapter~\ref{ch:22}.  This chapter supplies the stage: the
port, the identification $\Theta$, and the lift \emph{framework}.

\subsection{The Griess algebra in one page}

Let $V=\bigoplus_{n\ge0}V_n$ be a vertex operator algebra of CFT type ($V_0=\C\mathbb 1$) with
$V_1=0$.  For $u,v\in V_2$ the modes $u_{(n)}v$ have weight $2+2-n-1=3-n$, so $u_{(1)}v\in V_2$ and
$u_{(3)}v\in V_0=\C\mathbb1$.  The \emph{Griess algebra} of $V$ is $V_2$ with the product
$u\times v=u_{(1)}v$ and the form $\langle u,v\rangle$ defined by $u_{(3)}v=\langle u,v\rangle\mathbb1$.
Skew-symmetry $Y(u,z)v=e^{zL_{-1}}Y(v,-z)u$ gives
$u_{(1)}v=v_{(1)}u-L_{-1}(v_{(2)}u)+\cdots$, and $v_{(2)}u\in V_1=0$, so the product is
\emph{commutative}; the invariant bilinear form of a VOA of CFT type with $L_1V_1=0$ makes the form
\emph{associative}, $\langle u\times v,w\rangle=\langle u,v\times w\rangle$.  So $V_2$ is a
commutative, non-associative algebra with an associative form.  The conformal vector $\omega$ lies
in $V_2$ and acts by $\omega\times v=L_0v=2v$: the element $\omega/2$ is the identity.

For $V=\Vnat$ one has $\dim\Vnat_2=196884$ (Chapter~\ref{ch:05}), and Frenkel, Lepowsky and Meurman
proved that $\Vnat_2$ with this product is (a variant of) the algebra that Griess had constructed by
hand to build the Monster.  The decomposition
\begin{equation}\label{ch21:eq:vnat2}
  \Vnat_2=(V_\Leech^+)_2\oplus(T^+)_2,\qquad
  (V_\Leech^+)_2=S^2\mathfrak h\;\oplus\;\operatorname{span}\{x_\alpha:\langle\alpha,\alpha\rangle=4\},\qquad
  (T^+)_2=\mathfrak h(-\tfrac12)\otimes T,
\end{equation}
with dimensions $196884=(300+98280)+24\cdot4096$, is visible in Conway's description of the algebra:
\begin{equation}\label{ch21:eq:conwayB}
  \chXXIB=\chXXISym\;\oplus\;X\;\oplus\;W,\qquad
  \chXXISym=S^2(\Q^{24}),\quad X=\Q^{98280},\quad W=\Q^{4096}\otimes\Q^{24},
\end{equation}
$X$ indexed by the $98280$ pairs $\pm\alpha$ of minimal vectors of the Leech lattice and $W$ by a
$4096$-dimensional representation of the extraspecial group $2^{1+24}$ tensored with $\Q^{24}$.
AxiomMath implements \eqref{ch21:eq:conwayB} with rational coordinates; LeanVOA ports that
implementation and proves that, after extension of scalars, it \emph{is} \eqref{ch21:eq:vnat2}.

\begin{physics}
In two-dimensional conformal field theory the states of $V_2$ are the \emph{dimension-two chiral
fields}.  When $V_1=0$ (no spin-one currents, i.e.\ no continuous symmetry) the operator product
expansion of two such fields reads
\begin{equation*}
  a(z)\,b(w)\;\sim\;\frac{\langle a,b\rangle}{(z-w)^4}\;+\;\frac{(a\times b)(w)}{(z-w)^2}
  \;+\;\frac{\tfrac12\partial(a\times b)(w)}{z-w},
\end{equation*}
since a $(z-w)^{-3}$ term would be a spin-one field.  So the Griess algebra is ``the OPE of the
$196884$ spin-two fields of the Monster CFT'': the product is the coefficient of $(z-w)^{-2}$, the form
that of $(z-w)^{-4}$.  The stress tensor $T(z)$ is one of these fields ($\omega$), and
$T(z)T(w)\sim\frac{c/2}{(z-w)^4}+\frac{2T(w)}{(z-w)^2}+\cdots$ with $c=24$ says $\omega\times\omega=2\omega$,
$\langle\omega,\omega\rangle=12$.  The orbifold picture: the $\Z/2$-orbifold of $24$ bosons on the Leech
torus has untwisted spin-two fields (the $300$ bilinears $\partial X^i\partial X^j$ and the $98280$
combinations $e^{i\alpha X}+e^{-i\alpha X}$, $\langle\alpha,\alpha\rangle=4$, which we abbreviate as
``$\cos(\alpha X)$'' in the physics boxes; in the VOA they are the states $x_\alpha=e^\alpha+e^{-\alpha}$) and
twisted ones: the twisted
ground states $\sigma_s$ ($4096$ of them) have weight $24/16=\tfrac32$, and one half-integral oscillator
$\partial X^i_{-1/2}$ raises them to weight $2$, giving $24\cdot4096$ more.
\end{physics}

\subsection{The Monster, Tits's theorem, and the role of the form}\label{ch21:sec:tits}

The Monster $\mathbb M$ is the largest sporadic simple group, of order
$2^{46}3^{20}5^{9}7^{6}11^2 13^3\cdot17\cdot19\cdot23\cdot29\cdot31\cdot41\cdot47\cdot59\cdot71$.  Griess
constructed it as a group of automorphisms of a commutative non-associative algebra of dimension
$196883$ with an invariant form; Conway's simplified construction adjoins an identity element, giving
$196884=1+196883$ dimensions, and that is the algebra AxiomMath implements.  The two versions carry the
same automorphisms: an algebra automorphism of the $196884$-dimensional algebra fixes its identity $e$
(the identity of an algebra is unique), so if it also preserves $b$ it preserves the orthogonal
complement $e^\perp$, which with the product $u\cdot v$ followed by orthogonal projection to $e^\perp$ is
Griess's $196883$-dimensional algebra (up to normalisation); conversely, since
$u\cdot v=\mathrm{pr}_{e^\perp}(u\cdot v)+\tfrac{b(u,v)}{b(e,e)}e$ for $u,v\in e^\perp$, an automorphism of
Griess's algebra preserving its form extends by $e\mapsto e$.  Tits proved that the full
automorphism group of the algebra is exactly $\mathbb M$.

\emph{Why $b$ appears in \lean{AutO}.}  The form is the natural extra datum: $b$ is associative and
$b(u,v)=b(u\cdot v,e)$, so $b$ is determined by the product together with the linear functional
$b(\cdot,e)$; classically the associative form is unique up to scalars on each of the two irreducible
summands $\Q e$ and $e^\perp$ of the Monster, so imposing $b$ costs nothing classically.  But that is a
theorem about the Monster, not a consequence of the definition, and AxiomMath's definition
(Definition~\ref{ch21:def:AutO}) simply requires both conditions.  We follow it verbatim; LeanVOA proves
neither that $b$-preservation is automatic nor Tits's theorem.  ``\lean{AutO} is the Monster'' is
therefore a statement \emph{by reference} (to AxiomMath's library and to the literature).

Frenkel, Lepowsky and Meurman constructed $\Vnat$ together with a faithful action of $\mathbb M$ on it,
extending Griess's action on $\Vnat_2$; since $\Vnat$ is generated by $\Vnat_2$, Tits's theorem then gives
$\Aut(\Vnat)=\mathbb M$.

\begin{physics}
$\Vnat$ is the chiral algebra of the holomorphic conformal field theory with $c=24$ obtained as the
$\Z/2$-orbifold of $24$ free chiral bosons compactified on the Leech torus $\mathbb{R}^{24}/\Leech$ (Dixon,
Ginsparg and Harvey).  Its partition function is $J(q)=j(q)-744$, and classically its global symmetry
group is the Monster.  Each symmetry $g$ gives a twisted partition function $T_g(q)=\tr(g\,q^{L_0-1})$;
Conway--Norton's conjecture (Chapter~\ref{ch:16}) says these are genus-zero Hauptmoduln.  This chapter
makes ``the Monster'' a concrete Lean object acting on the spin-two fields of this CFT.
\end{physics}

\subsection{Provenance}\label{ch21:sec:provenance}

The $44$ files under \leanpath{VOA/\allowbreak Moonshine/\allowbreak MonsterGriess/} are ports of files
of the Lean library \emph{AxiomMath/CFSG}, at the pinned commit
\texttt{69cb972f587a3ecc035e092253167ff2e303c962}, released under the Apache~2.0 licence; its authors
are listed in \S\ref{ch21:sec:notes}.  Every ported file carries the Apache header naming the source
files, the commit and the original authors, and lists its modifications in a section ``Port notes'' of
its module docstring.  The port to Lean~4.30.0 / Mathlib~v4.30.0 is by F.~K.~Popov (with Claude).
The source library is written for Lean~v4.35.0-rc3; a port is therefore a \emph{re-elaboration} of the
source text by a different toolchain, not an import.  It contains the ``definitional closure'' of
AxiomMath's \lean{AutO} and \lean{Gmat} together with a ``light toolkit'' of theorems (the
parity-centralizer theorem, the triality element and its product invariance, the lattice facts they
need).  A declaration-by-declaration comparison with the source (\S\ref{ch21:sec:fidelityaudit}, after
all the definitions involved have been presented) found all definitions identical up to renaming of
namespaces, except one, the coordinate map \lean{phi0}, which is a definitionally equal restatement.

\subsection{How the chapter is organised}

Part~A (\S\S\ref{ch21:sec:golay}--\ref{ch21:sec:fidelityaudit}) presents AxiomMath's model: the
combinatorial data, the space $\chXXIB$, its product and form, the group \lean{AutO}, the visible
subgroup, the triality, the parity-centralizer theorem, what is not ported, and the port audit.  Part~B
(\S\S\ref{ch21:sec:vnatside}--\ref{ch21:sec:iso}) constructs the dictionary between $\chXXIB$ and
$\Vnat_2$ and proves the isomorphism.  Part~C (\S\S\ref{ch21:sec:action}--\ref{ch21:sec:transfer})
builds the framework in which elements of \lean{AutO} act on $\Vnat$.  Section~\ref{ch21:sec:filemap}
lists every file with its role, \S\ref{ch21:sec:audits} the audit files, and the chapter ends
with the referee's checklist.

\begin{remark}[Notation]\label{ch21:rem:notation}
$\chXXIF=\Z/2$ (Lean \lean{ZMod 2}).  $\chXXIC\subset\chXXIF^{24}$ is AxiomMath's Golay code
(\lean{Code.C}), $\chXXIC^*=\chXXIF^{24}/\chXXIC$ its cocode (\lean{Cocode}), $\chXXIC_0$ the
index-two subcode \lean{C₀}.  $\Leech$ is the Leech lattice, always in Conway's integer coordinates
scaled by $\sqrt8$: a lattice vector is an integer vector $u\in\Z^{24}$ and the inner product is
$u\cdot v/8$; $\chXXIL=\Leech/2\Leech$.  $\chXXIS$ (\lean{ShortIdx}) is the set of $98280$ classes of
$\chXXIL$ containing a minimal vector.  $I=\chXXIF\times\chXXIC_0$ (\lean{SpinIdx}) is the
$4096$-element spin index set.  $\chXXIB=\chXXISym\oplus X\oplus W$ with $\chXXISym$ the symmetric
$24\times24$ rational matrices, $X=\Q^{\chXXIS}$, $W=\Q^{I\times24}$.  $\chXXIT$ is the trilinear form of
$\chXXIB$; $T$ (roman) is always the $4096$-dimensional module of Chapter~\ref{ch:06}.  On the $\Vnat$
side, $\pi$ is the coordinate permutation \lean{piPerm} relating our Leech lattice to AxiomMath's.  As in
Chapter~\ref{ch:01}, \lean{V.smul u n v} is $u_{(n)}v$, so \lean{VnatVOA.smul u 1 v} is the Griess
product.  Status boxes cite the audit file in which the declaration's \texttt{\#print axioms} line
appears; every such line is expected to print \lean{[propext, Classical.choice, Quot.sound]} or a subset
(\S\ref{ch21:sec:audits}).
\end{remark}

%======================================================================
\section{AxiomMath's Golay code, its theta cocycle and the cocode}\label{ch21:sec:golay}

Part~A begins.  Throughout Part~A the Lean names live in the namespace
\lean{VOA.Moonshine.MonsterGriess} or one of its sub-namespaces (\lean{.Code}, \lean{.Spin},
\lean{.SpinLattice}, \lean{.Visible}, \lean{.ScalarKernel}, \lean{.Parity}, \lean{.Generated}),
which replace AxiomMath's \lean{CFSG.Monster}, \lean{CFSG.Monster.Griess}, \dots\ one for one.

\subsection{The binary Golay code in its MOG/hexacode form}

Chapter~\ref{ch:02} built the Golay code $\mathcal G$ as the row space of $[I_{12}\mid B]$ with a
Paley-type matrix $B$ (\lean{Code.golay}).  AxiomMath uses instead the description of Conway and
Curtis through the \emph{Miracle Octad Generator} (MOG) and the \emph{hexacode}.  The $24$
coordinates are arranged in a $4\times6$ array: coordinate $i=r+4j$ sits in row $r\in\{0,1,2,3\}$
and column $j\in\{0,\dots,5\}$.  Let $\mathbb F_4=\{0,1,\omega,\bar\omega\}$.  The hexacode
$H_6\subset\mathbb F_4^6$ is the $3$-dimensional $\mathbb F_4$-span of
$v_1=(1,0,0,1,\bar\omega,\omega)$, $v_2=(0,1,0,1,\omega,\bar\omega)$, $v_3=(0,0,1,1,1,1)$.  To a binary
word $x$ attach its \emph{hexacode projection} $h(x)\in\mathbb F_4^6$, whose $j$-th entry is
$x_{1,j}+\omega x_{2,j}+\bar\omega x_{3,j}$ (the rows $1,2,3$ are ``scored'' $1,\omega,\bar\omega$,
row $0$ scores $0$), and its \emph{top parity} $t(x)=\sum_jx_{0,j}$.

\begin{definition}[AxiomMath's Golay code]\label{ch21:def:C}
$\chXXIC=\{x\in\chXXIF^{24}: h(x)\in H_6\ \text{and every column has parity}\ t(x)\}$.
\end{definition}
\begin{leancode}
def hexProj (x : Fin 24 → ZMod 2) : Fin 6 → F4 := fun j =>
  algebraMap (ZMod 2) F4 (x (coord 1 j)) + α * algebraMap (ZMod 2) F4 (x (coord 2 j)) +
    ᾱ * algebraMap (ZMod 2) F4 (x (coord 3 j))

def topParity (x : Fin 24 → ZMod 2) : ZMod 2 := ∑ j, x (coord 0 j)
\end{leancode}
\begin{leancode}
def C : Submodule (ZMod 2) (Fin 24 → ZMod 2) where
  carrier := {x | hexProj x ∈ H₆ ∧ ∀ j, ∑ r, x (coord r j) = topParity x}
\end{leancode}
(\leanpath{VOA/\allowbreak Moonshine/\allowbreak MonsterGriess/\allowbreak Golay.lean}, \lean{VOA.Moonshine.MonsterGriess.Code.C}.)

\begin{dictionary}
\lean{F4} is a hand-made field structure on pairs $(\mathrm{re},\mathrm{im})\in\chXXIF^2$, with
\lean{α} $=(0,1)=\omega$ and \lean{ᾱ} $=(1,1)=\bar\omega$; the file proves that it is a field.
\lean{coord r j} is the coordinate $r+4j$, i.e.\ row $r$, column $j$.  The vectors \lean{v₁},
\lean{v₂}, \lean{v₃} and the span \lean{H₆} are defined in the same file exactly as in the text
(vector literals \lean{![1, 0, 0, 1, ᾱ, α]}, \dots).  The predicate is exactly the two MOG conditions.  The file constructs a basis \lean{basisC} of $12$
words and proves, by a kernel computation over the $4096$ codewords (\lean{decide +kernel}, no
\lean{native_decide}), that every codeword has weight in $\{0,8,12,16,24\}$ (\lean{four_dvd_wt}),
and that $\chXXIC$ is self-orthogonal (\lean{dotProduct_eq_zero});
\lean{CodeCocycle.lean} adds self-duality \lean{mem_C_iff_forall_dotProduct}.

\textbf{Could it be the wrong code?}  The identification with our $\mathcal G$ is a Lean theorem
(Theorem~\ref{ch21:thm:golaymatch} below), so $\chXXIC$ is a Golay code; there is no room for a
typo in $v_1,v_2,v_3$ that would survive that theorem.
\end{dictionary}
\begin{physics}
In the Leech-lattice CFT the Golay code is the ``code CFT'' skeleton: the $2^{12}$ codewords label
the ways of turning on $\Z/2$ momenta along the $24$ coordinate circles consistently with modular
invariance (Chapter~\ref{ch:02}).  The MOG coordinates are only a convenient labelling that makes the
symmetry group $M_{24}$ and the sextet structure visible; they play the role of a choice of basis
for the $24$ bosons.
\end{physics}
\begin{status}
Definition; covered by the audits of Part~A (\S\ref{ch21:sec:audits}).
\end{status}

\begin{theorem}[Our Golay code is AxiomMath's, up to a coordinate permutation]\label{ch21:thm:golaymatch}
Let $\pi$ be the permutation
$\pi=[0,1,2,3,4,8,13,15,23,7,20,18,19,10,17,14,11,21,5,16,22,6,12,9]$ of $\{0,\dots,23\}$ (our
coordinate $i$ goes to AxiomMath's coordinate $\pi(i)$).  Transporting words by $x\mapsto x\circ\pi^{-1}$
maps our code $\mathcal G$ onto $\chXXIC$.
\end{theorem}
\begin{leancode}
theorem golay_map_eq_C :
    _root_.Code.golay.map (permLin (ZMod 2) piPerm : (Fin 24 → ZMod 2) →ₗ[ZMod 2] _) = Code.C :=
\end{leancode}
(\leanpath{VOA/\allowbreak Moonshine/\allowbreak MonsterGriess/\allowbreak LeechMatch.lean},
\lean{VOA.Moonshine.MonsterGriess.LeechMatch.golay_map_eq_C}.)
\begin{dictionary}
\lean{_root_.Code.golay} is the code of Chapter~\ref{ch:02}; \lean{permLin R σ x = x ∘ σ⁻¹}, so
\lean{(permLin σ x) (σ i) = x i}; \lean{Submodule.map} is the image.  \lean{piPerm} is $\pi$, given
together with its explicit inverse.  The statement is an equality of subspaces, the strongest form.
\end{dictionary}
\begin{proof}[Proof on paper]
The twelve rows of $[I_{12}\mid B]$, moved by $\pi$, satisfy the two MOG conditions (a finite check,
\lean{golayRow_perm_mem_C}, by \lean{decide +kernel}); so the image of $\mathcal G$ lies in
$\chXXIC$.  Both have dimension $12$.  The permutation itself was found numerically by octad
backtracking (a search script that is not part of the release); only the check is in Lean.
\end{proof}
\begin{physics}
Two descriptions of the same $24$ compact bosons differing by a relabelling of the bosons: a
permutation of the fields of the free CFT, under which the Golay-code skeleton matches.
\end{physics}
\begin{status}
No hypotheses.  \leanpath{LeechMatch.lean} has no audit file of its own; it is used by, and so
covered by, the audit of \lean{griessIso} (\leanpath{MoonshineCNGriessAudit3.lean}).
\end{status}

\subsection{The Parker cocycle, the cocode and its pairing}

Why a cocycle is needed.  The extraspecial group $2^{1+24}$ of the Monster is built from the Golay code
and its cocode, and its ``code part'' must be a \emph{non-abelian} lift of the abelian group
$\chXXIC$: one wants signed matrices $X_a$ ($a\in\chXXIC$) with $X_aX_b=\pm X_{a+b}$, the sign being
$(-1)^{\theta(a,b)}$.  Associativity of matrix multiplication forces the coboundary of $\theta$ to be
the triple-intersection form, and the commutator of $X_a,X_b$ to be $(-1)^{|a\cap b|/2}$; Conway's
solution is the Parker loop.

\begin{definition}[Theta cocycle]\label{ch21:def:theta}
For binary words write $P(d)=\tfrac14|d|$, $K(d,e)=\tfrac12|d\cap e|$, $A(d,e,f)=|d\cap e\cap f|$,
all modulo $2$ (the divisions are exact on $\chXXIC$, which is doubly even and self-orthogonal).  Let
$B_1,\dots,B_{12}$ be the basis \lean{basisC} of $\chXXIC$ and, for $d=\sum\mu_iB_i$, $e=\sum\nu_jB_j$,
\begin{equation*}
  \theta(d,e)=\sum_{i,j}\mu_i\nu_j\beta_{ij}+\sum_{i<j}\sum_k\mu_i\mu_j\nu_k\,A(B_i,B_j,B_k),\qquad
  \beta_{ij}=\begin{cases}0&i<j,\\P(B_i)&i=j,\\K(B_i,B_j)&i>j.\end{cases}
\end{equation*}
Then $\theta$ is linear in $e$ and satisfies
\begin{equation*}
  \theta(d+e,f)=\theta(d,f)+\theta(e,f)+A(d,e,f),\qquad\theta(d,e)+\theta(e,d)=K(d,e),\qquad\theta(d,d)=P(d).
\end{equation*}
These three identities are what is used below; the explicit formula only serves to exhibit one
solution.  The \emph{Parker loop} is $\chXXIC\times\chXXIF$ with $(a,s)(b,t)=(a+b,s+t+\theta(a,b))$.
\end{definition}
\begin{leancode}
def θc (μ ν : Fin 12 → ZMod 2) : ZMod 2 :=
  ∑ i, ∑ j, μ i * ν j * β i j +
    ∑ i, ∑ j, ∑ k, if i < j then μ i * μ j * ν k * A (B i) (B j) (B k) else 0

noncomputable def θ (d e : C) : ZMod 2 := θc (basisC.repr d) (basisC.repr e)
\end{leancode}
(\leanpath{VOA/\allowbreak Moonshine/\allowbreak MonsterGriess/\allowbreak Golay.lean}.)
\begin{leancode}
theorem θ_add_left (d e f : C) : θ (d + e) f = θ d f + θ e f + A d e f := by
\end{leancode}
\begin{leancode}
theorem θ_add_θ_swap (d e : C) : θ d e + θ e d = K d e := by
\end{leancode}
\begin{leancode}
theorem θ_self (d : C) : θ d d = P d := by
\end{leancode}
(\leanpath{VOA/\allowbreak Moonshine/\allowbreak MonsterGriess/\allowbreak CodeCocycle.lean};
namespace \lean{VOA.Moonshine.MonsterGriess.Code}.)
\begin{dictionary}
\lean{β i j} is $\beta_{ij}$ and \lean{B i} the $i$-th basis word; \lean{basisC.repr d} the coordinate
vector of $d$.  \lean{A}, \lean{K}, \lean{P} are defined in \leanpath{Golay.lean} by the cardinality
formulas above.  The Parker loop is \lean{Code.Parker}, a structure with fields \lean{code} and
\lean{sign}.  Linearity in the second argument is \lean{θ_add_right}.
\end{dictionary}
\begin{physics}
The Parker loop plays for the Golay code the role that the cocycle $\varepsilon(\alpha,\beta)$ plays
for the Leech lattice in the vertex operators $e^\alpha e^\beta=\varepsilon(\alpha,\beta)e^{\alpha+\beta}$:
it is the sign bookkeeping that makes ``exponentials of codewords'' (the $X$-matrices on the spin
space) multiply consistently.
\end{physics}
\begin{status}
No hypotheses; the identities are audited in \leanpath{MoonshineCNGriessAudit2.lean}.
\end{status}

\begin{definition}[Cocode, pairing, $\chXXIC_0$, $\theta_D$]\label{ch21:def:cocode}
The cocode is $\chXXIC^*=\chXXIF^{24}/\chXXIC$; the pairing $\langle a,\delta\rangle$ of
$a\in\chXXIC$ and $\delta=[x]\in\chXXIC^*$ is $a\cdot x$, well defined by self-orthogonality.
$\Omega=(1,\dots,1)\in\chXXIC$; $\mu_0:\chXXIC\to\chXXIF$ is the first coordinate of a codeword in
the basis \lean{basisC}, with $\mu_0(\Omega)=1$; $\chXXIC_0=\ker\mu_0$ (a complement of
$\langle\Omega\rangle$, so $\chXXIC=\chXXIC_0\oplus\langle\Omega\rangle$ and $|\chXXIC_0|=2^{11}=2048$) and $\pi_0:\chXXIC\to\chXXIC_0$ the projection along
$\Omega$.  For $a\in\chXXIC$, $\theta_D(a)\in\chXXIC^*$ is a cocode element with
$\langle u,\theta_D(a)\rangle=\theta(a,u)$ for all $u$ (it exists because the pairing is perfect).
\end{definition}
\begin{leancode}
def μ₀ : C →ₗ[ZMod 2] ZMod 2 := (LinearMap.proj 0 ∘ₗ ψ) ∘ₗ C.subtype
\end{leancode}
\begin{leancode}
def C₀ : Submodule (ZMod 2) C := LinearMap.ker μ₀
\end{leancode}
\begin{leancode}
def π₀ : C →ₗ[ZMod 2] C₀ :=
  LinearMap.codRestrict C₀ (LinearMap.id + μ₀.smulRight ΩC) fun u => by
\end{leancode}
\begin{leancode}
abbrev Cocode : Type := (Fin 24 → ZMod 2) ⧸ C
\end{leancode}
\begin{leancode}
def pair (a : C) : Cocode →ₗ[ZMod 2] ZMod 2 :=
  C.liftQ (dotL a) fun _ hx => LinearMap.mem_ker.2 (dotProduct_eq_zero a.2 hx)
\end{leancode}
\begin{leancode}
noncomputable def thetaD (a : C) : Cocode := (exists_pair_eq (θL a)).choose
\end{leancode}
(\leanpath{VOA/\allowbreak Moonshine/\allowbreak MonsterGriess/\allowbreak Spin.lean}, namespace \lean{VOA.Moonshine.MonsterGriess.Spin}.)
\begin{dictionary}
\lean{ψ} is the coordinate map of \lean{basisC}; \lean{LinearMap.proj 0} takes its $0$-th entry.
\lean{μ₀.smulRight ΩC} is $u\mapsto\mu_0(u)\Omega$, so $\pi_0(u)=u+\mu_0(u)\Omega$.
\lean{C.liftQ} descends \lean{dotL a} $=(x\mapsto a\cdot x)$ to the quotient.  \lean{thetaD} is
defined by \lean{Classical.choose}; by nondegeneracy (\lean{eq_zero_of_forall_pair}) the chosen
element is unique, so the choice is harmless.
\end{dictionary}
\begin{physics}
The cocode is the dual of the code: in the lattice CFT it labels the classes of ``half-integral
windings'' $\tfrac12x$ modulo those already accounted for by the code; the pairing is the mutual
locality phase between the two kinds of charges.
\end{physics}
\begin{status}
Definitions; \lean{finrank_cocode} ($=12$) and \lean{card_C₀} ($=2048$) are proved.
\end{status}

%======================================================================
\section{The Leech lattice in Conway's coordinates, and \texorpdfstring{$\Lambda/2\Lambda$}{Lambda/2Lambda}}\label{ch21:sec:leech}

\subsection{The code lattice}

\begin{definition}[AxiomMath's Leech lattice]\label{ch21:def:codeLattice}
For a binary code $D$, $\Lambda(D)\subset\Z^{24}$ is the set of $u$ for which there are
$m\in\{0,1\}$ and $d\in D$ with $u_i\equiv m+2d_i\pmod 4$ for all $i$ and
$\sum_iu_i\equiv4m\pmod8$.  The form is $\langle u,v\rangle=u\cdot v/8$.  The minimal vectors are the
$u\in\Lambda(\chXXIC)$ with $\langle u,u\rangle=4$, i.e.\ $u\cdot u=32$.
\end{definition}
\begin{leancode}
def codeLattice (C : Submodule (ZMod 2) (Fin 24 → ZMod 2)) : Set (Fin 24 → ℤ) :=
  {u | ∃ m : ℤ, (m = 0 ∨ m = 1) ∧ ∃ d ∈ C,
    (∀ i, u i ≡ m + 2 * bitLift (d i) [ZMOD 4]) ∧ ∑ i, u i ≡ 4 * m [ZMOD 8]}

def codeForm (u v : Fin 24 → ℚ) : ℚ := (∑ i, u i * v i) / 8
\end{leancode}
\begin{leancode}
def minVecs : Set (Fin 24 → ℤ) :=
  {u | u ∈ codeLattice C ∧ codeForm (fun i => (u i : ℚ)) (fun i => (u i : ℚ)) = 4}
\end{leancode}
\begin{leancode}
noncomputable def classOf (u : Fin 24 → ℤ) : codeLatticeModTwo C allOnes_mem_C :=
  open Classical in
  if h : u ∈ codeLattice C then Submodule.Quotient.mk (⟨u, h⟩ : codeLatticeSubgroup allOnes_mem_C)
  else 0
\end{leancode}
\begin{leancode}
abbrev L2 : Type := codeLatticeModTwo C allOnes_mem_C
\end{leancode}
(\leanpath{VOA/\allowbreak Moonshine/\allowbreak MonsterGriess/\allowbreak Lattice.lean}; \lean{codeLattice}, \lean{codeForm} in
\lean{VOA.Moonshine.MonsterGriess}, \lean{minVecs}, \lean{classOf} in \lean{.Code}, \lean{L2} in
\lean{.SpinLattice}.)
\begin{dictionary}
\lean{bitLift t} $\in\{0,1\}$ is the integer lift of $t\in\chXXIF$; \lean{a ≡ b [ZMOD n]} is
congruence of integers.  \lean{codeLatticeSubgroup hΩ} is the same set as an additive subgroup (closure
under addition uses $\Omega\in D$); \lean{codeLatticeModTwo} is \lean{ModN (codeLatticeSubgroup hΩ) 2},
Mathlib's quotient of a $\Z$-module by $2$ times itself, i.e.\ $\Leech/2\Leech$ as an
$\chXXIF$-vector space.  \lean{classOf} is the class map, extended by $0$ off the lattice (a
convention that is only ever applied to lattice vectors).  This is Conway's standard description of
the Leech lattice scaled by $\sqrt8$: the vectors with $m=0$ are the even vectors $u\equiv2d\pmod4$
with $\sum u_i\equiv0\pmod8$, and those with $m=1$ the odd vectors.
\end{dictionary}
\begin{physics}
$\Leech$ is the momentum lattice of the $24$ chiral bosons; minimal vectors $\alpha$
($\langle\alpha,\alpha\rangle=4$) give the vertex operators $e^{i\alpha\cdot X}$ of conformal weight
$\tfrac12\langle\alpha,\alpha\rangle=2$, i.e.\ the charged spin-two fields of $\Vnat_2$.  The quotient
$\Leech/2\Leech$ is the group of ``charges modulo the $\Z/2$ orbifold'': the twist fields see momenta
only modulo $2\Leech$.
\end{physics}
\begin{status}
Definitions; the lattice is identified with ours in Theorem~\ref{ch21:thm:leechmatch}.
\end{status}

\begin{theorem}[Our Leech lattice is AxiomMath's, with no rescaling]\label{ch21:thm:leechmatch}
The map $x\mapsto x\circ\pi^{-1}$ is an isomorphism of groups from our Leech lattice
(Chapter~\ref{ch:02}) onto $\Lambda(\chXXIC)$, and it carries our form to \lean{codeForm} exactly.
\end{theorem}
\begin{leancode}
theorem perm_mem_codeLattice_C_iff (x : Fin 24 → ℤ) :
    permLin ℤ piPerm x ∈ codeLattice Code.C ↔ x ∈ Leech.carrier := by
\end{leancode}
\begin{leancode}
def leechEquiv : leech.carrier ≃+ ↥(codeLatticeSubgroup Code.allOnes_mem_C) := carrierEquiv
\end{leancode}
\begin{leancode}
theorem codeForm_leechEquiv (x y : leech.carrier) :
    codeForm (toQ (leechEquiv x)) (toQ (leechEquiv y)) = (leech.form x y : ℚ) := by
\end{leancode}
(\leanpath{VOA/\allowbreak Moonshine/\allowbreak MonsterGriess/\allowbreak LeechMatch.lean}, namespace
\lean{VOA.Moonshine.MonsterGriess.LeechMatch}.)
\begin{dictionary}
\lean{Leech.carrier} $=2A\cup(u+2A)\subset\Z^{24}$ with $u=(5,1,\dots,1)$ is our lattice in
Chapter~\ref{ch:02}'s coordinates, and \lean{leech} the integral lattice built from it, with
\lean{leech.form x y} $=x\cdot y/8$.  The key step \lean{mem_codeLattice_golay_iff} says that our
carrier is AxiomMath's congruence description applied to our code; then
Theorem~\ref{ch21:thm:golaymatch} and permutation-equivariance of \lean{codeLattice} finish.
\lean{toQ} reads an integer vector in $\Q^{24}$.  ``No rescaling'': both models use the same
$\sqrt8$-scaled integer coordinates, so the isometry statement has scale factor $1$.
\end{dictionary}
\begin{physics}
The same compact boson theory written in two bases of the momentum lattice related by a permutation of
coordinates; the CFT data (momenta, inner products, hence conformal weights) are unchanged.
\end{physics}
\begin{status}
No hypotheses; covered by the audit of \lean{griessIso} (\leanpath{MoonshineCNGriessAudit3.lean}),
standard axioms.
\end{status}

\subsection{Minimal vectors and the count \texorpdfstring{$98280$}{98280}}

\begin{theorem}[Minimal vectors]\label{ch21:thm:minvecs}
(a) Every nonzero $u\in\Lambda(\chXXIC)$ has $u\cdot u\ge32$ (the Leech lattice has no roots).
(b) Two minimal vectors are congruent modulo $2\Leech$ iff they are equal up to sign.
(c) The minimal vectors are the union of the three disjoint families
$(\pm4,\pm4,0^{22})$, $(\pm2^8,0^{16})$ on an octad with an even number of minus signs, and
$(\mp3,\pm1^{23})$; there are $1104+97152+98304=196560$ of them.
(d) $|\chXXIS|=98280$ (this part uses the bijectivity of the map $\varphi_0$ of Theorem~\ref{ch21:thm:phi0}, proved independently below).
\end{theorem}
\begin{leancode}
theorem normSq_ge {u : Fin 24 → ℤ} (hu : u ∈ codeLattice C) (h0 : u ≠ 0) : 32 ≤ normSq u := by
\end{leancode}
\begin{leancode}
theorem classOf_eq_iff {u v : Fin 24 → ℤ} (hu : u ∈ minVecs) (hv : v ∈ minVecs) :
    classOf u = classOf v ↔ v = u ∨ v = -u := by
\end{leancode}
(\leanpath{VOA/\allowbreak Moonshine/\allowbreak MonsterGriess/\allowbreak MinimalVectors.lean}, namespace \lean{.Code}.)
\begin{leancode}
theorem minVecs_eq : minVecs = minFamily₁ ∪ minFamily₂ ∪ minFamily₃ := by
\end{leancode}
\begin{leancode}
theorem card_minFin : minFin.card = 196560 := by
\end{leancode}
(\leanpath{VOA/\allowbreak Moonshine/\allowbreak MonsterGriess/\allowbreak LatticeFamilies.lean}, namespace \lean{.Code}.)
\begin{leancode}
theorem card_shortIdx : Nat.card ShortIdx = 98280 := by
\end{leancode}
(\leanpath{VOA/\allowbreak Moonshine/\allowbreak MonsterGriess/\allowbreak LatticeFamilies.lean}, \lean{VOA.Moonshine.MonsterGriess.card_shortIdx}.)
\begin{dictionary}
\lean{normSq u} $=\sum u_i^2$.  \lean{minFin} is a \lean{Finset} enumerating the three families
through explicit injective parametrisations, and \lean{minVecs_eq} shows that it is all of
\lean{minVecs}.  \lean{ShortIdx} is defined in Definition~\ref{ch21:def:short} below.  There is a
second, identical statement \lean{VOA.Moonshine.GriessDictionary.card_shortIdx}
(\leanpath{MoonshineCNGriessSign.lean}), derived from our lattice via \lean{LeechMatch}; and a second
\lean{classOf_eq_iff} in \lean{LeechMatch} with the disjunction written as \lean{u = v ∨ u = -v}.
These duplicates are harmless.
\end{dictionary}
\begin{proof}[Proof on paper]
(a) For a lattice vector $u$, $u\cdot u\equiv0\pmod{16}$ (a congruence computed from the
definition and the double-evenness of $\chXXIC$), so it suffices to exclude $u\cdot u=16$.  If $m=1$
all $24$ coordinates are odd and $u\cdot u\ge24$.  If $m=0$ and $u\equiv2d\pmod4$ with $d\ne0$, the
$\ge8$ coordinates in the support of $d$ are $\equiv2\pmod4$, so $u\cdot u\ge32$.  If $m=0$ and $d=0$,
all coordinates are divisible by $4$ with $\sum u_i\equiv0\pmod8$, so a single nonzero coordinate must
be $\pm8$ ($u\cdot u=64$), and otherwise there are at least two nonzero coordinates
($u\cdot u\ge32$).  (b) If $u\equiv v\pmod{2\Leech}$ with $u\neq\pm v$, then $(u-v)/2$ and $(u+v)/2$ are
nonzero lattice vectors with $\|(u-v)/2\|^2+\|(u+v)/2\|^2=\tfrac12(\|u\|^2+\|v\|^2)=32$, so one of them
has $u\cdot u<32$, contradicting (a).  (c) is a case analysis of the shapes allowed by the
congruences with $\sum u_i^2=32$ (the shapes are those tabulated in Conway--Sloane, \emph{Sphere
Packings, Lattices and Groups}, Ch.~4 \S11); the counts use $759$ octads (\lean{wtCount_eight}, kernel computation):
$1104=4\binom{24}2$, $97152=759\cdot2^7$, $98304=24\cdot2^{12}$.  (d) By (b) each class of a minimal
vector contains exactly two of them, and $\varphi_0$ (below) is a bijection, so $|\chXXIS|=196560/2$.
\end{proof}
\begin{physics}
$196560$ is the number of charged spin-two fields $e^{i\alpha\cdot X}$ of the Leech lattice CFT, and
$98280$ the number of $\theta$-invariant combinations $e^{i\alpha X}+e^{-i\alpha X}$ that survive the
orbifold projection.  Absence of roots ($\langle\alpha,\alpha\rangle=2$) is the absence of
spin-one currents: the Leech CFT has only the $24$ abelian currents $\partial X^i$, and these are
projected out by the orbifold, which is why $\Vnat_1=0$.
\end{physics}
\begin{status}
No hypotheses.  \lean{normSq_ge}, \lean{classOf_eq_iff}, \lean{minVecs_eq}, \lean{card_minFin},
\lean{card_shortIdx} are audited in \leanpath{ParityAudit.lean} (and \lean{card_shortIdx} also in
\leanpath{TrialityAudit.lean} and \leanpath{MoonshineCNGriessAudit2.lean}).
\end{status}

\subsection{Coordinates on \texorpdfstring{$\Lambda/2\Lambda$}{Lambda/2Lambda}: the map \texorpdfstring{$\varphi_0$}{phi0}}

The space $\chXXIL=\Leech/2\Leech$ is $24$-dimensional over $\chXXIF$; the quadratic form
$q(\bar u)=\tfrac12\langle u,u\rangle\bmod2=u\cdot u/16\bmod 2$ is well defined on it, with polar
form $\langle u,v\rangle\bmod2$.  AxiomMath coordinatises $\chXXIL$ by $\chXXIC\times\chXXIC^*$.

\paragraph{The idea.}  The quadratic space $(\chXXIL,q)$ has a ``hyperbolic'' splitting into two
complementary $12$-dimensional subspaces on which $q$ vanishes, one coming from the code and one from
the cocode, paired with each other by the code--cocode pairing.  In such coordinates
$q(a,\delta)=\langle a,\delta\rangle$, and the extraspecial group of $\chXXIL$ becomes a Heisenberg
group of $\chXXIC\times\chXXIC^*$, which is what the spin matrices of \S\ref{ch21:sec:spin} realise.
Concretely: (i) every binary word $v$ gives the lattice vector $(|v|-4v_k)_k$ (for $v=e_k$ it is the odd
minimal vector $(1^{23},-3)$, with the $-3$ in place $k$); its class $\lambda(v)$ is additive in $v$ and
vanishes for $v\in\chXXIC$, so it defines $\lambda_{\mathrm{coc}}$ on the cocode.  (ii) A codeword $a$ has
the canonical even lift $2a$ (or $2(\Omega-a)$ when $|a|\equiv4\pmod8$), whose class is $\lambda_X(a)$;
but $\lambda_X$ is additive only up to the defect $A$, and adding $\lambda_{\mathrm{coc}}(\theta_D(a))$
cancels this defect, which is why $\lambda_T$ rather than $\lambda_X$ is used.
For a binary word $v$ let $\lambda(v)$ be the class of the lattice vector $(|v|-4v_k)_k$; it vanishes
on $\chXXIC$, hence descends to $\lambda_{\mathrm{coc}}:\chXXIC^*\to\chXXIL$.  For a codeword $a$, let
$\lambda_X(a)$ be the class of $2a+2P(a)-4P(a)a$ (coordinatewise, with $P(a)=|a|/4\bmod 2$ lifted to
$\{0,1\}$; this is $2a$ if $P(a)=0$ and $2(\Omega-a)$ if $P(a)=1$), and $\lambda_T(a)=\lambda_X(a)+\lambda_{\mathrm{coc}}(\theta_D(a))$.

\begin{definition}[$\varphi_0$, short coordinates, $u_c$]\label{ch21:def:short}
$\varphi_0(a,\delta)=\lambda_T(a)+\lambda_{\mathrm{coc}}(\delta)$.  A pair $p\in\chXXIC\times\chXXIC^*$
is \emph{short} if $\varphi_0(p)$ is the class of a minimal vector; $\chXXIS$ is the set of short
pairs.  For $c\in\chXXIS$, $u_c$ is the lexicographically smaller of the two minimal vectors $\pm u$
in the class $\varphi_0(c)$.
\end{definition}
\begin{leancode}
noncomputable def lamT (a : C) : L2 := lamX a + lamCocode (thetaD a)
\end{leancode}
\begin{leancode}
noncomputable def phi0 (p : C × Cocode) : L2 := lamT p.1 + lamCocode p.2
\end{leancode}
\begin{leancode}
def IsShortCoord (p : C × Cocode) : Prop := ∃ u ∈ minVecs, classOf u = phi0 p

abbrev ShortIdx : Type := {p : C × Cocode // IsShortCoord p}
\end{leancode}
\begin{leancode}
def LexLE (u v : Fin 24 → ℤ) : Prop := toLex u ≤ toLex v

noncomputable def uvec (c : ShortIdx) : Fin 24 → ℤ :=
  open Classical in
  if LexLE c.2.choose (-c.2.choose) then c.2.choose else -c.2.choose
\end{leancode}
(\leanpath{VOA/\allowbreak Moonshine/\allowbreak MonsterGriess/\allowbreak Short.lean}; \lean{lamT}, \lean{phi0} in
\lean{.SpinLattice}, the rest in \lean{VOA.Moonshine.MonsterGriess}.)
\begin{dictionary}
\lean{lamX a} $=\lambda_X(a)$, \lean{lamCocode} $=\lambda_{\mathrm{coc}}$ (\leanpath{Lattice.lean},
\leanpath{Short.lean}).  \lean{c.2.choose} is \emph{some} minimal vector in the class $\varphi_0(c)$
(extracted from the existential \lean{IsShortCoord}); \lean{toLex} compares lexicographically.  By
Theorem~\ref{ch21:thm:minvecs}(b) the class contains exactly $\pm$ that vector, so \lean{uvec c} does
not depend on which representative the choice returns.  This matters: $u_c$ enters the product of
$\chXXIB$ through $U_c=u_cu_c^{\mathsf T}$ and $q_c$, which are even in $u_c$, but also through the
identification with $\Vnat$ (\lean{shortVec}), where a sign choice is a basis choice.
\end{dictionary}
\begin{physics}
$\chXXIS$ indexes the $98280$ orbifold-invariant charged spin-two fields
$e^{i\alpha X}+e^{-i\alpha X}$; $u_c$ is a canonical choice of one of the two momenta $\pm\alpha$.
\end{physics}
\begin{status}
Definitions (ported verbatim except that \lean{phi0} is a function; see the warning in
\S\ref{ch21:sec:fidelityaudit}).  Covered by the audits of Part~A.
\end{status}

\begin{theorem}[$\varphi_0$ is a linear isomorphism; $q\circ\varphi_0$ is the pairing]\label{ch21:thm:phi0}
$\varphi_0$ is additive and bijective, so it is an $\chXXIF$-linear isomorphism
$\chXXIC\times\chXXIC^*\cong\chXXIL$; in these coordinates $q(\varphi_0(a,\delta))=\langle a,\delta\rangle$
and the polar form is $\langle a,\varepsilon\rangle+\langle b,\delta\rangle$.  In particular
$\langle a,\delta\rangle=0$ for every short pair $(a,\delta)$.
\end{theorem}
\begin{leancode}
theorem phi0_bijective : Function.Bijective phi0 := phi0Lin_bijective
\end{leancode}
\begin{leancode}
noncomputable def phi0Equiv : (C × Cocode) ≃ₗ[ZMod 2] L2 :=
  LinearEquiv.ofBijective phi0Lin phi0Lin_bijective
\end{leancode}
\begin{leancode}
theorem qL_phi0 (a : C) (δ : Cocode) : qL (phi0 (a, δ)) = pair a δ := by
\end{leancode}
\begin{leancode}
theorem bL_phi0 (a b : C) (δ ε : Cocode) :
    QuadraticMap.polarBilin qL (phi0 (a, δ)) (phi0 (b, ε)) = pair a ε + pair b δ := by
\end{leancode}
(namespace \lean{.SpinLattice}) and
\begin{leancode}
theorem pair_eq_zero_of_isShortCoord {p : C × Cocode} (h : IsShortCoord p) : pair p.1 p.2 = 0 := by
\end{leancode}
(\leanpath{VOA/\allowbreak Moonshine/\allowbreak MonsterGriess/\allowbreak LatticeQuadratic.lean}.)
\begin{dictionary}
\lean{qL} is the quadratic form $q$ on \lean{L2}: by \lean{qL_cls}, $q$ of the class of $u$ is
$\tfrac12(u\cdot u/8)\bmod2$.  \lean{QuadraticMap.polarBilin} is Mathlib's polar form
$B(x,y)=q(x+y)-q(x)-q(y)$.  The rank of $\chXXIL$ is obtained by a port-specific argument (a basis
from \lean{Submodule.basisOfPid} and the $24$ independent vectors $8e_i$), recorded in the port notes;
the statement \lean{finrank_L2} ($=24$) is the source's.
\end{dictionary}
\begin{proof}[Proof on paper]
Additivity of $\lambda_T$ is the cocycle identity $\theta_D(a+b)=\theta_D(a)+\theta_D(b)+A(a,b)$
(\lean{thetaD_add}, from $\theta(a+b,f)=\theta(a,f)+\theta(b,f)+A(a,b,f)$) combined with the
corresponding defect of $\lambda_X$ (\lean{lamCodeCls_add}): the two defects cancel.  Injectivity:
if $\varphi_0(a,\delta)=0$, then its polar pairing with every $\varphi_0(b,\varepsilon)$ vanishes,
i.e.\ $\langle a,\varepsilon\rangle+\langle b,\delta\rangle=0$ for all $(b,\varepsilon)$, and
nondegeneracy of the code--cocode pairing gives $a=0$, $\delta=0$; since both sides have dimension
$24$ ($12+12$, \lean{finrank_L2}) the map is bijective.  The quadratic-form identity is a
computation on representatives: $q(\lambda_T(a))=0$, $q(\lambda_{\mathrm{coc}}(\delta))=0$ and the
cross term is $\langle a,\delta\rangle$.  Finally a minimal vector has $q=\tfrac12\cdot4\bmod2=0$.
\end{proof}
\begin{physics}
$\Leech/2\Leech$ with its quadratic form is the ``charge space'' of the twist fields: the
$2^{24}$ classes label the generators of the extraspecial group that acts on the $2^{12}$ twisted
ground states, and $q$ records whether two such generators commute or anticommute with themselves
squared.  Code${}\times{}$cocode coordinates split this charge into a ``code part'' and a ``cocode
part'', like position and momentum in a Heisenberg group.
\end{physics}
\begin{status}
No hypotheses.  \lean{phi0Equiv}, \lean{qL_phi0}, \lean{pair_eq_zero_of_isShortCoord},
\lean{phi0_bijective} are audited in \leanpath{MoonshineCNGriessAudit2.lean}.
\end{status}

\begin{remark}[The other classes of $\Leech/2\Leech$]\label{ch21:rem:classes}
Classically, the $2^{24}=16777216$ classes split as $1$ (zero) $+\,98280$ (classes of $\pm$ a minimal
vector) $+\,8386560$ (classes of $48$ vectors of norm $6$) $+\,8292375$ (classes of a frame of $48$
vectors of norm $8$).  Only the first two kinds enter the Griess algebra.  This count is quoted for
orientation and is not formalised.
\end{remark}

%======================================================================
\section{The spin module and the extraspecial group \texorpdfstring{$2^{1+24}$}{2\^{}(1+24)}}\label{ch21:sec:spin}

\subsection{Signed permutation matrices on \texorpdfstring{$4096$}{4096} coordinates}

\paragraph{The idea.}  By \S\ref{ch21:sec:leech}, $\chXXIL\cong\chXXIC\times\chXXIC^*$ with
$q(a,\delta)=\langle a,\delta\rangle$.  The extraspecial group attached to $(\chXXIL,q)$ is then a
Heisenberg group of $\chXXIC\times\chXXIC^*$, and it has a ``Schr\"odinger'' representation on functions
on the code: codewords act by translations, cocode elements by the characters
$c\mapsto(-1)^{\langle c,\delta\rangle}$.  The $2^{12}=4096$ codewords are parametrised by
$I=\chXXIF\times\chXXIC_0$ through $(\eta,c)\mapsto c+\eta\Omega$ (because
$\chXXIC=\chXXIC_0\oplus\langle\Omega\rangle$, $|\chXXIC_0|=2048$).  AxiomMath's matrices $R(s,a,\delta)$ are
this representation written in the coordinates $(\eta,c)$: the $\chXXIC_0$-part of a translation by $a$
moves $c$ to $\pi_0(c+a)$, its $\Omega$-part is absorbed in the sign $(-1)^{\eta\mu_0(a)}$, an odd cocode
element ($\langle\Omega,\delta\rangle=1$) flips $\eta$, and the character contributes
$(-1)^{\langle\pi_0(c+a),\delta\rangle}$; $s$ is a global sign.

\begin{definition}[Spin index and spin matrices]\label{ch21:def:qmatR}
$I=\chXXIF\times\chXXIC_0$ ($2\cdot2048=4096$ elements).  For $s\in\chXXIF$, $a\in\chXXIC$, $\delta\in\chXXIC^*$
and a commutative ring $R$, $R(s,a,\delta)$ is the $I\times I$ signed permutation matrix with
\begin{equation*}
  R(s,a,\delta)_{i,j}=\begin{cases}(-1)^{\sigma(i)}&j=p(i),\\0&\text{otherwise,}\end{cases}\qquad
  p(\eta,c)=(\eta+\langle\Omega,\delta\rangle,\ \pi_0(c+a)),\quad
  \sigma(\eta,c)=s+\eta\,\mu_0(a)+\langle\pi_0(c+a),\delta\rangle .
\end{equation*}
\end{definition}
\begin{leancode}
abbrev SpinIdx : Type := ZMod 2 × C₀

def shift (a : C) (c : C₀) : C₀ := π₀ ((c : C) + a)
\end{leancode}
\begin{leancode}
def qPerm (a : C) (δ : Cocode) (i : SpinIdx) : SpinIdx :=
  (i.1 + pair ΩC δ, shift a i.2)

def qSign (s : ZMod 2) (a : C) (δ : Cocode) (i : SpinIdx) : ZMod 2 :=
  s + i.1 * μ₀ a + pair (shift a i.2) δ
\end{leancode}
\begin{leancode}
def qmatR (s : ZMod 2) (a : C) (δ : Cocode) : Matrix SpinIdx SpinIdx R :=
  spm R (qPerm a δ) (qSign s a δ)
\end{leancode}
\begin{leancode}
def spm (σ : ι → ι) (ε : ι → ZMod 2) : Matrix ι ι R :=
  Matrix.of fun i j => if j = σ i then sgnR R (ε i) else 0
\end{leancode}
(\leanpath{VOA/\allowbreak Moonshine/\allowbreak MonsterGriess/\allowbreak Spin.lean}, namespace \lean{.Spin}.)
\begin{dictionary}
\lean{sgnR R t} is $1$ for $t=0$ and $-1$ for $t=1$; \lean{i.1}, \lean{i.2} are the two components of
$i=(\eta,c)$.  The matrices have entries in $\{0,\pm1\}$, so the ring $R$ is irrelevant:
\lean{SpinMatrixMap.lean} proves \lean{map_qmatR : (qmatR R s a δ).map f = qmatR S s a δ} for every ring
homomorphism $f$, which is how AxiomMath's rational matrices are read over $\C$.
\end{dictionary}
\begin{physics}
A signed permutation of the $4096$ twist ground states: the matrices by which lattice exponentials
$e^\alpha$ act on the degenerate ground state of the twisted sector, in a basis where they permute
basis states up to sign.
\end{physics}
\begin{status}
Definitions, verbatim from the source.
\end{status}
\begin{theorem}[Product law]\label{ch21:thm:qmatRmul}
$R(s,a,\delta)R(t,b,\varepsilon)=R(s+t+\langle b,\delta\rangle,a+b,\delta+\varepsilon)$ and
$R(s,a,\delta)^{\mathsf T}=R(s+\langle a,\delta\rangle,a,\delta)$; each $R(s,a,\delta)$ is orthogonal.
\end{theorem}
\begin{leancode}
theorem qmatR_mul (s t : ZMod 2) (a b : C) (δ ε : Cocode) :
    qmatR R s a δ * qmatR R t b ε = qmatR R (s + t + pair b δ) (a + b) (δ + ε) := by
\end{leancode}
\begin{leancode}
theorem qmatR_transpose (s : ZMod 2) (a : C) (δ : Cocode) :
    (qmatR R s a δ)ᵀ = qmatR R (s + pair a δ) a δ := by
\end{leancode}
(\leanpath{VOA/\allowbreak Moonshine/\allowbreak MonsterGriess/\allowbreak SpinMatrixLaws.lean}.)
\begin{proof}[Proof on paper]
Products of signed permutation matrices compose the permutations and add the signs along the path;
the permutations compose as $(a,\delta)\mapsto(a+b,\delta+\varepsilon)$ because $\pi_0$ is linear and
$\pi_0(\Omega)=0$, and the signs add up to $s+t+\langle b,\delta\rangle$ by bilinearity of the pairing.
\end{proof}
\begin{physics}
These are the $2^{12}\times2^{12}$ ``gamma matrices'' of the twisted sector: in the $\Z/2$ orbifold
the $24$ bosons have half-integral modes and no zero modes, the ground state is degenerate, and the
lattice exponentials $e^\alpha$ act on the $2^{12}$ ground states through a Clifford-like extraspecial
group.  The product law is the statement that $e^{\alpha}e^{\beta}=\pm e^{\alpha+\beta}$ with a sign
whose commutator is $(-1)^{\langle\alpha,\beta\rangle}$.
\end{physics}
\begin{status}
No hypotheses; \lean{qmatR_mul} audited in \leanpath{MoonshineCNGriessAudit2.lean}.
\end{status}

\subsection{The extraspecial group \texorpdfstring{$Q$}{Q}}

\begin{definition}[Parker loop, $X$- and $D$-matrices, $Q$]\label{ch21:def:Q}
For $e=(a,s)$ in the Parker loop, $X_e$ is the signed permutation matrix on $I$ sending
$(\eta,c)\mapsto(\eta,\pi_0(c+a))$ with sign $(-1)^{s+P(a)+\theta(a,c)+\eta\mu_0(a)}$; for
$\delta\in\chXXIC^*$, $D_\delta$ sends $(\eta,c)\mapsto(\eta+\langle\Omega,\delta\rangle,c)$ with sign
$(-1)^{\langle c,\delta\rangle}$.  In terms of Definition~\ref{ch21:def:qmatR},
$X_{(a,s)}=R(s,a,\theta_D(a))$ and $D_\delta=R(0,0,\delta)$.  $Q\subset GL_I(\Q)$ is the subgroup
generated by all $X_e$ and $D_\delta$.
\end{definition}
\begin{leancode}
@[ext] structure Parker where
  code : C
  sign : ZMod 2
\end{leancode}
\begin{leancode}
def XmatR (e : Parker) : Matrix SpinIdx SpinIdx R := spm R (xPerm e.code) (xSign e)

def DmatR (δ : Cocode) : Matrix SpinIdx SpinIdx R := spm R (dPerm δ) (dSign δ)
\end{leancode}
\begin{leancode}
def xSign (e : Parker) (i : SpinIdx) : ZMod 2 :=
  e.sign + P e.code + θ e.code i.2 + i.1 * μ₀ e.code
\end{leancode}
\begin{leancode}
theorem XmatR_eq_qmatR (e : Parker) : XmatR R e = qmatR R e.sign e.code (thetaD e.code) := by
\end{leancode}
\begin{leancode}
theorem DmatR_eq_qmatR (δ : Cocode) : DmatR R δ = qmatR R 0 0 δ := by
\end{leancode}
\begin{leancode}
def Q : Subgroup (GL SpinIdx ℚ) := Subgroup.closure (Set.range XU ∪ Set.range DU)
\end{leancode}
(\leanpath{VOA/\allowbreak Moonshine/\allowbreak MonsterGriess/\allowbreak SpinGroupQ.lean}; \lean{Parker} in \lean{.Code}, the rest in
\lean{.Spin}.)  The structure of $Q$ is proved in \leanpath{SpinExtraspecial.lean}:
\begin{leancode}
theorem natCard_Q : Nat.card Q = 2 ^ 25 := by
\end{leancode}
\begin{leancode}
theorem frattini_Q : frattini Q = zpowers zQ :=
\end{leancode}
\begin{leancode}
noncomputable def phi : (Q ⧸ center Q) ≃* Multiplicative L2 :=
  quotCenterEquiv.trans (AddEquiv.toMultiplicative phi0Equiv.toAddEquiv)
\end{leancode}
(\leanpath{VOA/\allowbreak Moonshine/\allowbreak MonsterGriess/\allowbreak SpinExtraspecial.lean};
\lean{natCard_Q}, \lean{frattini_Q} in \lean{.Spin}, \lean{phi} in \lean{.SpinLattice}.)
\begin{dictionary}
\lean{xPerm}, \lean{dPerm}, \lean{dSign} are the permutations and signs in the text.
\lean{XU e}, \lean{DU δ} are the matrices \lean{Xmat e}, \lean{Dmat δ} as units of
\lean{GL SpinIdx ℚ} (their inverses are their transposes).  By \lean{XmatR_eq_qmatR} and
\lean{DmatR_eq_qmatR} every generator is a spin matrix, so by the product law
$Q=\{R(s,a,\delta)\}$.  The
group isomorphism \lean{spinEquiv : SpinCoord ≃* Q} identifies $Q$ with the ``coordinate group''
$\chXXIF\times\chXXIC\times\chXXIC^*$ with product
$(s,a,\delta)(t,b,\varepsilon)=(s+t+\langle b,\delta\rangle,a+b,\delta+\varepsilon)$; \lean{zQ} is
$-1$, \lean{center Q} $=$ \lean{frattini Q} $=\{\pm1\}$.  \lean{phi} identifies $Q/Z(Q)$ with
$\chXXIL$ (as a multiplicative group) through $\varphi_0$.  \lean{coordOf r} and
\lean{centralCoord r} (\leanpath{GriessTrilinear.lean}) are the $(a,\delta)$- and $s$-coordinates of
$r\in Q$; \lean{rep p} is the element with coordinates $(0,p)$.
\end{dictionary}
\begin{physics}
$Q\cong2^{1+24}$ is the group generated by the lattice exponentials acting on the twist-field ground
states: it is a central extension of the charge group $\Leech/2\Leech$ by the sign $\pm1$, with
commutator form $(-1)^{\langle\alpha,\beta\rangle}$.  In the Monster it becomes the normal subgroup
$O_2(C(z))$ of the centraliser of the orbifold involution.
\end{physics}
\begin{status}
No hypotheses; audited in \leanpath{ParityAudit.lean}.  The source's statement that $Q$ is extraspecial
(\lean{isExtraspecialPGroup_Q}) is not ported because its predicate lives in a source prelude that pulls
in Feit--Thompson material; its content is the conjunction of \lean{natCard_Q}, \lean{center_Q},
\lean{frattini_Q}.
\end{status}

\subsection{Coordinates on \texorpdfstring{$Q$}{Q} and the vectors attached to short elements}

\begin{definition}[Coordinates, representatives, $\mathrm{Xlift}$]\label{ch21:def:Xlift}
Every $r\in Q$ is $R(s,a,\delta)$ for a unique triple; write $\mathrm{coord}(r)=(a,\delta)$ and
$\mathrm{cc}(r)=s$ (``central coordinate''), and $\mathrm{rep}(p)=R(0,p)$ for $p\in\chXXIC\times\chXXIC^*$.
$r$ is \emph{short} if $\mathrm{coord}(r)\in\chXXIS$.  For $r\in Q$ let
$\mathrm{Xlift}(r)=(-1)^{\mathrm{cc}(r)}e_{\mathrm{coord}(r)}\in\Q^{\chXXIS}$ if $r$ is short, and $0$
otherwise; so $\mathrm{Xlift}(\mathrm{rep}(c))=e_c$ and $\mathrm{Xlift}(-r)=-\mathrm{Xlift}(r)$.
\end{definition}
\begin{leancode}
noncomputable def coordOf (r : Q) : C × Cocode :=
  ((spinEquiv.symm r).a, (spinEquiv.symm r).δ)

noncomputable def centralCoord (r : Q) : ZMod 2 := (spinEquiv.symm r).s
\end{leancode}
\begin{leancode}
noncomputable def rep (p : C × Cocode) : Q := spinEquiv ⟨0, p.1, p.2⟩
\end{leancode}
\begin{leancode}
noncomputable def Xlift (r : Q) : SpaceX :=
  open Classical in
  if h : IsShortCoord (coordOf r) then sgn (centralCoord r) • Pi.single ⟨coordOf r, h⟩ 1 else 0
\end{leancode}
(\leanpath{VOA/\allowbreak Moonshine/\allowbreak MonsterGriess/\allowbreak GriessTrilinear.lean}.)
\begin{dictionary}
\lean{spinEquiv.symm r} is the coordinate triple of $r$ in the group
$\chXXIF\times\chXXIC\times\chXXIC^*$; \lean{Pi.single c 1} is the basis vector $e_c$ of
\lean{SpaceX} $=\Q^{\chXXIS}$ (defined in \S\ref{ch21:sec:spaceB}).  Thus the $98280$-dimensional
block of the Griess algebra is spanned by vectors attached to the short elements of $Q$, each pair
$\pm r$ giving $\pm e_c$.
\end{dictionary}
\begin{physics}
$\mathrm{Xlift}$ says that a charged spin-two field $\cos(\alpha X)$ is labelled by a pair of opposite
elements of the extraspecial group, the sign recording which of the two is used.
\end{physics}
\begin{status}
Definitions, verbatim from the source.
\end{status}

%======================================================================
\section{The rational Griess algebra \texorpdfstring{$\mathcal B$}{B}}\label{ch21:sec:spaceB}

\subsection{The space and the form}

\begin{definition}[The space $\chXXIB$ and the form $b$]\label{ch21:def:spaceB}
$\chXXISym$ is the space of symmetric $24\times24$ rational matrices, $X=\Q^{\chXXIS}$ with standard basis
$(e_c)_{c\in\chXXIS}$, $W$ the space of $I\times24$ rational matrices, and
$\chXXIB=\chXXISym\times X\times W$.  For $u=(A,x,N)$, $v=(A',x',N')$,
\begin{equation*}
  b(u,v)=\tr(AA')+\sum_{c\in\chXXIS}x_cx'_c+\sum_{j\in I}\sum_{i<24}N_{ji}N'_{ji}.
\end{equation*}
\end{definition}
\begin{leancode}
def symMat : Submodule ℚ (Matrix (Fin 24) (Fin 24) ℚ) where
  carrier := {A | Aᵀ = A}
\end{leancode}
\begin{leancode}
abbrev SpaceX : Type := ShortIdx → ℚ

abbrev SpaceW : Type := Matrix SpinIdx (Fin 24) ℚ

abbrev SpaceB : Type := symMat × SpaceX × SpaceW
\end{leancode}
\begin{leancode}
noncomputable def bform : LinearMap.BilinForm ℚ SpaceB :=
  LinearMap.mk₂ ℚ (fun u v => trace ((u.1 : Matrix (Fin 24) (Fin 24) ℚ) *
      (v.1 : Matrix (Fin 24) (Fin 24) ℚ)) +
      ∑ c, u.2.1 c * v.2.1 c + ∑ j, ∑ i, u.2.2 j i * v.2.2 j i)
\end{leancode}
(\leanpath{VOA/\allowbreak Moonshine/\allowbreak MonsterGriess/\allowbreak Algebra.lean}, namespace
\lean{VOA.Moonshine.MonsterGriess}.)
\begin{dictionary}
For $u$ of type \lean{SpaceB}, \lean{u.1} is the $\chXXISym$-component (a subtype element, coerced to a
matrix), \lean{u.2.1} the $X$-component and \lean{u.2.2} the $W$-component.  \lean{LinearMap.mk₂}
packages a function with four linearity proofs (the omitted lines) as a bilinear map.  Since $A$ is
symmetric, $\tr(AA')=\sum_{k,l}A_{kl}A'_{kl}$, so $b$ is the sum of the standard dot products of the
three blocks and is positive definite (\lean{bform_self_pos}).  The dimensions are proved in
\leanpath{TrialityBasis.lean}: $\dim\chXXISym=300$, $\dim X=98280$, $\dim W=98304$, and
\end{dictionary}
\begin{leancode}
theorem finrank_B : finrank ℚ SpaceB = 196884 := by
\end{leancode}
(\leanpath{VOA/\allowbreak Moonshine/\allowbreak MonsterGriess/\allowbreak TrialityBasis.lean}.)
\begin{physics}
The three blocks are the three kinds of spin-two fields of the orbifold CFT (\S\ref{ch21:sec:intro}):
$\chXXISym$ ($300$) the bilinears $\partial X^i\partial X^j$, $X$ ($98280$) the orbifold-invariant
combinations $e^{i\alpha X}+e^{-i\alpha X}$, and $W$ ($24\cdot4096$) the twisted states
$\partial X^i_{-1/2}\sigma_s$.  The positive definite form is the two-point function, positive because
the CFT is unitary.
\end{physics}
\begin{status}
Definitions; \lean{finrank_B} audited in \leanpath{TrialityAudit.lean}.
\end{status}

\subsection{The product}

The product uses three pieces of data attached to a short class $c=(a_c,\delta_c)\in\chXXIS$: the
matrix $U_c=u_cu_c^{\mathsf T}$ ($u_c$ the minimal vector of Definition~\ref{ch21:def:short}, so
$U_c^2=32U_c$), the quadratic form $q_c(M)=u_c^{\mathsf T}Mu_c$, and the spin matrix
$R_c=R(0,a_c,\delta_c)$; and the structure constants
\begin{equation*}
  \tau(d,e,c)=\begin{cases}(-1)^{\langle a_e,\delta_d\rangle}&c=d+e\ \text{in }\chXXIC\times\chXXIC^*,\\
  0&\text{otherwise.}\end{cases}
\end{equation*}

\begin{definition}[The Griess product of $\chXXIB$]\label{ch21:def:gmul}
The product is the bilinear map determined by the following table (and commutativity), for
$A,A'\in\chXXISym$, $x,y\in X$, $N,M\in W$ (``$\cdot$'' is the product, each entry lists the
$(\chXXISym,X,W)$-components):

\begin{center}\small
\begin{tabular}{lll}
\toprule
product & $\chXXISym$-part & $X$-part\quad/\quad $W$-part\\\midrule
$A\cdot A'$ & $2(AA'+A'A)$ & $0$\quad/\quad$0$\\
$A\cdot x$ & $0$ & $c\mapsto\tfrac18x_cq_c(A)$\quad/\quad$0$\\
$x\cdot y$ & $\tfrac18\sum_cx_cy_cU_c$ & $c\mapsto\sum_{d,e}x_dy_e\tau(d,e,c)$\quad/\quad$0$\\
$A\cdot N$ & $0$ & $0$\quad/\quad$\tfrac{\tr A}8N+NA$\\
$x\cdot N$ & $0$ & $0$\quad/\quad$\tfrac18\sum_cx_c(R_cN-\tfrac14R_cNU_c)$\\
$N\cdot M$ & $\tfrac18\tr(N^{\mathsf T}M)1+\tfrac12(N^{\mathsf T}M+M^{\mathsf T}N)$ &
  $c\mapsto\tfrac18(\tr(N^{\mathsf T}R_cM)-\tfrac14q_c(N^{\mathsf T}R_cM))$\quad/\quad$0$\\
\bottomrule
\end{tabular}
\end{center}
\end{definition}
In Lean the product is defined at once on general elements, by three component functions:
\begin{leancode}
noncomputable def Umat (c : ShortIdx) : Matrix (Fin 24) (Fin 24) ℚ := vecMulVec (uQ c) (uQ c)

noncomputable def Rmat (c : ShortIdx) : Matrix SpinIdx SpinIdx ℚ := qmatR ℚ 0 c.1.1 c.1.2

noncomputable def qf (c : ShortIdx) (M : Matrix (Fin 24) (Fin 24) ℚ) : ℚ :=
  uQ c ⬝ᵥ (M *ᵥ uQ c)

noncomputable def tau (d e c : ShortIdx) : ℚ :=
  open Classical in if c.1 = d.1 + e.1 then sgn (pair e.1.1 d.1.2) else 0
\end{leancode}
\begin{leancode}
noncomputable def gmulA (u v : SpaceB) : Matrix (Fin 24) (Fin 24) ℚ :=
  (2 : ℚ) • (Am u * Am v + Am v * Am u) + (8 : ℚ)⁻¹ • ∑ c, (xs u c * xs v c) • Umat c +
    (trace ((Nm u)ᵀ * Nm v) / 8) • (1 : Matrix (Fin 24) (Fin 24) ℚ) +
    (2 : ℚ)⁻¹ • ((Nm u)ᵀ * Nm v + (Nm v)ᵀ * Nm u)
\end{leancode}
\begin{leancode}
noncomputable def gmulX (u v : SpaceB) (c : ShortIdx) : ℚ :=
  (xs v c * qf c (Am u) + xs u c * qf c (Am v)) / 8 +
    ∑ d, ∑ e, xs u d * xs v e * tau d e c +
    (trace ((Nm u)ᵀ * Rmat c * Nm v) - qf c ((Nm u)ᵀ * Rmat c * Nm v) / 4) / 8

noncomputable def gmulWHalf (u v : SpaceB) : SpaceW :=
  (trace (Am u) / 8) • Nm v + Nm v * Am u +
    (8 : ℚ)⁻¹ • ∑ c, xs u c • (Rmat c * Nm v - (4 : ℚ)⁻¹ • (Rmat c * Nm v * Umat c))

noncomputable def gmulW (u v : SpaceB) : SpaceW := gmulWHalf u v + gmulWHalf v u

noncomputable def gmulFun (u v : SpaceB) : SpaceB :=
  (⟨gmulA u v, gmulA_symm u v⟩, gmulX u v, gmulW u v)
\end{leancode}
\begin{leancode}
noncomputable def gmul : SpaceB →ₗ[ℚ] SpaceB →ₗ[ℚ] SpaceB :=
  LinearMap.mk₂ ℚ gmulFun gmulFun_add_left gmulFun_smul_left gmulFun_add_right
    gmulFun_smul_right
\end{leancode}
(\leanpath{VOA/\allowbreak Moonshine/\allowbreak MonsterGriess/\allowbreak Algebra.lean}.)

\begin{dictionary}
\lean{uQ c} is $u_c$ read in $\Q^{24}$; \lean{vecMulVec a b} is $ab^{\mathsf T}$, \lean{M *ᵥ v} is $Mv$
and \lean{a ⬝ᵥ b} the dot product, so \lean{Umat c} $=U_c$ and \lean{qf c M} $=q_c(M)$.  \lean{c.1} is the
pair $(a_c,\delta_c)$; \lean{sgn : ZMod 2 → ℚ} is $t\mapsto(-1)^t$.  \lean{Am u}, \lean{xs u}, \lean{Nm u}
are the three components of $u$; \lean{(8 : ℚ)⁻¹ •} is multiplication by $\frac18$; \lean{gmulA_symm}
proves that the $\chXXISym$-component is symmetric.  Setting two of the three components of $u$ and of
$v$ to zero in \lean{gmulA}, \lean{gmulX}, \lean{gmulW} gives exactly the rows of the table; the file
\leanpath{MoonshineCNGriessTable.lean} proves each row as a separate theorem, for instance
\end{dictionary}
\begin{leancode}
theorem gmul_spinBlock_spinBlock_snd_fst (N M : SpaceW) :
    (gmul (spinBlock N) (spinBlock M)).2.1
      = fun c => (trace (Nᵀ * Rmat c * M) - qf c (Nᵀ * Rmat c * M) / 4) / 8 := by
\end{leancode}
(\leanpath{VOA/\allowbreak Moonshine/\allowbreak MoonshineCNGriessTable.lean}, where
\lean{spinBlock N} $=(0,0,N)$, and similarly \lean{symBlock A} $=(A,0,0)$, \lean{shortBlock x} $=(0,x,0)$;
\lean{.2.1} is the $X$-component.)
\begin{dictionary}
\textbf{The port adaptation in \lean{gmul}.}  The source builds \lean{gmul} with the same
\lean{LinearMap.mk₂ ℚ gmulFun …}, but proves linearity in the second argument from commutativity.  The
port proves right linearity directly.  The data argument \lean{gmulFun} is identical; only proofs of
propositions differ, and by proof irrelevance the two \lean{gmul} are the same bilinear map
(\S\ref{ch21:sec:fidelityaudit}).

\textbf{Where the coefficients come from.}  The factors $2$, $\frac18$, $\frac14$ are those of Conway's
normalisation.  Their origin is clearest on the vertex-algebra side (Part~B): Lean coordinates of lattice
vectors are $\sqrt8$ times orthonormal coordinates, so a quadratic expression $y^{\mathsf T}Ay$ in
lattice coordinates is $8$ times the orthonormal one, which produces the $\frac18$'s; the $2$ in
$2(AA'+A'A)$ is the double Wick contraction of two boson bilinears; and $\frac{\tr A}8$ is the
ground-state energy $\tfrac{24}{16}$ of the twisted sector when $A=\tfrac12I$.  Theorems
\ref{ch21:thm:conway} and \ref{ch21:thm:twistedTop} prove that the vertex-algebra products have exactly
these coefficients.
\end{dictionary}

\begin{example}[Two worked products]\label{ch21:ex:products}
(1) $e_c\cdot e_c=(\tfrac18U_c,0,0)$: the $X$-part is $\tau(c,c,\cdot)$, supported at $c+c=0$, which is
not short.  (2) Put $P_c=\tfrac18U_c$.  Then $P_c\cdot P_c=4P_c^2=\tfrac4{64}U_c^2=\tfrac{4\cdot32}{64}U_c=16P_c$
and $P_c\cdot e_c=\tfrac18q_c(P_c)e_c=\tfrac18\cdot\tfrac{32^2}8e_c=16e_c$.  So $P_c$ is, up to the factor
$16$, an idempotent with $e_c$ in its eigenvalue-$16$ eigenspace.  Under the dictionary of Part~B, $P_c$
is the state $\alpha(-1)^2\mathbb1$ (Theorem~\ref{ch21:thm:conway}(c)), and indeed
$(\alpha(-1)^2\mathbb1)_{(1)}(\alpha(-1)^2\mathbb1)=4\langle\alpha,\alpha\rangle\,\alpha(-1)^2\mathbb1=16\,\alpha(-1)^2\mathbb1$
in the Heisenberg VOA.
\end{example}

\begin{remark}[Identity element]\label{ch21:rem:unit}
$e=(\tfrac14I,0,0)$ (Lean \lean{oneB}) is the identity: $2(\tfrac14A+\tfrac14A)=A$;
$\tfrac18q_c(\tfrac14I)=\tfrac18\cdot\tfrac14\cdot32=1$; and $\tfrac{\tr(I/4)}8N+\tfrac14N=\tfrac68N+\tfrac14N=N$.
This is a Lean theorem, and so is the identification of its image under the dictionary of Part~B with
$\omega/2$: with $e_1=(I,0,0)=4e$,
\begin{leancode}
theorem gmul_oneB_left (u : SpaceB) : gmul oneB u = u := by
\end{leancode}
\begin{leancode}
theorem griessNat_griessIdentityVec :
    griessNat griessIdentityVec = (2 : ℂ) • vnatConformalTwo := by
\end{leancode}
(\leanpath{VOA/\allowbreak Moonshine/\allowbreak MoonshineCNRationalGriess.lean}, a file documented in
Chapter~\ref{ch:22}; \lean{griessIdentityVec} $=e_1$, \lean{vnatConformalTwo} is $\omega$ as an element of
$\Vnat_2$, \lean{griessNat} is the dictionary of Definition~\ref{ch21:def:griessNat}.)  So
$\mathrm{griessNat}(e)=\omega/2$, and $2b(e,e)=3=\langle\omega/2,\omega/2\rangle$.
\end{remark}

\begin{physics}
The table lists the $(z-w)^{-2}$ OPE coefficients of the spin-two fields in a basis adapted to the
orbifold: boson bilinears close on themselves; a bilinear acts diagonally on $e^{i\alpha X}+e^{-i\alpha X}$
with eigenvalue $\propto\alpha^{\mathsf T}A\alpha$ (the change of its conformal weight under a change of
metric); two exponentials fuse to the exponential of the summed momentum when
$\langle\alpha,\beta\rangle=-2$, with the cocycle sign $\tau$, and to $\alpha(-1)^2$ when $\beta=-\alpha$;
untwisted$\times$twisted $\to$ twisted and twisted$\times$twisted $\to$ untwisted (the orbifold fusion
rules).  The twisted-twisted coefficients $R_c$ are matrix elements of $e^{\alpha}$ between twist ground
states.
\end{physics}
\begin{status}
Definitions, verbatim up to the documented proof-term difference in \lean{gmul}; the table is audited
in \leanpath{MoonshineCNGriessAudit.lean}.  \lean{gmul_oneB_left} is audited in
\leanpath{MoonshineCNRationalAudit.lean}.
\end{status}

\subsection{The trilinear form and its uniqueness}

\begin{definition}[Griess trilinear form]\label{ch21:def:T}
$\chXXIT(u,v,w)=b(u\cdot v,w)$.
\end{definition}
\begin{leancode}
noncomputable def T : SpaceB →ₗ[ℚ] SpaceB →ₗ[ℚ] SpaceB →ₗ[ℚ] ℚ := gmul.compr₂ bform
\end{leancode}
(\leanpath{VOA/\allowbreak Moonshine/\allowbreak MonsterGriess/\allowbreak GriessTrilinear.lean}.)
\begin{dictionary}
\lean{gmul.compr₂ bform} is $(u,v)\mapsto b(u\cdot v,\cdot)$.  The file proves that $\chXXIT$ is fully
symmetric (\lean{T_swap12}, \lean{T_cycle}, \lean{T_swap23}) through an explicit formula; full symmetry
says at once that the product is commutative (by nondegeneracy of $b$) and that $b$ is associative,
$b(u\cdot v,w)=b(u,v\cdot w)$.
\end{dictionary}
\begin{physics}
$\chXXIT(u,v,w)$ is the three-point function $\langle u(z_1)v(z_2)w(z_3)\rangle$ of three spin-two
fields; its total symmetry is the Bose symmetry of a correlator of bosonic fields.
\end{physics}
\begin{status}
No hypotheses; audited in \leanpath{ParityAudit.lean}.
\end{status}

\begin{theorem}[Uniqueness of the Griess trilinear form]\label{ch21:thm:Tunique}
Let $\chXXIT'$ be a trilinear form on $\chXXIB$, symmetric in its first two and in its last two
arguments, with the following values.  For $A,A',A''\in\chXXISym$; short $r,s,t\in Q$ (with
$c=\mathrm{coord}(r)$); $q,q'\in\Q^I$ and $v,w\in\Q^{24}$ (so $q\otimes v=qv^{\mathsf T}\in W$):
\begin{center}\small
\begin{tabular}{ll}
\toprule
arguments & value of $\chXXIT'$\\\midrule
$(A,A',A'')$ & $4\tr(AA'A'')$\\
$(\mathrm{Xlift}\,r,\mathrm{Xlift}\,s,A)$ & $(-1)^{\epsilon}q_c(A)/8$ if $rs=(-1)^\epsilon$; $0$ if $rs\notin\{\pm1\}$\\
$(\mathrm{Xlift}\,r,\mathrm{Xlift}\,s,\mathrm{Xlift}\,t)$ & $(-1)^\epsilon$ if $rst=(-1)^\epsilon$; $0$ if $rst\notin\{\pm1\}$\\
$(q\otimes v,q'\otimes w,A)$ & $(q\cdot q')\bigl(\tfrac{(v\cdot w)\tr A}8+v^{\mathsf T}Aw\bigr)$\\
$(q\otimes v,q'\otimes w,\mathrm{Xlift}\,r)$ & $(qr\cdot q')\bigl(v\cdot w-\tfrac14(v\cdot u_c)(w\cdot u_c)\bigr)/8$\\
$(A,A',x)$, $(A,A',N)$, $(A,x,N)$, $(x,x',N)$, $(N,N',N'')$ & $0$\\
\bottomrule
\end{tabular}
\end{center}
Then $\chXXIT'=\chXXIT$.
\end{theorem}
\begin{leancode}
theorem T_unique (T' : SpaceB →ₗ[ℚ] SpaceB →ₗ[ℚ] SpaceB →ₗ[ℚ] ℚ)
    (h12 : ∀ u v w, T' u v w = T' v u w) (h23 : ∀ u v w, T' u v w = T' u w v)
    (hT' : IsGriessForm T') : T' = T := by
\end{leancode}
\begin{leancode}
def IsGriessForm (T' : SpaceB →ₗ[ℚ] SpaceB →ₗ[ℚ] SpaceB →ₗ[ℚ] ℚ) : Prop :=
  (∀ A A' A'' : symMat, T' (inA A) (inA A') (inA A'') =
    4 * trace ((A : Matrix (Fin 24) (Fin 24) ℚ) * (A' : Matrix (Fin 24) (Fin 24) ℚ) *
      (A'' : Matrix (Fin 24) (Fin 24) ℚ))) ∧
  (∀ (r s : Q) (hr : IsShortCoord (coordOf r)) (_ : IsShortCoord (coordOf s)) (A : symMat),
    (∀ e : ZMod 2, r * s = zQ ^ e.val →
      T' (inX (Xlift r)) (inX (Xlift s)) (inA A) = sgn e * qf ⟨coordOf r, hr⟩ A / 8) ∧
    (r * s ∉ zpowers zQ → T' (inX (Xlift r)) (inX (Xlift s)) (inA A) = 0)) ∧
  (∀ r s t : Q, IsShortCoord (coordOf r) → IsShortCoord (coordOf s) →
    IsShortCoord (coordOf t) →
    (∀ e : ZMod 2, r * s * t = zQ ^ e.val →
      T' (inX (Xlift r)) (inX (Xlift s)) (inX (Xlift t)) = sgn e) ∧
    (r * s * t ∉ zpowers zQ → T' (inX (Xlift r)) (inX (Xlift s)) (inX (Xlift t)) = 0)) ∧
  (∀ (q q' : SpinIdx → ℚ) (v w : Fin 24 → ℚ) (A : symMat),
    T' (inW (tensor q v)) (inW (tensor q' w)) (inA A) =
      (q ⬝ᵥ q') * ((v ⬝ᵥ w) * trace (A : Matrix (Fin 24) (Fin 24) ℚ) / 8 +
        v ⬝ᵥ ((A : Matrix (Fin 24) (Fin 24) ℚ) *ᵥ w))) ∧
  (∀ (q q' : SpinIdx → ℚ) (v w : Fin 24 → ℚ) (r : Q) (hr : IsShortCoord (coordOf r)),
    T' (inW (tensor q v)) (inW (tensor q' w)) (inX (Xlift r)) =
      ((q ᵥ* ((r : GL SpinIdx ℚ) : Matrix SpinIdx SpinIdx ℚ)) ⬝ᵥ q') *
        ((v ⬝ᵥ w) - (v ⬝ᵥ uQ ⟨coordOf r, hr⟩) * (w ⬝ᵥ uQ ⟨coordOf r, hr⟩) / 4) / 8) ∧
  (∀ (A A' : symMat) (x : SpaceX), T' (inA A) (inA A') (inX x) = 0) ∧
  (∀ (A A' : symMat) (N : SpaceW), T' (inA A) (inA A') (inW N) = 0) ∧
  (∀ (A : symMat) (x : SpaceX) (N : SpaceW), T' (inA A) (inX x) (inW N) = 0) ∧
  (∀ (x x' : SpaceX) (N : SpaceW), T' (inX x) (inX x') (inW N) = 0) ∧
  (∀ N N' N'' : SpaceW, T' (inW N) (inW N') (inW N'') = 0)
\end{leancode}
(\leanpath{VOA/\allowbreak Moonshine/\allowbreak MonsterGriess/\allowbreak GriessTrilinear.lean}.)
\begin{dictionary}
\lean{inA A}, \lean{inX x}, \lean{inW N} are $(A,0,0)$, $(0,x,0)$, $(0,0,N)$; \lean{tensor q v} is
$qv^{\mathsf T}$; \lean{Xlift}, \lean{coordOf} are Definition~\ref{ch21:def:Xlift}; \lean{zQ} $=-1\in Q$
and \lean{zQ ^ e.val} $=(-1)^e$; \lean{zpowers zQ} $=\{\pm1\}$; \lean{q ᵥ* r} is the row vector $qr$.  The
ten conjuncts are exactly the ten rows of the table.  The hypotheses of \lean{T_unique} are precisely
``symmetric'' and ``has these values''; nothing else is assumed.
\end{dictionary}
\begin{proof}[Proof on paper]
Matrices span $\chXXISym$; the $\mathrm{Xlift}(\mathrm{rep}\,c)=e_c$ span $X$; the $q\otimes v$ span $W$.
A trilinear form symmetric under all permutations (the two transpositions generate $S_3$) is determined
by its values on unordered triples from a spanning set, and the table prescribes a value for every
unordered triple of block types.
\end{proof}
\begin{physics}
A CFT three-point function of primaries is fixed by its values on a basis; this lemma is the
finite-dimensional bookkeeping that lets one prove ``$g$ is a symmetry'' by checking that $g$ maps the
three-point couplings of a basis to themselves.  This is how $2^{1+24}.\mathrm{Co}_1$ is classically
shown to preserve the algebra.
\end{physics}
\begin{status}
No hypotheses beyond those displayed; audited in \leanpath{ParityAudit.lean}.
\end{status}

%======================================================================
\section{The automorphism group \texorpdfstring{\texttt{AutO}}{AutO}: AxiomMath's Monster}\label{ch21:sec:auto}

\begin{definition}[Algebra automorphisms; \lean{AutO}]\label{ch21:def:AutO}
A $\Q$-linear automorphism $f$ of $\chXXIB$ is an \emph{algebra automorphism} if
$f(u\cdot v)=f(u)\cdot f(v)$ for all $u,v$.  \lean{AutO} is the subgroup of $GL_\Q(\chXXIB)$ of algebra
automorphisms preserving $b$.
\end{definition}
\begin{leancode}
def IsAlgAut (f : SpaceB ≃ₗ[ℚ] SpaceB) : Prop := ∀ u v, f (gmul u v) = gmul (f u) (f v)
\end{leancode}
\begin{leancode}
def AutO : Subgroup (SpaceB ≃ₗ[ℚ] SpaceB) where
  carrier := {g | IsAlgAut g ∧ ∀ u v, bform (g u) (g v) = bform u v}
\end{leancode}
\begin{leancode}
theorem mem_AutO {g : SpaceB ≃ₗ[ℚ] SpaceB} :
    g ∈ AutO ↔ (∀ u v, g (gmul u v) = gmul (g u) (g v)) ∧ ∀ u v, bform (g u) (g v) = bform u v :=
  Iff.rfl
\end{leancode}
(\leanpath{VOA/\allowbreak Moonshine/\allowbreak MonsterGriess/\allowbreak AutO.lean},
\lean{VOA.Moonshine.MonsterGriess.AutO}.)

\begin{dictionary}
\lean{SpaceB ≃ₗ[ℚ] SpaceB} is the group of $\Q$-linear bijections of $\chXXIB$ under composition.
\lean{Subgroup} bundles a carrier with closure proofs (omitted).  \lean{mem_AutO} unfolds membership by
\lean{Iff.rfl}: the definition contains nothing beyond the two displayed conditions.

\textbf{Is this ``the Monster''?}  In AxiomMath's library the Monster is \emph{defined} as this group
(\lean{MGriess := ↥CFSG.Monster.Griess.AutO}), and AxiomMath prove there, in Lean~4.35, their formal
Monster-existence statement for it.  None of that is ported.  In LeanVOA, \lean{AutO} is ``the
Monster'' \emph{by reference}: the definition is AxiomMath's verbatim (\S\ref{ch21:sec:fidelityaudit}),
and its identification with $\mathbb M$ is a theorem of AxiomMath's library and of the classical
literature (\S\ref{ch21:sec:tits}), not of LeanVOA.  The book calls it ``\lean{AutO}'' or ``AxiomMath's
Monster''.

\textbf{Could \lean{AutO} be trivial or degenerate?}  No: it contains $G_{\mathrm{vis}}$ and $\tau_B$
(Theorems~\ref{ch21:thm:gvis} and~\ref{ch21:thm:tauBgmul}), which contain explicit nontrivial elements such
as $x_B$.  Finiteness and the order of \lean{AutO} are not LeanVOA theorems; Chapter~\ref{ch:22} needs only
that \lean{AutO} is generated by $G_{\mathrm{vis}}$ and $\tau_B$, and states that as a hypothesis.
\end{dictionary}
\begin{physics}
\lean{AutO} is the group of all linear maps of the $196884$ spin-two fields preserving their OPE and
two-point function, i.e.\ the largest group that could act as a symmetry of the Monster CFT and be visible
on this sector.  Classically (Griess, Conway, Tits, Frenkel--Lepowsky--Meurman) it \emph{is} the global
symmetry group of the CFT; in Lean the corresponding statement $\Aut(\Vnat)\cong$ \lean{AutO} is proved in
Chapter~\ref{ch:22} only under the hypothesis \lean{GmatEqAutO}.
\end{physics}
\begin{status}
Definition, verbatim from the source; it occurs in statements audited in
\leanpath{ProductInvarianceAudit.lean} and \leanpath{ParityAudit.lean}.
\end{status}

%======================================================================
\section{The visible subgroup \texorpdfstring{$G_{\mathrm{vis}}$}{Gvis} (classically \texorpdfstring{$2^{1+24}.\mathrm{Co}_1$}{2\^{}(1+24).Co1})}\label{ch21:sec:gvis}

\subsection{\texorpdfstring{$\mathrm{Co}_0$}{Co0} as rational matrices}

\begin{definition}[$\mathrm{Co}_0$]\label{ch21:def:latticeO}
\lean{latticeO} is the group of rational orthogonal $24\times24$ matrices $g$ ($gg^{\mathsf T}=1$) such
that $x\mapsto xg$ maps the rational points of $\Lambda(\chXXIC)$ onto themselves.
\end{definition}
\begin{leancode}
def latticeO : Subgroup (GL (Fin 24) ℚ) where
  carrier := {g | (g : Matrix (Fin 24) (Fin 24) ℚ) ∈ orthogonalGroup (Fin 24) ℚ ∧
    (fun x => x ᵥ* (g : Matrix (Fin 24) (Fin 24) ℚ)) '' latticeQ = latticeQ}
\end{leancode}
\begin{leancode}
theorem mem_ker_modTwoHom_iff (g : latticeO) : g ∈ modTwoHom.ker ↔ mat g = 1 ∨ mat g = -1 :=
\end{leancode}
(\leanpath{VOA/\allowbreak Moonshine/\allowbreak MonsterGriess/\allowbreak ScalarKernel.lean}, namespace \lean{.ScalarKernel}.)
\begin{dictionary}
\lean{x ᵥ* g} is the row vector $xg$; \lean{latticeQ} is the image of \lean{codeLattice C} in
$\Q^{24}$; \lean{orthogonalGroup} is Mathlib's $\{g: gg^{\mathsf T}=1\}$ (equivalently
$g^{\mathsf T}g=1$).  Since orthogonality is for the standard dot product and the lattice form is
$x\cdot y/8$, \lean{latticeO} is exactly $O(\Leech)=\mathrm{Co}_0$ (in these coordinates).
\lean{modTwoEnd g} is the induced map on $\chXXIL$ and \lean{modTwoHom} the resulting homomorphism to
$GL(\chXXIL)^{\mathrm{op}}$; the theorem says its kernel is $\{\pm1\}$, the classical fact
$\mathrm{Co}_0\to\mathrm{Co}_1=\mathrm{Co}_0/\{\pm1\}\hookrightarrow GL(\Leech/2\Leech)$.
\end{dictionary}
\begin{physics}
$\mathrm{Co}_0$ is the group of rotations of the $24$ bosons preserving the momentum lattice: the
symmetry group of the Leech lattice CFT that commutes with the $U(1)^{24}$ currents up to rotation.
Only $\mathrm{Co}_1=\mathrm{Co}_0/\{\pm1\}$ survives in the Monster, because $-1$ is the orbifold
reflection, which acts trivially on the orbifold theory.
\end{physics}
\begin{status}
Both audited in \leanpath{ParityAudit.lean}.
\end{status}

\subsection{Visible pairs and their action}

\begin{definition}[Visible pairs]\label{ch21:def:Pvis}
A \emph{visible pair} is $(s,g)$ with $s\in GL_I(\Q)$ orthogonal, with dyadic entries (in $\Z[\tfrac12]$),
normalising $Q$, and $g\in\mathrm{Co}_0$ \emph{compatible} with $s$: for all $r,r'\in Q$ with
$r'=s^{-1}rs$, the class of $r'$ in $Q/Z(Q)\cong\chXXIL$ is $g$ applied to the class of $r$.
\end{definition}
\begin{leancode}
def IsDyadicSpinNormalizer (s : GL SpinIdx ℚ) : Prop :=
  (s : Matrix SpinIdx SpinIdx ℚ) ∈ orthogonalGroup SpinIdx ℚ ∧
    (∀ i j, (s : Matrix SpinIdx SpinIdx ℚ) i j ∈ dyadic) ∧
    (fun r => s⁻¹ * r * s) '' (Q : Set (GL SpinIdx ℚ)) = Q
\end{leancode}
\begin{leancode}
def IsCompatible (s : GL SpinIdx ℚ) (g : latticeO) : Prop :=
  ∀ r r' : Q, (r' : GL SpinIdx ℚ) = s⁻¹ * r * s →
    Multiplicative.toAdd (phi (QuotientGroup.mk r')) =
      modTwoEnd g (Multiplicative.toAdd (phi (QuotientGroup.mk r)))
\end{leancode}
\begin{leancode}
def Pvis : Subgroup (GL SpinIdx ℚ × latticeO) where
  carrier := {p | IsDyadicSpinNormalizer p.1 ∧ IsCompatible p.1 p.2}
\end{leancode}
(\leanpath{VOA/\allowbreak Moonshine/\allowbreak MonsterGriess/\allowbreak VisiblePair.lean}, namespace \lean{.Visible}.)
\begin{dictionary}
\lean{GL SpinIdx ℚ} is the group of invertible rational $I\times I$ matrices; \lean{dyadic} is the subring
$\Z[\tfrac12]\subset\Q$; \lean{(fun r => s⁻¹ * r * s) '' Q = Q} says conjugation by $s$ maps $Q$ onto $Q$.
In \lean{IsCompatible}, \lean{QuotientGroup.mk r} is the class of $r$ in $Q/Z(Q)$, \lean{phi} the
isomorphism $Q/Z(Q)\cong\chXXIL$ of \S\ref{ch21:sec:spin} (written multiplicatively, hence
\lean{Multiplicative.toAdd}), and \lean{modTwoEnd g} the action of $g$ on $\chXXIL$.  \lean{Pvis} is a
subgroup of the product group, so visible pairs compose componentwise.
\end{dictionary}
\begin{physics}
A visible pair is a rotation $g$ of the Leech lattice together with a compatible rotation $s$ of the
$4096$ twist ground states: the data needed to extend a lattice symmetry of the untwisted theory to the
twisted sector.
\end{physics}
\begin{status}
Definitions, verbatim from the source (\leanpath{VisiblePair.lean} is ported in full); \lean{Pvis}
audited in \leanpath{ParityAudit.lean}.
\end{status}

\begin{definition}[The action $\theta_p$ and $G_{\mathrm{vis}}$]\label{ch21:def:Gvis}
For $p=(s,g)$: $\theta_p(A,x,N)=(g^{\mathsf T}Ag,\ \theta^X_p(x),\ s^{\mathsf T}Ng)$, where
$\theta^X_p(e_c)=\mathrm{Xlift}(s^{-1}\,\mathrm{rep}(c)\,s)$ ($=\pm e_{c'}$ with $c'$ the class of
$g$ applied to $c$).  $G_{\mathrm{vis}}=\{\theta_p:p\text{ visible}\}$.
\end{definition}
\begin{leancode}
def thetaA (g : Matrix (Fin 24) (Fin 24) ℚ) : symMat →ₗ[ℚ] symMat where
  toFun A := ⟨gᵀ * (A : Matrix (Fin 24) (Fin 24) ℚ) * g, by
\end{leancode}
\begin{leancode}
def thetaW (s : Matrix SpinIdx SpinIdx ℚ) (g : Matrix (Fin 24) (Fin 24) ℚ) :
    SpaceW →ₗ[ℚ] SpaceW where
  toFun N := sᵀ * N * g
\end{leancode}
\begin{leancode}
def thetaX (p : Pvis) : SpaceX →ₗ[ℚ] SpaceX :=
  (Pi.basisFun ℚ ShortIdx).constr ℚ fun c => Xlift (conjQ p (rep c.1))
\end{leancode}
\begin{leancode}
def theta (p : Pvis) : SpaceB →ₗ[ℚ] SpaceB :=
  (thetaA (mat p.1.2)).prodMap
    ((thetaX p).prodMap (thetaW (p.1.1 : Matrix SpinIdx SpinIdx ℚ) (mat p.1.2)))
\end{leancode}
\begin{leancode}
def Gvis : Subgroup (SpaceB ≃ₗ[ℚ] SpaceB) where
  carrier := Set.range thetaEquiv
\end{leancode}
(\leanpath{VOA/\allowbreak Moonshine/\allowbreak MonsterGriess/\allowbreak VisiblePair.lean}; \lean{thetaEquiv p} is \lean{theta p} with
inverse \lean{theta p⁻¹}.)
\begin{dictionary}
\lean{conjQ p} is conjugation $r\mapsto s^{-1}rs$ as an automorphism of $Q$ (well defined because $s$
normalises $Q$).  \lean{Pi.basisFun … .constr} defines a linear map by its values on the standard basis
$e_c$.  \lean{prodMap} acts blockwise.  \lean{mat g} is the matrix of $g\in$ \lean{latticeO}.
Note the conventions: lattice vectors are \emph{row} vectors acted on by $x\mapsto xg$, hence
$A\mapsto g^{\mathsf T}Ag$ on $S$ and $N\mapsto s^{\mathsf T}Ng$ on $W$.  \lean{thetaEquiv} is
anti-multiplicative (\lean{thetaEquiv_mul : thetaEquiv (p * q) = thetaEquiv q * thetaEquiv p}), which
is why the group homomorphism \lean{thetaHom} lands in the opposite group; $G_{\mathrm{vis}}$ is its
image (\lean{Gvis_eq_unop_range}).  Nothing ``opposite'' survives in $G_{\mathrm{vis}}$ itself: it
is an ordinary subgroup of $GL_\Q(\chXXIB)$, namely the set of all $\theta_p$.
\end{dictionary}

\begin{example}[Visible pairs]\label{ch21:ex:visible}
(1) For $r\in Q$, $(r,1)$ is visible: $r$ is a signed permutation matrix (orthogonal, entries $0,\pm1$),
it normalises $Q$, and conjugation by $r$ is trivial on $Q/Z(Q)$, which is compatible with $g=1$.  Then
$\theta_{(r,1)}(A,x,N)=(A,\theta^X(x),r^{\mathsf T}N)$ with $\theta^X(e_c)=\pm e_c$, the sign being the
commutator of $r$ with $\mathrm{rep}(c)$; this is \lean{rhoEquiv r}.  For $r=z=-1$ it is
$x_B:(A,x,N)\mapsto(A,x,-N)$.  (2) $(1,-1)$ is visible ($-1\in\mathrm{Co}_0$ acts trivially on
$\chXXIL$) and $\theta_{(1,-1)}(A,x,N)=(A,x,-N)$ as well: $(1,-1)=(-1,1)\cdot(-1,-1)$, and $(-1,-1)$ acts
trivially.  (3) A coordinate permutation $g$ in $M_{24}\subset\mathrm{Co}_0$ needs a compatible $s$, a
$4096\times4096$ matrix permuting (up to sign) the basis of $\Q^I$ in a way that depends on how $g$ moves
the code; such $s$ exist classically, but no specific one is written down in this chapter.
\end{example}

\paragraph{Which elements of $\mathrm{Co}_0$ occur.}  Classically every $g\in\mathrm{Co}_0$ occurs in a
visible pair, and then $|G_{\mathrm{vis}}|=|Q|\cdot|\mathrm{Co}_0|/2=2^{25}|\mathrm{Co}_1|$, which is why
$G_{\mathrm{vis}}$ is ``$2^{1+24}.\mathrm{Co}_1$''.  The library does \emph{not} prove this: it proves no
surjectivity of $\mathrm{Pvis}\to\mathrm{Co}_0$ and no order of $G_{\mathrm{vis}}$.  Nothing in this chapter
or in Chapter~\ref{ch:22} needs it: $G_{\mathrm{vis}}$ is \emph{defined} as the set of all $\theta_p$, and
the statements about it (it lies in \lean{AutO}; it is the centraliser of $x_B$; each of its elements lifts
to $\Vnat$) hold for whatever it is.
\begin{physics}
$\theta_p$ acts on the three kinds of spin-two fields: by rotating the bosons in the bilinears, by
permuting (with signs) the exponentials $\cos(\alpha X)$, and by rotating both the oscillator index and
the ground-state index of the twisted fields.
\end{physics}
\begin{status}
Definitions, verbatim; \lean{thetaEquiv}, \lean{Gvis} audited in \leanpath{ParityAudit.lean}.
\end{status}

\begin{theorem}[$G_{\mathrm{vis}}$ preserves the algebra; kernel; $x_B$]\label{ch21:thm:gvis}
(a) Every $f\in G_{\mathrm{vis}}$ preserves the product and the form; hence $G_{\mathrm{vis}}\le$ \lean{AutO}.
(b) $\theta_p=1$ iff $p=1$ or $p=(-1,-1)$.
(c) For the central element $z=-1\in Q$, $\theta_{(z,1)}=x_B$, the parity involution
$x_B(A,x,N)=(A,x,-N)$.
(d) $Q$ is generated by its short elements.
\end{theorem}
\begin{leancode}
theorem mem_Gvis_gmul {f : SpaceB ≃ₗ[ℚ] SpaceB} (hf : f ∈ Gvis) (u v : SpaceB) :
    f (gmul u v) = gmul (f u) (f v) ∧ bform (f u) (f v) = bform u v := by
\end{leancode}
\begin{leancode}
theorem Gvis_le_AutO : Gvis ≤ AutO := fun _ hf =>
\end{leancode}
\begin{leancode}
theorem mem_ker_thetaHom_iff (p : Pvis) : p ∈ thetaHom.ker ↔ p = 1 ∨ p = negPair := by
\end{leancode}
\begin{leancode}
theorem rhoEquiv_zQ : rhoEquiv zQ = xB := by
\end{leancode}
\begin{leancode}
theorem closure_isShortCoord_eq_top : closure {r : Q | IsShortCoord (coordOf r)} = ⊤ := by
\end{leancode}
(\leanpath{VOA/\allowbreak Moonshine/\allowbreak MonsterGriess/\allowbreak VisibleAction.lean}, namespace \lean{.Visible}.)  The parity
involution itself:
\begin{leancode}
noncomputable def parityMap : SpaceB →ₗ[ℚ] SpaceB :=
  LinearMap.prodMap LinearMap.id (LinearMap.prodMap LinearMap.id (-LinearMap.id))
\end{leancode}
\begin{leancode}
noncomputable def xB : SpaceB ≃ₗ[ℚ] SpaceB :=
  LinearEquiv.ofInvolutive parityMap fun u => by simp [parityMap_apply]
\end{leancode}
(\leanpath{VOA/\allowbreak Moonshine/\allowbreak MonsterGriess/\allowbreak GriessTrilinear.lean}.)
\begin{dictionary}
\lean{negPair} is the pair $(-1,-1)\in$ \lean{Pvis} (\lean{zU} $=-1$, \lean{negOneO} $=-1\in\mathrm{Co}_0$).
\lean{rhoEquiv r} $=\theta_{(r,1)}$, the image of $r\in Q$ (``Axiom's extraspecial generators'').
\lean{LinearEquiv.ofInvolutive} turns an involutive linear map into a linear equivalence.
In (a) the conclusion is exactly membership in \lean{AutO} (Definition~\ref{ch21:def:AutO}).
\end{dictionary}
\begin{proof}[Proof on paper]
(a) Let $\chXXIT_p(u,v,w)=\chXXIT(\theta_pu,\theta_pv,\theta_pw)$.  It is symmetric, and its block values are
those of $\chXXIT$: $g^{\mathsf T}(\cdot)g$ with $g$ orthogonal preserves $\tr(AA'A'')$; conjugation by $s$
permutes the short elements of $Q$ compatibly with $g$, so it preserves the $XXA$, $XXX$ values (the
products $rs$, $rst$ and the norms $q_c$ transform covariantly by compatibility); on $W$,
$s^{\mathsf T}(q\otimes v)g=(qs)\otimes(vg)$ and the $WWX$ value
$(qr\cdot q')(\dots)$ is preserved because $s$ is orthogonal and conjugates $r$ to $s^{-1}rs$.  By
Theorem~\ref{ch21:thm:Tunique}, $\chXXIT_p=\chXXIT$, i.e.\ $b(\theta_p(u\cdot v),\theta_pw)=b(\theta_pu\cdot\theta_pv,\theta_pw)$;
nondegeneracy of $b$ and surjectivity of $\theta_p$ give multiplicativity, and preservation of $b$ is
the orthogonality of $g$ and $s$.  (b) If $\theta_p=1$ then $g^{\mathsf T}Ag=A$ for all symmetric $A$
forces $g=\pm1$, and $s^{\mathsf T}Ng=N$ for all $N$ forces $s=\pm1$ with the same sign.  (c) For
$p=(-1,1)$: $A$ and $x$ are fixed (conjugation by $-1$ is trivial) and $N\mapsto-N$.  (d) The short
classes span $\chXXIL$ (the minimal vectors generate $\Leech$, \lean{closure_latticeGens}), so the
short elements generate $Q$ modulo $Z(Q)=\Phi(Q)$, and elements generating modulo the Frattini subgroup
generate.
\end{proof}
\begin{physics}
$G_{\mathrm{vis}}$ is the part of the Monster that is \emph{manifest} in the orbifold construction:
the symmetries of the Leech lattice ($\mathrm{Co}_0$) together with the extraspecial group of
``momentum shifts'' $e^{\pi i\langle\mu,p\rangle}$ and their action on the twist ground states.  It
commutes with the orbifold involution $z$ (which is $x_B$: $+1$ on untwisted, $-1$ on twisted
fields), and classically (not formalised) it is the centraliser $C_{\mathbb M}(z)\cong2^{1+24}.\mathrm{Co}_1$
of a $2B$ involution of the Monster.  The kernel $\{1,(-1,-1)\}$ says that the reflection $X\mapsto-X$ combined with
$-1$ on the ground states acts trivially on the orbifold theory.
\end{physics}
\begin{status}
No hypotheses; audited in \leanpath{ParityAudit.lean}.
\end{status}

%======================================================================
\section{Triality: the element \texorpdfstring{$\tau_B$}{tauB}, the four-group and \texorpdfstring{$G_{\mathrm{mat}}$}{Gmat}}\label{ch21:sec:triality}

\subsection{The triality basis}

$G_{\mathrm{vis}}$ preserves the decomposition $\chXXIB=(\chXXISym\oplus X)\oplus W$ into the
$\pm1$-eigenspaces of $x_B$, i.e.\ into ``untwisted'' and ``twisted'' parts.  Classically the Monster is
generated by $G_{\mathrm{vis}}$ and one more element, which must mix the two.

\paragraph{Why a triality basis.}  Conway's idea: $x_B$ should be one of three commuting involutions
$x_B,y_B,z_B$ (with $1$, a four-group), conjugate under an element of order $3$ that permutes them
cyclically.  Each involution splits $\chXXIB$ into $\pm1$
eigenspaces; $x_B$ is $-1$ exactly on $W$ ($98304$-dimensional).  To permute the three involutions, one
needs a basis in which all three are diagonal, grouped into triples on which $x_B,y_B,z_B$ have sign
patterns $(1,-1,-1)$, $(-1,1,-1)$, $(-1,-1,1)$, plus vectors fixed by all three; the order-$3$ element
then cycles each triple.  Counting dimensions: $W$ splits as two halves $\eta=0,1$ of $49152$ each,
paired with the $49152$ classes of odd minimal vectors $(\mp3,\pm1^{23})$; the off-diagonal matrix
units $E_{ij}+E_{ji}$ are paired with the two classes $4e_i\pm4e_j$; and the remaining vectors (diagonal
units, octad classes) are fixed by the four-group, on which the order-$3$ element must act by something
else.

\begin{definition}[Triality basis]\label{ch21:def:tbasis}
The index set is
\begin{equation*}
  T_{\mathrm{idx}}=\{0..23\}\ \sqcup\ (\mathrm{LtPair}\times\{0,1,2\})\ \sqcup\ (\chXXIC_0\times\{0..23\}\times\{0,1,2\})\ \sqcup\ \mathrm{OctIdx},
\end{equation*}
of size $24+276\cdot3+2048\cdot24\cdot3+759\cdot64=196884$, where $\mathrm{LtPair}$ is the set of
pairs $i<j$ and $\mathrm{OctIdx}$ the pairs (octad $d$, even subset of $d$ modulo complement in $d$).
The basis vectors are: the diagonal matrix units $E_{ii}$; for each pair $i<j$ the triple
$(E_{ij}+E_{ji},\ e_{r^-}-e_{r^+},\ e_{r^-}+e_{r^+})$ where $r^\mp$ are the short classes of
$4e_i\mp4e_j$; for each odd minimal-vector class $(c,i)$ the triple $(e_{\mathrm{odd}(c,i)},\ \pm E_1(c)\otimes e_i,\ \pm E_0(c)\otimes e_i)$
with $E_\eta(c)$ the basis vector $(\eta,c)$ of $\Q^I$; and the vectors $e_c$ of the octad classes
(shape $(\pm2^8,0^{16})$).  The Gram matrix is diagonal with entries in $\{1,2\}$.
\end{definition}
\begin{leancode}
abbrev TIdx : Type := Fin 24 ⊕ (LtPair × Fin 3) ⊕ (C₀ × Fin 24 × Fin 3) ⊕ OctIdx
\end{leancode}
\begin{leancode}
theorem card_tIdx : Fintype.card TIdx = 196884 := by
\end{leancode}
\begin{leancode}
noncomputable def tbasis : Basis TIdx ℚ SpaceB :=
  basisOfLinearIndependentOfCardEqFinrank linearIndependent_tvec (card_tIdx.trans finrank_B.symm)
\end{leancode}
\begin{leancode}
theorem bform_tbasis (x y : TIdx) : bform (tbasis x) (tbasis y) = if x = y then tnorm x else 0 := by
\end{leancode}
(\leanpath{VOA/\allowbreak Moonshine/\allowbreak MonsterGriess/\allowbreak TrialityBasis.lean}.)
\begin{dictionary}
\lean{tbasis} is the basis whose vectors are \lean{tvec x}, made into a basis from linear independence
and the count $196884=\dim\chXXIB$ (\lean{coe_tbasis : ⇑tbasis = tvec}).  \lean{tnorm} is $2$ on the
pair triples and $1$ elsewhere.  The label bookkeeping (\leanpath{TrialityLabels.lean},
\leanpath{TrialityClasses.lean}) proves that the four label families \lean{rMinus}, \lean{rPlus},
\lean{rOdd}, \lean{rOct} of elements of $Q$ hit each short class exactly once
(\lean{labelEquiv : ShortLabel ≃ ShortIdx}, $552+49152+48576=98280$).  The octad part uses the
$6$-dimensional quadratic space of even subsets of an octad modulo complement
(\leanpath{OctadEvenSubset.lean}).
\end{dictionary}
\begin{physics}
The triality basis groups the spin-two fields into triples that are permuted by a hidden $\Z/3$
symmetry: off-diagonal oscillator bilinears with the two combinations of $\cos((4e_i\mp4e_j)X/\sqrt8)$;
odd-momentum exponentials with the two halves of the first excited twisted states.  This is the
algebraic shadow of the fact that the $\Z/2$ orbifold of the Leech CFT has a larger symmetry that mixes
twisted and untwisted sectors (FLM's triality, Chapter~\ref{ch:11}).
\end{physics}
\begin{status}
Audited in \leanpath{TrialityAudit.lean}.
\end{status}

\subsection{The element \texorpdfstring{$\tau_B$}{tauB}}

For an octad $d$ let $E_d$ be the $\chXXIF$-space of even subsets of $d$ modulo the complement in $d$
($\delta\sim d\setminus\delta$); it has dimension $6$ and $64$ elements, carries the quadratic form
$q_d(\delta)=\tfrac12|\delta|\bmod2$ with polar form $b_d(\delta,\varepsilon)=|\delta\cap\varepsilon|\bmod2$
(nondegenerate, with Gauss sum $\sum_\delta(-1)^{q_d(\delta)}=8$), and indexes the $64$ octad classes of
shape $(\pm2^8,0^{16})$ supported on $d$.  Put
\begin{equation*}
  (H_d)_{\delta\varepsilon}=\tfrac18(-1)^{b_d(\delta,\varepsilon)},\qquad
  D_d=\mathrm{diag}\bigl((-1)^{q_d(\delta)}\bigr)\qquad(64\times64).
\end{equation*}
Then $H_d^2=D_d^2=1$, $H_d,D_d$ are orthogonal, and $(H_dD_d)^3=1$ (\lean{binary_quadratic_fourier_triality}):
the ``Fourier'' and ``sign'' operators generate a copy of $S_3$, the finite analogue of $S$ and $T$ in
$SL_2(\Z)$ acting on theta functions.

\begin{definition}[$\tau_B$]\label{ch21:def:tauB}
$\tau_B$ is the linear map that fixes the $24$ diagonal units, sends the $k$-th vector of every triple to the
$(k+1)$-st (indices mod $3$), and sends the octad vector $v_{d,\delta}$ to
$\sum_\varepsilon(H_dD_d)_{\delta\varepsilon}v_{d,\varepsilon}$.  Equivalently: its matrix in \lean{tbasis} is the
transpose of the block-diagonal matrix \lean{tauMat} with blocks $1$, the $3$-cycle, and $H_dD_d$.
\end{definition}
\begin{leancode}
noncomputable def tauMat : Matrix TIdx TIdx ℚ :=
  diag2 1 (diag2 pairBlock (diag2 oddBlock octBlock))
\end{leancode}
\begin{leancode}
noncomputable def tauB : SpaceB ≃ₗ[ℚ] SpaceB :=
  Matrix.toLinOfInv tbasis tbasis tauMat_transpose_mul tauMat_mul_transpose
\end{leancode}
\begin{leancode}
theorem tauB_pow_three : tauB ^ 3 = 1 := by
\end{leancode}
\begin{leancode}
theorem orderOf_tauB : orderOf tauB = 3 := by
\end{leancode}
\begin{leancode}
theorem trace_tauB : LinearMap.trace ℚ SpaceB tauB = 783 := by
\end{leancode}
\begin{leancode}
theorem finrank_ker_tauB_sub_one :
    finrank ℚ (LinearMap.ker ((tauB : SpaceB →ₗ[ℚ] SpaceB) - 1)) = 66150 := by
\end{leancode}
\begin{leancode}
theorem bform_tauB (u v : SpaceB) : bform (tauB u) (tauB v) = bform u v :=
\end{leancode}
(\leanpath{VOA/\allowbreak Moonshine/\allowbreak MonsterGriess/\allowbreak TrialityMatrix.lean}.)
\begin{dictionary}
\lean{Matrix.toLinOfInv b b h₁ h₂} is the linear equivalence with matrix $M^{\mathsf T}$ in the basis
$b$ given proofs that $M^{\mathsf T}M=MM^{\mathsf T}=1$ (\lean{coe_tauB : … = toLin tbasis tbasis tauMatᵀ}).
\lean{diag2} forms block-diagonal matrices on a sum type.  The orthogonality proofs use that $H_d$,
$D_d$ are orthogonal with dyadic entries and $(H_dD_d)^3=1$ (\leanpath{OctadEvenSubset.lean}).
The action on basis vectors is \lean{tauB_diag}, \lean{tauB_pair}, \lean{tauB_odd}, \lean{tauB_oct}.

\textbf{The numbers.}  Trace: the diagonal units contribute $24$, each $3$-cycle contributes $0$, and
each octad block contributes $\tr(H_dD_d)=\tfrac18\sum_\delta(-1)^{b_d(\delta,\delta)+q_d(\delta)}=\tfrac18\sum_\delta(-1)^{q_d(\delta)}=1$
(as $b_d(\delta,\delta)=2q_d(\delta)=0$); so $\tr\tau_B=24+759=783$.  Fixed space: for an element of order $3$
with eigenvalues $1,\omega,\bar\omega$ of multiplicities $m_1,m,m$ one has $m_1+2m=196884$ and
$m_1-m=\tr\tau_B=783$, so $m_1=\tfrac13(196884+2\cdot783)=66150$.  These numbers are consistency checks
(the class $3A$ of $\mathbb M$ has character value $782$ on the $196883$-dimensional representation).
\end{dictionary}
\begin{physics}
$\tau_B$ is an order-three symmetry rotating ``oscillator bilinear $\to$ untwisted exponential $\to$
twisted state''.  Its trace $783=1+782$ on $\Vnat_2$ equals the coefficient of $q$ in the
McKay--Thompson series $T_{3A}=q^{-1}+783q+\cdots$, as it must if $\tau_B$ is in the class $3A$ of
the Monster (this identification is classical and not formalised).
\end{physics}
\begin{status}
No hypotheses; audited in \leanpath{TrialityAudit.lean}.
\end{status}

\begin{definition}[Four-group and the generated group]\label{ch21:def:Gmat}
$y_B=\tau_Bx_B\tau_B^{-1}$, $z_B=\tau_B^2x_B\tau_B^{-2}$;
$G_{\mathrm{mat}}=\langle G_{\mathrm{vis}}\cup\{\tau_B\}\rangle$.
\end{definition}
\begin{leancode}
noncomputable def yB : SpaceB ≃ₗ[ℚ] SpaceB := tauB * xB * tauB⁻¹
\end{leancode}
\begin{leancode}
noncomputable def zB : SpaceB ≃ₗ[ℚ] SpaceB := tauB ^ 2 * xB * (tauB ^ 2)⁻¹
\end{leancode}
\begin{leancode}
noncomputable def Gmat : Subgroup (SpaceB ≃ₗ[ℚ] SpaceB) :=
  Subgroup.closure ((Gvis : Set (SpaceB ≃ₗ[ℚ] SpaceB)) ∪ {tauB})
\end{leancode}
(\leanpath{VOA/\allowbreak Moonshine/\allowbreak MonsterGriess/\allowbreak TrialityFourGroup.lean}; \lean{yB}, \lean{zB} in
\lean{VOA.Moonshine.MonsterGriess}, \lean{Gmat} in \lean{VOA.Moonshine.MonsterGriess.Generated}.)
\begin{dictionary}
\lean{tauB⁻¹} is the inverse in the group \lean{SpaceB ≃ₗ[ℚ] SpaceB}; \lean{Subgroup.closure S} is the
subgroup generated by $S$; the coercion \lean{(Gvis : Set _)} views the subgroup as a set.  These are the
source's definitions verbatim, with the fully qualified names
\lean{CFSG.Monster.Generated.Gmat} $\mapsto$ \lean{VOA.Moonshine.MonsterGriess.Generated.Gmat}; it is
\emph{this} \lean{Gmat} that appears in \lean{GmatEqAutO} (Chapter~\ref{ch:22}).
\end{dictionary}
\begin{physics}
$G_{\mathrm{mat}}$ is the group generated by the manifest symmetries of the orbifold and the hidden
triality: classically the whole Monster.  That it is all of \lean{AutO} is AxiomMath's generation theorem.
\end{physics}
\begin{status}
Definitions; audited in \leanpath{TrialityAudit.lean}.
\end{status}
\begin{theorem}[Rational triality four-group]\label{ch21:thm:fourgroup}
In \lean{tbasis}, $x_B,y_B,z_B$ are diagonal with signs $(1,-1,-1)$, $(-1,1,-1)$, $(-1,-1,1)$ on the
three slots of each odd triple and $+1$ elsewhere; $\{1,x_B,y_B,z_B\}$ is a Klein four-group, $\tau_B$
permutes $x_B\mapsto y_B\mapsto z_B\mapsto x_B$ by conjugation, all three preserve $b$, and each has
a $98304$-dimensional $(-1)$-eigenspace.
\end{theorem}
\begin{leancode}
theorem rational_triality_four_group :
    yB = tauB * xB * tauB⁻¹ ∧ zB = tauB ^ 2 * xB * (tauB ^ 2)⁻¹ ∧
    (∀ t : TIdx, xB (tbasis t) = oddSign ![1, -1, -1] t • tbasis t) ∧
    (∀ t : TIdx, yB (tbasis t) = oddSign ![-1, 1, -1] t • tbasis t) ∧
    (∀ t : TIdx, zB (tbasis t) = oddSign ![-1, -1, 1] t • tbasis t) ∧
\end{leancode}
(the statement continues with the group relations, the eigenspace dimensions and the dyadic
lattice; \leanpath{TrialityFourGroup.lean}.)
\begin{dictionary}
\lean{oddSign v t} is $v_k$ if $t$ is slot $k$ of an odd triple and $1$ otherwise.  The source's conjunct
\lean{IsElementaryAbelian 2 fourGroupB} is replaced by its content (commutativity and exponent $2$),
because the predicate lives in a source prelude that is not ported.  Also proved: \lean{Hmat_le_Gmat}
($G_{\mathrm{vis}}\le G_{\mathrm{mat}}$), \lean{tauB_mem_Gmat}, \lean{xmat_eq_xB}, and the port
addition \lean{Gmat_eq_closure_insert : Gmat = Subgroup.closure (insert tauB (Gvis : Set _))}.
\end{dictionary}
\begin{physics}
$x_B$ is the orbifold involution $z$; $y_B$ and $z_B$ are its triality images, which act as $-1$ on
the odd-momentum exponentials and on half of the twisted states.  In FLM's language $\{1,x_B,y_B,z_B\}$
is the four-group whose nontrivial elements are permuted by triality; the three $2^{1+24}$ subgroups
they centralise are conjugate in the Monster.
\end{physics}
\begin{status}
No hypotheses; audited in \leanpath{TrialityAudit.lean}.
\end{status}

\subsection{Product invariance of \texorpdfstring{$\tau_B$}{tauB}}

\begin{theorem}[$\tau_B$ is an automorphism; $G_{\mathrm{mat}}\le$ \lean{AutO}]\label{ch21:thm:tauBgmul}
$\tau_B(u\cdot v)=\tau_B(u)\cdot\tau_B(v)$ for all $u,v$; hence $\tau_B\in$ \lean{AutO} and
$G_{\mathrm{mat}}\le$ \lean{AutO}.
\end{theorem}
\begin{leancode}
theorem tauB_gmul (u v : SpaceB) : tauB (gmul u v) = gmul (tauB u) (tauB v) := by
\end{leancode}
\begin{leancode}
theorem tauB_mem_AutO : tauB ∈ AutO := ⟨isAlgAut_tauB, bform_tauB⟩
\end{leancode}
\begin{leancode}
theorem Gmat_le_AutO : Gmat ≤ AutO := Gmat_le_AutO_of_isAlgAut isAlgAut_tauB
\end{leancode}
(\leanpath{VOA/\allowbreak Moonshine/\allowbreak MonsterGriess/\allowbreak ProductInvariance.lean}; \lean{Gmat_le_AutO} in
\lean{.Generated}.)
\begin{dictionary}
The three statements say exactly what they appear to say; \lean{Gmat_le_AutO} combines
\lean{tauB_mem_AutO} with \lean{Gvis_le_AutO} (\lean{Gmat_le_AutO_iff : Gmat ≤ AutO ↔ tauB ∈ AutO}).
In the source, \lean{tauB_mem_AutO} sits in \lean{Griess/AutFinite.lean} and \lean{Gmat_le_AutO} in
\lean{Generated/Centralizer.lean}; the port proves the same statements here.
\end{dictionary}
\begin{proof}[Proof on paper]
Since $\tau_B$ preserves $b$, it suffices that the trilinear form is invariant,
$\chXXIT(\tau_Bu,\tau_Bv,\tau_Bw)=\chXXIT(u,v,w)$, and by trilinearity it suffices to check this on triples of
\lean{tbasis} vectors.  There are $196884^3$ such triples, so the proof is organised:
\begin{enumerate}[label=(\roman*)]
\item \emph{Diagonal units.}  Each basis vector is an eigenvector of multiplication by every diagonal
  $E_{kk}$, and $\tau_B$ permutes basis vectors with equal eigenvalues; so any triple containing a
  diagonal unit is invariant (\leanpath{InvarianceDiagonal.lean}).
\item \emph{Grading.}  $\chXXIB$ is graded by $\Gamma=(\Z/2)^2\times\chXXIC/\langle\Omega\rangle$; every
  basis vector is homogeneous, $\chXXIT$ vanishes on triples whose degrees do not sum to $0$, and $\tau_B$
  maps degree $g$ to degree $\alpha(g)$ for an automorphism $\alpha$ of $\Gamma$
  (\leanpath{GriessGrading.lean}).  So only degree-zero triples remain.
\item \emph{Sectors.}  The degree-zero triples fall into six types, each with a closed-form value of
  $\chXXIT$ shown to be invariant under the index shift induced by $\tau_B$: pair${}^3$
  (\leanpath{InvarianceTwoCoord.lean}); odd${}\cdot{}$odd${}\cdot{}$pair
  (\leanpath{InvarianceOddTwoCoord.lean}); the odd triple, one vector from each slot $0,1,2$
  (\leanpath{InvarianceXYZ.lean}); octad${}\cdot{}$octad${}\cdot{}$pair
  (\leanpath{InvarianceOctadTwoCoord.lean}); octad${}\cdot{}$octad${}\cdot{}$octad; and
  odd${}\cdot{}$odd${}\cdot{}$octad, by character sums of $(-1)^{b_d(\delta,\cdot)}$
  (\leanpath{InvarianceOddValues.lean}, \leanpath{InvarianceOddOctad.lean}).  Shared lemmas are in
  \leanpath{InvarianceOctadBasic.lean}.  In each case $\tau_B$ acts on the relevant indices by the shift
  $k\mapsto k+1$ or by $H_dD_d$, and the closed form is checked to be invariant.
\item \emph{The octad${}^3$ case in more detail.}  For three octads $d_1,d_2,d_3$, $\chXXIT$ vanishes on
  their octad vectors unless $d_1+d_2+d_3=0$ in $\chXXIC$.  In that case consider the $18$-dimensional
  quadratic space $E=E_{d_1}\oplus E_{d_2}\oplus E_{d_3}$ with $q=q_{d_1}+q_{d_2}+q_{d_3}$.  A
  \emph{Lagrangian} is a subspace $K$ with $K^\perp=K$ for the polar form (so $\dim K=9$) on which $q$
  vanishes.  The relevant $K$ is the kernel of the map $(\delta_1,\delta_2,\delta_3)\mapsto\delta_1+\delta_2+\delta_3$
  to the cocode, and there is an isotropic vector $a_0$ (built from $d_1\cap d_2$) such that the values
  $\chXXIT(v_{d_1,\delta_1},v_{d_2,\delta_2},v_{d_3,\delta_3})$ form, up to normalisation, a function on $E$
  supported on the coset $a_0+K$ with signs given by $q$ (the ``Lagrangian function'';
  \lean{three_octad_lagrangian_correspondence}, \leanpath{OctadThreeOctad.lean}).  The tensor product of
  the three Fourier matrices $H_{d_i}$ maps the indicator of a Lagrangian coset to itself (the Fourier
  transform of the indicator of $K$ is a multiple of the indicator of $K^\perp=K$), and the tensor product
  of the $D_{d_i}$ multiplies it by $(-1)^q$, which is constant on the coset.  So both, hence $\tau_B$,
  fix $\chXXIT$ on octad${}^3$ (\lean{fourier_lagrangianFun}, \lean{signOp_lagrangianFun},
  \leanpath{InvarianceThreeOctad.lean}).  The precise signs are in the Lean files; this paragraph gives
  the mechanism.
\end{enumerate}
Then $b(\tau_B(u\cdot v),w)=b(u\cdot v,\tau_B^{-1}w)=\chXXIT(u,v,\tau_B^{-1}w)=\chXXIT(\tau_Bu,\tau_Bv,w)=b(\tau_Bu\cdot\tau_Bv,w)$
for all $w$, and $b$ is nondegenerate.
\end{proof}
\begin{physics}
That a hand-built $\Z/3$ rotation mixing the twisted and untwisted sectors preserves all OPE
coefficients is the finite-dimensional form of the statement that the orbifold theory has an extra
symmetry not visible in its construction.  The Fourier-transform step on octads is the algebraic
analogue of a modular $S$-transformation exchanging sectors.
\end{physics}
\begin{status}
No hypotheses; audited in \leanpath{ProductInvarianceAudit.lean} and
\leanpath{MoonshineCNMonsterFinalAudit.lean}.  The heartbeat budget is raised in
\lean{isGriessForm_thetaT} (\leanpath{VisibleAction.lean}) only; no \lean{native_decide} anywhere in
the port.
\end{status}

%======================================================================
\section{The parity-centralizer theorem: \texorpdfstring{$\mathrm{AutO}\cap C(x_B)=G_{\mathrm{vis}}$}{AutO meet C(xB) = Gvis}}\label{ch21:sec:parity}

\begin{theorem}[AxiomMath's rational parity-centralizer theorem]\label{ch21:thm:parity}
(1) A $\Q$-linear automorphism $f$ of $\chXXIB$ preserves the product and the form and commutes with
$x_B$ if and only if $f\in G_{\mathrm{vis}}$.  (2) Every such $f$ is $\theta_p$ for a visible pair $p$,
unique up to $p\mapsto p\cdot(-1,-1)$.  (3) $p\cdot(-1,-1)=(-s,-g)$.  Equivalently,
\lean{AutO} $\cap\,C(x_B)=G_{\mathrm{vis}}$.
\end{theorem}
\begin{leancode}
theorem rational_parity_centralizer_equals_visible :
    (∀ f : SpaceB ≃ₗ[ℚ] SpaceB, ((∀ u v, f (gmul u v) = gmul (f u) (f v)) ∧
        (∀ u v, bform (f u) (f v) = bform u v) ∧ ∀ u, f (xB u) = xB (f u)) ↔ f ∈ Gvis) ∧
    (∀ f : SpaceB ≃ₗ[ℚ] SpaceB, (∀ u v, f (gmul u v) = gmul (f u) (f v)) →
      (∀ u v, bform (f u) (f v) = bform u v) → (∀ u, f (xB u) = xB (f u)) →
      ∃ p : Pvis, thetaEquiv p = f ∧
        ∀ p' : Pvis, thetaEquiv p' = f ↔ p' = p ∨ p' = p * negPair) ∧
    ∀ p : Pvis, (((p * negPair : Pvis).1.1 : GL SpinIdx ℚ) : Matrix SpinIdx SpinIdx ℚ) =
        -((p.1.1 : GL SpinIdx ℚ) : Matrix SpinIdx SpinIdx ℚ) ∧
      mat (p * negPair : Pvis).1.2 = -mat p.1.2 := by
\end{leancode}
\begin{leancode}
theorem AutO_inf_centralizer_xB_eq_Gvis :
    AutO ⊓ Subgroup.centralizer {xB} = Gvis := by
\end{leancode}
(\leanpath{VOA/\allowbreak Moonshine/\allowbreak MonsterGriess/\allowbreak ParityEqualsVisible.lean}, namespace \lean{.Parity}.)
\begin{dictionary}
\lean{⊓} is intersection of subgroups; \lean{Subgroup.centralizer {xB}} is $\{f: fx_B=x_Bf\}$.  The
corollary is a port addition (the subgroup form of part (1)).  Part (1) has the three conditions on the
left exactly as in \lean{AutO} plus commutation with $x_B$; the equivalence is an ``iff'', so it also
contains $G_{\mathrm{vis}}\le$ \lean{AutO}.  This theorem does \emph{not} depend on
$G_{\mathrm{mat}}=$ \lean{AutO}: its AxiomMath import closure is $82$ files, and the port contains all of
it that the proof uses.
\end{dictionary}
\begin{proof}[Proof on paper]
Let $f$ be a ``parity automorphism'' (\lean{IsParityAut}: product, form, commutes with $x_B$).
\begin{enumerate}[label=(\alph*)]
\item \emph{Blocks.}  $f$ preserves $W$ (the $-1$-eigenspace of $x_B$) and $\chXXISym\oplus X$.  For
  $u=(A,x,0)$ let $L_u:W\to W$ be multiplication by $u$; by the table,
  $L_u(N)=\tfrac{\tr A}8N+NA+\tfrac18\sum_cx_cR_cN(1-\tfrac14U_c)$.  Since $W=\Q^I\otimes\Q^{24}$ and $L_u$ is
  a sum of operators $M\otimes M'$ (left factor on $\Q^I$, right factor on $\Q^{24}$), traces factor:
  $\tr(M\otimes M')=\tr M\cdot\tr M'$.  For $u=(A,0,0)$, $v=(A',0,0)$: $L_u=1\otimes(\tfrac{\tr A}8+A)$, so
  \begin{equation*}
    \tr(L_uL_v)=4096\bigl(24\tfrac{\tr A\tr A'}{64}+2\tfrac{\tr A\tr A'}8+\tr(AA')\bigr)
    =4096\tr(AA')+512\cdot5\tr A\tr A',
  \end{equation*}
  since $4096(\tfrac{24}{64}+\tfrac28)=4096\cdot\tfrac58=2560$.  For $u=e_c$, $v=e_d$:
  $L_{e_c}=\tfrac18R_c\otimes(1-\tfrac14U_c)$; $\tr(R_cR_d)=0$ for $c\neq d$ (a non-scalar element of $Q$
  has trace $0$) and $R_c^2=1$ for short $c$, while
  $\tr\bigl((1-\tfrac14U_c)^2\bigr)=\tr(1+\tfrac32U_c)=24+48=72$ (using $U_c^2=32U_c$, $\tr U_c=32$); so
  $\tr(L_{e_c}L_{e_d})=\tfrac1{64}\cdot4096\cdot72\,\delta_{cd}=4608\,\delta_{cd}=(4096+512)\delta_{cd}$.  Mixed
  terms vanish because $\tr R_c=0$.  Altogether
  $\tr(L_uL_v)=4096\,b(u,v)+512\,\sigma(u,v)$ with $\sigma(u,v)=5\tr A\tr A'+x\cdot x'$.  Since $f$ conjugates
  $L_u$ to $L_{fu}$, it preserves $\sigma$ as well as $b$.  Let $\Phi$ be the operator with
  $b(\Phi u,v)=\sigma(u,v)$: $\Phi(A,x)=(5\tr(A)I,x)$.  It commutes with $f$, and its eigenspaces are
  $\Q I$ (eigenvalue $120$), the traceless matrices (eigenvalue $0$) and $X$ (eigenvalue $1$); so $f$
  preserves $\Q I$, $\chXXISym=\Q I\oplus\chXXISym_0$ and $X$ (\leanpath{ParityTraceSplitting.lean}).
\item \emph{$\chXXISym$-block.}  On $\chXXISym$ the product is the Jordan product $2(AB+BA)$ and $b$ is the trace form;
  a Jordan automorphism of real symmetric matrices preserving the trace form is $A\mapsto g^{\mathsf T}Ag$
  with $g$ orthogonal (\leanpath{ParityRealExtension.lean}; the proof extends scalars to $\mathbb{R}$).
\item \emph{Short lines and $\mathrm{Co}_0$.}  $f$ maps each $e_c$ to $\pm e_{\ell(c)}$ for a bijection
  $\ell$ with $q_{\ell(c)}(g^{\mathsf T}Ag)=q_c(A)$; hence $g$ maps minimal vectors to minimal vectors up
  to sign, is rational with denominators dividing $8$, and preserves $\Leech$: $g\in\mathrm{Co}_0$
  (\leanpath{ParityShortLines.lean}, \leanpath{ScalarKernel.lean}).
\item \emph{$W$-block.}  On $W$, $f(N)=s^{\mathsf T}Ng$ for an orthogonal $s$, and comparing the $X\cdot W$
  products shows that $s$ conjugates every $R_c$ to $\pm R_{\ell(c)}$; since $Q$ is generated by its short
  elements, $s$ normalises $Q$, hence has dyadic entries (\lean{dyadic_of_conj_Q},
  \leanpath{SpinExtraspecial.lean}), and $(s,g)$ is a visible pair with $\theta_{(s,g)}=f$.
\end{enumerate}
The converse is Theorem~\ref{ch21:thm:gvis}(a) together with the observation that each $\theta_p$
commutes with $x_B$ ($N\mapsto s^{\mathsf T}Ng$ is linear in $N$).
\end{proof}
\begin{physics}
``Every symmetry of the Monster CFT that commutes with the orbifold involution is a manifest
symmetry of the orbifold construction.''  Physically this is plausible: a symmetry commuting with $z$
preserves the untwisted sector, hence acts on the Leech lattice theory's $\Z/2$-invariant part, and
such symmetries are rotations of the lattice combined with momentum shifts.  The theorem makes this
exact at the level of the Griess algebra.
\end{physics}
\begin{status}
No hypotheses.  Both audited in \leanpath{ParityAudit.lean} (795 \texttt{\#print axioms} lines over the
parity closure), standard axioms.  Two proof steps are re-proved in the port with the same statements
(\lean{phi0_coordOf_inv}, \lean{phi0_coordOf_mul}) because of Mathlib-version instance problems; the
port notes list them.
\end{status}

%======================================================================
\section{What is not ported: the generation theorem}\label{ch21:sec:notported}

AxiomMath's library contains, in Lean~v4.35:
\begin{plaincode}
theorem gmat_eq_griess :
    CFSG.Monster.Generated.Gmat = CFSG.Monster.Griess.AutO :=
  CFSG.YGroups.Y443.Gmat_eq_AutO_of_stroth
\end{plaincode}
This block is \emph{not} LeanVOA source, which is why it is set as plain text rather than as a
checked Lean quotation; it is quoted from \leanpath{Notes/\allowbreak monster\_aut/\allowbreak CFSG\_PORT\_FIDELITY.md} (source
file \lean{CFSG/Sporadic/Bridge/MonsterModelsGriess.lean}, namespace
\lean{CFSG.Blueprint.NativeSporadic.MonsterModel}).  Both sides are written with fully qualified names,
and they are the source declarations of which our \lean{Generated.Gmat} and \lean{AutO} are ports.

\begin{warning}[title={Dead end: porting \texttt{gmat\_eq\_griess} to Lean 4.30}]
The authors' reconnaissance measured the file-level import closure
of \lean{gmat_eq_griess}: about $16{,}361$ files, $208$\,MB, $4.19$ million lines, with $1{,}720$
certificate files, $13{,}487$ \lean{decide +kernel} proofs and heartbeat budgets up to $1.6\cdot10^8$.
The proof cannot avoid the Stroth--Leon--Sims recognition of the Baby Monster: the alternative
``two involution classes'' route (\lean{GmatGriess_eq_top_of}) has a hypothesis that is discharged only
through the same recognition.  ``A Stroth-free proof of $G_{\mathrm{mat}}=$ \lean{AutO} inside CFSG'' is
therefore a recorded dead end.  A port of that closure to Lean~4.30, or a move of the whole library to Lean~4.35, was judged far too
costly; the library stays on Lean~4.30 and carries $G_{\mathrm{mat}}=$ \lean{AutO} as one named hypothesis,
\lean{GmatEqAutO} (Chapter~\ref{ch:22}).  The axiom footprint of \lean{gmat_eq_griess} in AxiomMath's
toolchain was not checked.
\end{warning}

\begin{warning}[title={Recorded diagnosis: AxiomMath's finiteness proof does not transfer to $\mathbb C$}]
AxiomMath prove \lean{Finite AutO} analytically over the ordered field $\Q$ ($\ell^1$-bounded and
separated orbits of $2A$-axes, \lean{finite_of_isolated_axes_basis}).  Over $\C$ the form is not a norm
and that argument fails; the same note records that the step ``real and finite
$\Rightarrow$ rational'' for automorphisms of $\C\otimes\chXXIB$ is Tits's theorem, for which no cheap
route was found.  The converse direction actually used in Chapter~\ref{ch:22} avoids Tits by a
different argument (rationality of traces and a commutant lemma).
\end{warning}

So the port provides the \emph{statements} $G_{\mathrm{vis}}\le G_{\mathrm{mat}}\le$ \lean{AutO},
\lean{AutO} $\cap\,C(x_B)=G_{\mathrm{vis}}$, $\tau_B\in$ \lean{AutO}, and the definitions on both sides
of \lean{gmat_eq_griess}, but not \lean{gmat_eq_griess} itself, nor finiteness, order or simplicity of
\lean{AutO}.


%======================================================================
\section{The port-fidelity audit}\label{ch21:sec:fidelityaudit}

All the definitions of Part~A have now been presented, so we can state precisely in what sense they are
AxiomMath's.  Because \lean{GmatEqAutO} (Chapter~\ref{ch:22}) is to be read as ``AxiomMath's theorem
\lean{gmat_eq_griess}'', it matters that our \lean{Generated.Gmat} and \lean{AutO} are the \emph{same
definitions} as theirs.  A read-only audit
(\leanpath{Notes/\allowbreak monster\_aut/\allowbreak CFSG\_PORT\_FIDELITY.md}) compared, declaration by
declaration, the definitional closure of the two constants on both sides.  Its method and verdict, which
this book reports as the authors' claim (a source-text comparison, not a Lean check):
\begin{itemize}
\item Starting from \lean{Gmat} and \lean{AutO}, every identifier occurring in a definition body, and in
  the \emph{statement} of every theorem referenced from a definition body (theorems can fix data: \lean{tauB}
  gets its matrix from the statement of \lean{tauMat_transpose_mul}, \lean{tbasis} its vectors from
  \lean{linearIndependent_tvec}), was followed down to Mathlib.  This gave $199$ definition-like
  declarations and $104$ theorem statements.
\item $198$ definition bodies were matched to the source: $194$ are textually identical after removing
  comments, attributes, visibility keywords and the namespace prefixes; three differ only in proof terms
  of propositions (\lean{bform}, \lean{gmul}) or by a parser artefact (\lean{signMatrix}); one differs in
  substance: \lean{SpinLattice.phi0}.
\item All $88$ data-fixing theorem statements that have a source counterpart are identical; $44$ of the
  $46$ instances are identical to source instances and the other $2$ lie outside the closure; the relevant
  \lean{variable} contexts match line for line.
\end{itemize}

\begin{warning}[title={The one substantive difference: \texttt{phi0} is a function, not a bundled linear map}]
In the source, $\varphi_0$ (Definition~\ref{ch21:def:short}) is the bundled $\chXXIF$-linear map
\lean{(AddMonoidHom.toZModLinearMap 2 lamTHom).coprod lamCocode}; in the port
(\leanpath{VOA/\allowbreak Moonshine/\allowbreak MonsterGriess/\allowbreak Short.lean}) it is the plain
function \lean{fun p => lamT p.1 + lamCocode p.2}, because bundling needs the additivity of
$\lambda_T$, which an earlier version of the port did not have.  The two are \emph{definitionally}
equal (\lean{LinearMap.coprod f g} unfolds pointwise to \lean{f p.1 + g p.2}, and
\lean{lamTHom.toFun := lamT}), so \lean{IsShortCoord}, hence \lean{ShortIdx}, \lean{SpaceX},
\lean{SpaceB}, \lean{gmul}, \lean{bform}, \lean{AutO} and \lean{Gmat} have the same \emph{values}.  But the
elaborated terms are not syntactically identical, so a literal name-for-name transport of AxiomMath's
\lean{gmat_eq_griess} across toolchains would need a bridge step closed by \lean{rfl}.  The port now
proves additivity (\lean{phi0_add}) and bundles the map as \lean{phi0Lin}, with
\lean{phi0Lin_apply : phi0Lin p = phi0 p} by \lean{rfl}.
\end{warning}

\begin{warning}[title={Out-of-date port docstrings}]
Some port docstrings still say that something is ``not ported'' that now is.
\leanpath{Short.lean} says this of \lean{phi0Equiv} and \lean{card_shortIdx}, which are in
\leanpath{LatticeQuadratic.lean} and \leanpath{LatticeFamilies.lean}.  \leanpath{LeechMatch.lean} calls
bijectivity of \lean{phi0} ``the open statement'' (its \lean{shortIdxEquivPm} and
\lean{card_shortIdx_of_bijective} take \lean{Function.Bijective phi0} as a hypothesis); that hypothesis
is the theorem \lean{SpinLattice.phi0_bijective}.  The docstring of \lean{ShortSignDictionary} in
\leanpath{MoonshineCNGriessUntwisted.lean} says \lean{phi0Equiv} is not ported.
\leanpath{TrialityFourGroup.lean} says that \lean{tauB ∈ AutO} is ``not proved here'', true of that file;
it is proved in \leanpath{ProductInvariance.lean}.  The Lean statements are unaffected.
\end{warning}

%======================================================================
\section{The Griess algebra of \texorpdfstring{$V^\natural$}{V-natural}: the untwisted half}\label{ch21:sec:vnatside}

Part~B begins.  Its Lean names live in the namespace \lean{VOA.Moonshine.GriessDictionary}.  We use
the moonshine module \lean{VnatVOA} of Chapter~\ref{ch:20} (\leanpath{VOA/\allowbreak Moonshine/\allowbreak AutVnat.lean}),
which is by definition the vertex operator algebra of Chapter~\ref{ch:14} on the carrier
\lean{Vnat := SkewVP × TP} $=V_\Leech^+\oplus(V_\Leech^T)^+$, with \lean{VnatVOA.Y = Ynat wbarL},
\lean{VnatVOA.grading = gradingNat} (\lean{rfl}).  Its Griess product is
\begin{leancode}
def griessMul (W : VertexOperatorAlgebra 𝒱) (u v : W.grading 2) : W.grading 2 :=
  ⟨W.smul (u : 𝒱) 1 (v : 𝒱), by
\end{leancode}
(\leanpath{VOA/\allowbreak Core/\allowbreak LieAlgebraV1.lean}), i.e.\ $u\times v=u_{(1)}v$ on $W_2$, and its invariant form is
\lean{SimpleNat.natForm} (Chapter~\ref{ch:10}): \lean{natForm x y = leechBV x.1 y.1 + TPForm x.2 y.2},
the sum of the invariant form $B_V$ of $V_\Leech^+$ (= \lean{plusForm}, Chapter~\ref{ch:04}) and the
form \lean{TPForm} of $(V_\Leech^T)^+$ (Chapter~\ref{ch:09}), normalised by $\langle\mathbb1,\mathbb1\rangle=1$.
On $\Vnat_2$ it coincides with the intrinsic form $u_{(3)}v=\langle u,v\rangle\mathbb1$
(Theorem~\ref{ch21:thm:griessFormNat}).

\subsection{The \texorpdfstring{$S^2\mathfrak h$}{S2h} block in Conway's coordinates}\label{ch21:sec:conway}

The library writes oscillator states in an orthonormal frame $h_a$ of $\mathfrak h=\C\otimes\Leech$
(\lean{chargeCoord}, Chapter~\ref{ch:03}), whose entries are irrational.  To compare with AxiomMath's
rational matrices, the $S^2\mathfrak h$ block is rewritten in Conway's integer coordinates, using the
lattice vectors $8e_i\in\Leech$ (norm $8$, $\langle8e_i,\beta\rangle=\beta_i$).

\begin{definition}[Conway axes and quadratic states]\label{ch21:def:conwayQuad}
$\mathrm{ax}_i=8e_i\in\Leech$.  For a complex $24\times24$ matrix $A$,
$Q_C(A)=\tfrac18\sum_{i,j}A_{ij}\,\mathrm{ax}_i(-1)\mathrm{ax}_j(-1)\mathbb1\in(V_\Leech^+)_2$.
\end{definition}
\begin{leancode}
def conwayAxis (i : Fin 24) : leech.carrier :=
  (⟨(8 : ℤ) • Pi.single i 1, eight_smul_mem_leechCarrier _⟩ : ↥Leech.carrier)
\end{leancode}
\begin{leancode}
noncomputable def conwayQuad : Matrix (Fin 24) (Fin 24) ℂ →ₗ[ℂ] VPlus := frameQuad ∘ₗ conwayToFrame
\end{leancode}
\begin{leancode}
theorem conwayQuad_eq_sum_dV (A : Matrix (Fin 24) (Fin 24) ℂ) :
    conwayQuad A = (8 : ℂ)⁻¹ • ∑ i, ∑ j, A i j • (dV (conwayAxis i) (conwayAxis j) : VPlus) := by
\end{leancode}
(\leanpath{VOA/\allowbreak Moonshine/\allowbreak MoonshineCNGriessConway.lean}.)
\begin{dictionary}
\lean{frameQuad M} $=\sum_{a,b}M_{ab}h_a(-1)h_b(-1)\mathbb1$ in the orthonormal frame;
\lean{conwayToFrame A} $=\tfrac18FAF^{\mathsf T}$ with $F$ the frame coordinates of the axes
(\lean{conwayFrame}, $F^{\mathsf T}F=8\cdot1$).  \lean{dV α β} is the state $\alpha(-1)\beta(-1)\mathbb1$.
So the definition through the frame and the displayed formula agree (\lean{conwayQuad_eq_sum_dV}).
\lean{VPlus} is $V_\Leech^+$.
\end{dictionary}
\begin{physics}
$Q_C(A)$ is the composite operator $\tfrac18A_{ij}\partial X^i\partial X^j$ written in Conway's
lattice coordinates; $Q_C(\tfrac12I)=\tfrac12\sum_a h_a(-1)^2\mathbb1=\omega$ is the stress tensor
(since $\mathrm{ax}_i/\sqrt8$ is a unit vector).
\end{physics}
\begin{status}
Definition; audited in \leanpath{MoonshineCNGriessAudit.lean}.
\end{status}

\begin{theorem}[The untwisted Griess products in Conway coordinates]\label{ch21:thm:conway}
For symmetric $A,B$ and $\beta\in\Leech$ with integer coordinates $y_\beta$:
(a) $Q_C(A)_{(1)}Q_C(B)=2\,Q_C(AB+BA)$;
(b) $Q_C(A)_{(1)}x_\beta=\tfrac18(y_\beta^{\mathsf T}Ay_\beta)\,x_\beta$, where $x_\beta=e^\beta+e^{-\beta}$;
(c) $\beta(-1)^2\mathbb1=Q_C(\tfrac18y_\beta y_\beta^{\mathsf T})$;
(d) $\langle Q_C(A),Q_C(B)\rangle=2\tr(AB)$; $\langle x_\alpha,x_\beta\rangle=2\varepsilon(\alpha,\alpha)$ if
$\beta=\pm\alpha$ and $0$ otherwise ($\alpha$ minimal).
\end{theorem}
\begin{leancode}
theorem conwayQuad_one_conwayQuad_of_symm {A B : Matrix (Fin 24) (Fin 24) ℂ} (hA : Aᵀ = A)
    (hB : Bᵀ = B) :
    leechVOAPlus.smul (conwayQuad A) 1 (conwayQuad B) = (2 : ℂ) • conwayQuad (A * B + B * A) := by
\end{leancode}
\begin{leancode}
theorem conwayQuad_one_frameX0P (A : Matrix (Fin 24) (Fin 24) ℂ) (β : leech.carrier) :
    leechVOAPlus.smul (conwayQuad A) 1 (frameX0P β)
      = ((8 : ℂ)⁻¹ * ((fun i => (vec β i : ℂ)) ⬝ᵥ (A *ᵥ fun i => (vec β i : ℂ)))) • frameX0P β := by
\end{leancode}
\begin{leancode}
theorem dV_self_eq_conwayQuad (β : leech.carrier) :
    (dV β β : VPlus)
      = conwayQuad ((8 : ℂ)⁻¹ • vecMulVec (fun i => (vec β i : ℂ)) (fun i => (vec β i : ℂ))) := by
\end{leancode}
\begin{leancode}
theorem plusForm_conwayQuad_conwayQuad_of_symm {A B : Matrix (Fin 24) (Fin 24) ℂ}
    (hB : Bᵀ = B) : plusForm (conwayQuad A) (conwayQuad B) = 2 * trace (A * B) := by
\end{leancode}
\begin{leancode}
theorem plusForm_frameX0P_frameX0P {α : leech.carrier} (hα : leech.form α α = 4)
    (β : leech.carrier) :
    plusForm (frameX0P α) (frameX0P β)
      = if β = α ∨ β = -α then 2 * latticeCocycle leech α α else 0 := by
\end{leancode}
(\leanpath{VOA/\allowbreak Moonshine/\allowbreak MoonshineCNGriessConway.lean}.)
\begin{dictionary}
\lean{leechVOAPlus.smul u 1 v} is $u_{(1)}v$ in $V_\Leech^+$; \lean{frameX0P β} is $x_\beta$
(Chapter~\ref{ch:04}); \lean{vec β} the integer coordinates of $\beta$; \lean{latticeCocycle leech α α}
is $\varepsilon(\alpha,\alpha)\in\{\pm1\}$ (Chapter~\ref{ch:06}; it need not be $1$).  The products of the
$x_\alpha$ among themselves are lemmas of \leanpath{MoonshineLRFinOrderIsing.lean} (Chapter~\ref{ch:20}),
used unchanged.  Compare with the table of \S\ref{ch21:sec:spaceB}: (a) is $A\cdot A'=2(AA'+A'A)$; (b)
is $A\cdot x=\tfrac18q_c(A)x_c$; (c) is the source of $\tfrac18U_c$ in $x\cdot y$; (d) is $b=\tr(AA')$
up to the factor $2$.
\end{dictionary}
\begin{proof}[Proof on paper]
(a) Wick's theorem for the Heisenberg VOA: $(h_ah_b)_{(1)}(h_ch_d)$ is the sum over single contractions,
giving $\sum M_{ab}N_{cd}(\delta_{bc}h_ah_d+\dots)$, i.e.\ $Q((M+M^{\mathsf T})(N+N^{\mathsf T}))$; then
transport through $\tfrac18F(\cdot)F^{\mathsf T}$ using $F^{\mathsf T}F=8$.  (b) $h_a(1)h_b(1)e^\beta$ produces
$\langle h_a,\beta\rangle\langle h_b,\beta\rangle e^\beta$.  (c), (d) are similar contractions.
\end{proof}
\begin{physics}
(a) is the OPE of two bilinears in free bosons; (b) says $\cos(\beta X)$ is an eigen-operator of the
``metric deformation'' $A_{ij}\partial X^i\partial X^j$ with eigenvalue proportional to
$\beta^{\mathsf T}A\beta$ (the change of its conformal weight); (d) are free-boson two-point functions.
\end{physics}
\begin{status}
No hypotheses beyond those displayed; all audited in \leanpath{MoonshineCNGriessAudit.lean}.
\end{status}

\subsection{The untwisted dictionary}

\begin{definition}[Matrix and short dictionaries]\label{ch21:def:dict}
For a rational matrix $M$ in AxiomMath's coordinates let $M^\pi=(M_{\pi i,\pi j})_{ij}$ (a complex
matrix in our coordinates) and $\mathrm{matDict}(M)=Q_C(M^\pi)$.  For $c\in\chXXIS$ let $\alpha_c$ be the
minimal vector of our lattice that $\pi$ carries to $u_c$.  For a phase function $\zeta:\chXXIS\to\C$,
$\mathrm{shortDict}_\zeta(x)=\sum_cx_c\zeta_c\,x_{\alpha_c}$, and
$\Phi_\zeta(A,x,N)=\mathrm{matDict}(A)+\mathrm{shortDict}_\zeta(x)$.
\end{definition}
\begin{leancode}
noncomputable def axiomToConway (M : Matrix (Fin 24) (Fin 24) ℚ) : Matrix (Fin 24) (Fin 24) ℂ :=
  Matrix.of fun i j => (M (piPerm i) (piPerm j) : ℂ)
\end{leancode}
\begin{leancode}
noncomputable def matDict : Matrix (Fin 24) (Fin 24) ℚ →+ VPlus where
  toFun M := conwayQuad (axiomToConway M)
\end{leancode}
\begin{leancode}
noncomputable def shortVec (c : ShortIdx) : leech.carrier :=
  leechEquiv.symm ⟨uvec c, (uvec_mem_minVecs c).1⟩
\end{leancode}
\begin{leancode}
noncomputable def shortDict (ζ : ShortIdx → ℂ) : SpaceX →+ VPlus where
  toFun x := ∑ c, ((x c : ℂ) * ζ c) • frameX0P (shortVec c)
\end{leancode}
\begin{leancode}
noncomputable def untwistedDict : SpaceB →+ VPlus where
  toFun u := matDict u.1 + shortDict ζ u.2.1
\end{leancode}
(\leanpath{VOA/\allowbreak Moonshine/\allowbreak MoonshineCNGriessUntwisted.lean}.)
\begin{dictionary}
The maps are additive (\lean{→+}) and $\Q$-homogeneous (\lean{matDict_smul}, \lean{untwistedDict_smul});
they are not declared $\Q$-linear because $\Q$-linearity into a $\C$-space is automatic for additive maps.
\lean{untwistedDict} ignores the $W$-component.  \lean{shortVec_norm} says
$\langle\alpha_c,\alpha_c\rangle=4$ and \lean{quadForm_shortVec} that
$y_{\alpha_c}^{\mathsf T}A^\pi y_{\alpha_c}=q_c(A)$, which is how (b) of
Theorem~\ref{ch21:thm:conway} becomes AxiomMath's $\tfrac18q_c(A)$.
\end{dictionary}
\begin{physics}
The dictionary writes each element of AxiomMath's $\chXXISym\oplus X$ as a concrete spin-two field of the
untwisted sector: matrices become boson bilinears, the basis vector $e_c$ becomes
$\zeta_c(e^{i\alpha_cX}+e^{-i\alpha_cX})$.  The phase $\zeta_c$ is a choice of normalisation of the
vertex operator, which is ambiguous up to the cocycle.
\end{physics}
\begin{status}
Definitions; audited in \leanpath{MoonshineCNGriessAudit.lean}.
\end{status}

\begin{definition}[Sign dictionary]\label{ch21:def:ssd}
A phase function $\zeta:\chXXIS\to\C$ is a \emph{sign dictionary} if the following six conditions hold.
The first three are about supports and do not involve $\zeta$; they follow from additivity and
injectivity of $\varphi_0$ (Theorem~\ref{ch21:thm:ofphi0}), so in practice only the last three matter.
\begin{itemize}
\item (pm\_injective) distinct short classes have distinct pairs $\pm\alpha_c$: $\alpha_d=\pm\alpha_c\Rightarrow d=c$;
\item (sum\_shape) if $c=d+e$ in $\chXXIC\times\chXXIC^*$, then either $\langle\alpha_d,\alpha_e\rangle=-2$ and
  $\alpha_c=\pm(\alpha_d+\alpha_e)$, or $\langle\alpha_d,\alpha_e\rangle=2$ and $\alpha_c=\pm(\alpha_e-\alpha_d)$;
\item (sum\_exists) if $\langle\alpha_d,\alpha_e\rangle=\pm2$ then $d+e$ is short;
\item (sq) $\zeta_c^2\varepsilon(\alpha_c,\alpha_c)=1$ (normalisation of the form);
\item (sign\_neg\_two), (sign\_two) for $c=d+e$ with $\langle\alpha_d,\alpha_e\rangle=\mp2$:
  $(-1)^{\langle a_e,\delta_d\rangle}\zeta_c=\zeta_d\zeta_e\varepsilon(\alpha_d,\pm\alpha_e)$ (AxiomMath's sign
  against our cocycle).
\end{itemize}
\end{definition}
\begin{leancode}
structure ShortSignDictionary (ζ : ShortIdx → ℂ) : Prop where
  sq : ∀ c, ζ c ^ 2 * latticeCocycle leech (shortVec c) (shortVec c) = 1
  pm_injective : ∀ c d : ShortIdx, (shortVec d = shortVec c ∨ shortVec d = -shortVec c) → d = c
  sum_shape : ∀ d e c : ShortIdx, c.1 = d.1 + e.1 →
    (leech.form (shortVec d) (shortVec e) = -2 ∧
      (shortVec c = shortVec d + shortVec e ∨ shortVec c = -(shortVec d + shortVec e))) ∨
    (leech.form (shortVec d) (shortVec e) = 2 ∧
      (shortVec c = shortVec e - shortVec d ∨ shortVec c = -(shortVec e - shortVec d)))
  sum_exists : ∀ d e : ShortIdx, (leech.form (shortVec d) (shortVec e) = -2 ∨
    leech.form (shortVec d) (shortVec e) = 2) → ∃ c : ShortIdx, c.1 = d.1 + e.1
  sign_neg_two : ∀ d e c : ShortIdx, c.1 = d.1 + e.1 → leech.form (shortVec d) (shortVec e) = -2 →
    (MonsterGriess.sgn (pair e.1.1 d.1.2) : ℂ) * ζ c = ζ d * ζ e * latticeCocycle leech (shortVec d) (shortVec e)
  sign_two : ∀ d e c : ShortIdx, c.1 = d.1 + e.1 → leech.form (shortVec d) (shortVec e) = 2 →
    (MonsterGriess.sgn (pair e.1.1 d.1.2) : ℂ) * ζ c = ζ d * ζ e * latticeCocycle leech (shortVec d) (-shortVec e)
\end{leancode}
(\leanpath{VOA/\allowbreak Moonshine/\allowbreak MoonshineCNGriessUntwisted.lean}.)
\begin{dictionary}
A \lean{structure … : Prop} is a conjunction with named fields.  Every field is a statement about
lattice vectors, AxiomMath's coordinates and the cocycle; no vertex algebra enters.  The structure is
used as a \emph{hypothesis} in \lean{untwistedDict_gmul} and \lean{untwistedDict_form}, and then
\emph{discharged} for an explicit $\zeta$ (Theorem~\ref{ch21:thm:griessShortSign}), so it is not an open
input of any final statement.  Its docstring's remark that \lean{phi0Equiv} ``is not ported'' is stale.
\end{dictionary}
\begin{physics}
These conditions say that, after normalising the vertex operators by $\zeta$, the OPE coefficients
$\varepsilon(\alpha,\beta)$ of the lattice CFT become AxiomMath's signs $(-1)^{\langle a_e,\delta_d\rangle}$.
Two normalisations of vertex operators related by a coboundary give isomorphic OPEs; this is that
statement made explicit.
\end{physics}
\begin{status}
Definition; audited in \leanpath{MoonshineCNGriessAudit.lean}.
\end{status}

\begin{theorem}[The untwisted dictionary is an isometric algebra map]\label{ch21:thm:untwisted}
If $\zeta$ is a sign dictionary, then for all $u,v\in\chXXIB$ with zero $W$-part,
$\Phi_\zeta(u\cdot v)=\Phi_\zeta(u)_{(1)}\Phi_\zeta(v)$ and $\langle\Phi_\zeta u,\Phi_\zeta v\rangle=2\,b(u,v)$.
\end{theorem}
\begin{leancode}
theorem untwistedDict_gmul (hζ : ShortSignDictionary ζ) {u v : SpaceB} (hu : u.2.2 = 0)
    (hv : v.2.2 = 0) :
    untwistedDict ζ (gmul u v) = leechVOAPlus.smul (untwistedDict ζ u) 1 (untwistedDict ζ v) := by
\end{leancode}
\begin{leancode}
theorem untwistedDict_form (hζ : ShortSignDictionary ζ) {u v : SpaceB} (hu : u.2.2 = 0)
    (hv : v.2.2 = 0) :
    plusForm (untwistedDict ζ u) (untwistedDict ζ v) = 2 * (bform u v : ℂ) := by
\end{leancode}
(\leanpath{VOA/\allowbreak Moonshine/\allowbreak MoonshineCNGriessUntwisted.lean}; the same statements in $\Vnat$ are
\lean{untwistedDictNat_gmul}, \lean{untwistedDictNat_form}.)
\begin{dictionary}
\lean{ζ} is a section variable \lean{{ζ : ShortIdx → ℂ}}.  The
hypotheses: the sign dictionary and the vanishing of the $W$-components.  The products are in
$V_\Leech^+$, which is a subalgebra of $\Vnat$ (\lean{vnat_smul_untwisted}), and the form \lean{plusForm}
is the restriction of \lean{natForm} (\lean{natForm_untwisted}).
\end{dictionary}
\begin{proof}[Proof on paper]
Bilinearity reduces to the channels $\chXXISym\cdot\chXXISym$, $\chXXISym\cdot X$, $X\cdot X$ and to the forms.  The first two and
the $S$-forms are Theorem~\ref{ch21:thm:conway} (a), (b), (d) transported by $\pi$.  For $X\cdot X$, take
basis vectors $e_d,e_e$.  In $V_\Leech^+$, $x_{\alpha}{}_{(1)}x_\beta$ (both minimal) is: a multiple of
$\alpha(-1)^2\mathbb1$ if $\beta=\pm\alpha$ (the $\alpha(-2)\mathbb1$ terms of
$e^\alpha_{(1)}e^{-\alpha}$ and $e^{-\alpha}_{(1)}e^{\alpha}$ cancel); a cocycle multiple of
$x_{\alpha+\beta}$ if $\langle\alpha,\beta\rangle=-2$; of $x_{\alpha-\beta}$ if $\langle\alpha,\beta\rangle=2$;
and $0$ otherwise (the vertex operator $e^\alpha_{(1)}e^\beta$ has leading power
$z^{\langle\alpha,\beta\rangle}$).  On AxiomMath's side $e_d\cdot e_e$ has $S$-part $\tfrac18\delta_{de}U_d$ and
$X$-part $\tau(d,e,\cdot)$, nonzero only at $c=d+e$.  The fields \lean{pm_injective}, \lean{sum_shape},
\lean{sum_exists} match the supports, and \lean{sq}, \lean{sign_neg_two}, \lean{sign_two} match the
coefficients.
\end{proof}
\begin{physics}
The OPE of untwisted spin-two fields of the orbifold equals AxiomMath's formulas for $\chXXISym\oplus X$ once
the vertex operators are normalised by $\zeta$.
\end{physics}
\begin{status}
Conditional on \lean{ShortSignDictionary ζ} (discharged below); audited in
\leanpath{MoonshineCNGriessAudit.lean}.
\end{status}

\begin{theorem}[Three fields from $\varphi_0$]\label{ch21:thm:ofphi0}
Additivity and injectivity of $\varphi_0$ imply \lean{pm_injective}, \lean{sum_shape}, \lean{sum_exists};
so a sign dictionary is the same as a phase function satisfying \lean{sq}, \lean{sign_neg_two},
\lean{sign_two}.  The lattice input: if three minimal vectors satisfy $\gamma\equiv\alpha+\beta$ modulo
$2\Leech$, then $\alpha\cdot\beta=\mp16$ (i.e.\ $\langle\alpha,\beta\rangle=\mp2$) and
$\gamma=\pm(\alpha\pm\beta)$ correspondingly.
\end{theorem}
\begin{leancode}
theorem minVecs_class_add_cases {a b c : Fin 24 → ℤ} (ha : a ∈ minVecs) (hb : b ∈ minVecs)
    (hc : c ∈ minVecs) (h : classOf c = classOf a + classOf b) :
    (a ⬝ᵥ b = -16 ∧ (c = a + b ∨ c = -(a + b))) ∨ (a ⬝ᵥ b = 16 ∧ (c = b - a ∨ c = -(b - a))) := by
\end{leancode}
\begin{leancode}
theorem ShortSignDictionary.of_phi0 {ζ : ShortIdx → ℂ}
    (hadd : ∀ p q, phi0 (p + q) = phi0 p + phi0 q) (hinj : Function.Injective phi0)
\end{leancode}
(\leanpath{VOA/\allowbreak Moonshine/\allowbreak MoonshineCNGriessShortClasses.lean}; the remaining hypotheses of
\lean{of_phi0} are the three phase conditions, verbatim as in the structure.)
\begin{proof}[Proof on paper]
If $\gamma\equiv\alpha+\beta$, the four vectors $\gamma\pm\alpha\pm\beta$ lie in $2\Leech$ (since
$2\alpha,2\beta\in2\Leech$), whose nonzero vectors have $u\cdot u\ge4\cdot32$.  The parallelogram law
$\sum_\pm\|\gamma\pm\alpha\pm\beta\|^2=4(\|\gamma\|^2+\|\alpha\|^2+\|\beta\|^2)=384$ leaves room for at most
three nonzero terms of size $\ge128$, so one vanishes: $\gamma=\pm(\alpha\pm\beta)$, and $\|\gamma\|^2=32$
forces $\alpha\cdot\beta=\mp16$.
\end{proof}
\begin{physics}
Two charged spin-two fields can fuse into a third charged spin-two field only if their momenta have
inner product $\mp2$ (then $e^{i\alpha X}e^{i\beta X}\sim(z-w)^{-2}e^{i(\alpha+\beta)X}$, a
$(z-w)^{-2}$ term, i.e.\ a Griess product); this lemma says the $\Z/2$ charge arithmetic of
$\Leech/2\Leech$ sees exactly these fusions.
\end{physics}
\begin{status}
No hypotheses beyond those displayed; both audited in \leanpath{MoonshineCNGriessAudit.lean}.
\end{status}

\subsection{The sign refinement and the unconditional untwisted dictionary}

\begin{theorem}[Sign refinements exist]\label{ch21:thm:signref}
Let $V$ be a finite-dimensional $\chXXIF$-vector space and $B$ a symmetric bilinear form on $V$.  There is
$Z:V\to\C$ with $Z(0)=1$ and $Z(x+y)=Z(x)Z(y)(-1)^{B(x,y)}$.
\end{theorem}
\begin{leancode}
theorem exists_sign_refinement {V : Type*} [AddCommGroup V] [Module (ZMod 2) V]
    [Module.Finite (ZMod 2) V] (B : V →ₗ[ZMod 2] V →ₗ[ZMod 2] ZMod 2)
    (hB : ∀ x y, B x y = B y x) :
    ∃ Z : V → ℂ, Z 0 = 1 ∧ ∀ x y, Z (x + y) = Z x * Z y * Extraspecial.sgn (B x y) := by
\end{leancode}
(\leanpath{VOA/\allowbreak Moonshine/\allowbreak MoonshineCNGriessSign.lean}.)
\begin{dictionary}
\lean{Extraspecial.sgn t} $=(-1)^t\in\C$ (Chapter~\ref{ch:06}).  \lean{[Module.Finite (ZMod 2) V]}:
finite dimension.  No other hypothesis.
\end{dictionary}
\begin{proof}[Proof on paper]
Choose a basis $b_1,\dots,b_n$ and put
$Z(x)=\prod_i(i^{B(b_i,b_i)})^{x_i}\cdot(-1)^{\sum_{i<j}x_ix_jB(b_i,b_j)}$, with $\sqrt{-1}$ written $i$
in the base and $x_i\in\{0,1\}$ the coordinates.  The cross terms of $Z(x+y)$ against $Z(x)Z(y)$ produce
$(-1)^{\sum_{i\ne j}x_iy_jB(b_i,b_j)}$ from the second factor and $(-1)^{\sum_ix_iy_iB(b_i,b_i)}$ from
the first ($i^{2}=-1$ when both $x_i=y_i=1$), together $(-1)^{B(x,y)}$ by symmetry.
\end{proof}
\begin{physics}
A ``quadratic refinement'' of a $\Z/2$-valued form: the standard device for fixing phases of vertex
operators (or of fermion zero-mode representations) consistently with given commutation signs.
\end{physics}
\begin{status}
No hypotheses beyond those displayed; audited in \leanpath{MoonshineCNGriessAudit2.lean}.
\end{status}

\begin{definition}[The sign map $\zeta$]\label{ch21:def:griessShortSign}
On $V=\chXXIC\times\chXXIC^*$ let $E(p,q)$ be the exponent of our lattice cocycle,
$\varepsilon(\alpha,\beta)=(-1)^{E(p,q)}$ for $\alpha,\beta$ in the classes $\varphi_0(p),\varphi_0(q)$,
and $\tau(p,q)=\langle q_a,p_\delta\rangle$.  Both have commutator $\langle\alpha,\beta\rangle\bmod2$
(for $E$ by the commutator property of $\varepsilon$; for $\tau$ by Theorem~\ref{ch21:thm:phi0}), so
$B=E+\tau$ is symmetric.  Let $Z$ be a sign refinement of $B$ and $\zeta_c=Z(c)$.
\end{definition}
\begin{leancode}
noncomputable def cocycleExp (p q : MonsterGriess.Code.C × Cocode) : ZMod 2 :=
  latRedLin p ⬝ᵥ (cocMat leech *ᵥ latRedLin q)
\end{leancode}
\begin{leancode}
noncomputable def signForm :
    (MonsterGriess.Code.C × Cocode) →ₗ[ZMod 2] (MonsterGriess.Code.C × Cocode) →ₗ[ZMod 2] ZMod 2 :=
  LinearMap.mk₂ (ZMod 2) (fun p q => cocycleExp p q + pair q.1 p.2)
\end{leancode}
\begin{leancode}
noncomputable def signRefinement : MonsterGriess.Code.C × Cocode → ℂ :=
  (exists_sign_refinement signForm signForm_symm).choose
\end{leancode}
\begin{leancode}
noncomputable def griessShortSign (c : ShortIdx) : ℂ := signRefinement c.1
\end{leancode}
(\leanpath{VOA/\allowbreak Moonshine/\allowbreak MoonshineCNGriessSign.lean}.)
\begin{dictionary}
\lean{latRedLin p} is the reduction modulo $2$, in our lattice basis, of any lattice vector in the
class $\varphi_0(p)$ (through \lean{leechEquiv}); \lean{cocMat leech} is the matrix of the cocycle
(Chapter~\ref{ch:06}: $\varepsilon(\alpha,\beta)=(-1)^{\bar\alpha^{\mathsf T}\mathrm{cocMat}\,\bar\beta}$).
\lean{signForm_symm} proves symmetry.  \lean{signRefinement} is chosen by \lean{Classical.choose}: any
refinement would do, and different choices differ by a character $(-1)^{\ell(x)}$; nothing downstream
depends on which one is chosen except the identification of bases, and all final statements are
invariant.
\end{dictionary}
\begin{physics}
$\zeta$ is the normalisation of the $98280$ vertex operators $\cos(\alpha X)$ that turns the lattice
CFT's OPE signs into AxiomMath's.  That it takes values in $\{\pm1,\pm i\}$ rather than $\{\pm1\}$
reflects that $\varepsilon(\alpha,\alpha)$ can be $-1$ in the chosen cocycle gauge.
\end{physics}
\begin{status}
Definitions; audited in \leanpath{MoonshineCNGriessAudit2.lean}.
\end{status}

\begin{theorem}[The untwisted dictionary, unconditionally]\label{ch21:thm:griessShortSign}
\lean{griessShortSign} is a sign dictionary; hence for all $u,v$ with zero $W$-part the dictionary
$\Phi_\zeta$ (into $\Vnat$) is multiplicative and $\langle\Phi_\zeta u,\Phi_\zeta v\rangle=2b(u,v)$, and its
$\C$-linear extension is an isomorphism $\C\otimes_\Q(\chXXISym\oplus X)\cong(V_\Leech^+)_2$.
\end{theorem}
\begin{leancode}
theorem shortSignDictionary_griessShortSign : ShortSignDictionary griessShortSign :=
\end{leancode}
\begin{leancode}
theorem untwistedGriessNat_gmul {u v : SpaceB} (hu : u.2.2 = 0) (hv : v.2.2 = 0) :
    untwistedDictNat griessShortSign (gmul u v)
      = VnatVOA.griessMul (untwistedDictNat griessShortSign u)
          (untwistedDictNat griessShortSign v) :=
\end{leancode}
\begin{leancode}
theorem untwistedGriessNat_form {u v : SpaceB} (hu : u.2.2 = 0) (hv : v.2.2 = 0) :
    SimpleNat.natForm (untwistedDictNat griessShortSign u : Vnat)
        (untwistedDictNat griessShortSign v : Vnat) = 2 * (bform u v : ℂ) :=
\end{leancode}
(\leanpath{VOA/\allowbreak Moonshine/\allowbreak MoonshineCNGriessSign.lean}) and
\begin{leancode}
noncomputable def untwistedIsoC : ℂ ⊗[ℚ] untwistedSub ≃ₗ[ℂ] leechVOAPlus.grading 2 :=
  (Algebra.TensorProduct.basis ℂ untwistedBasisQ).equiv untwistedBasisC (Equiv.refl _)
\end{leancode}
\begin{leancode}
theorem untwistedIsoC_tmul (u : untwistedSub) :
    (untwistedIsoC (1 ⊗ₜ u) : VPlus) = untwistedDict griessShortSign (u : SpaceB) := by
\end{leancode}
(\leanpath{VOA/\allowbreak Moonshine/\allowbreak MoonshineCNGriessUntwistedBasis.lean}.)
\begin{dictionary}
\lean{untwistedDictNat ζ u} is $(\Phi_\zeta u,0)\in\Vnat_2$.  \lean{untwistedSub} $=\{u:u.2.2=0\}$.
\lean{ℂ ⊗[ℚ] M} is the extension of scalars; \lean{1 ⊗ₜ u} the image of $u$.  \lean{untwistedIsoC}
is built by sending the basis $1\otimes b_k$ of the tensor product to the basis $\Phi(b_k)$ of
$(V_\Leech^+)_2$, and \lean{untwistedIsoC_tmul} says it extends $\Phi$.
\end{dictionary}
\begin{proof}[Proof on paper]
\lean{of_phi0} reduces to \lean{sq}, \lean{sign_neg_two}, \lean{sign_two}.  From
$Z(c)=Z(d)Z(e)(-1)^{E(d,e)+\tau(d,e)}$ for $c=d+e$ and $Z(d),Z(e)\in\{\pm1,\pm i\}$:
$(-1)^{\tau}Z(c)=Z(d)Z(e)\varepsilon(\alpha_d,\alpha_e)$, which is \lean{sign_neg_two} (and \lean{sign_two}
after $\varepsilon(\alpha,-\beta)=\varepsilon(\alpha,\beta)$).  From $1=Z(0)=Z(c+c)=Z(c)^2(-1)^{E(c,c)+\langle a_c,\delta_c\rangle}$ and
$\langle a_c,\delta_c\rangle=0$ for short $c$: $Z(c)^2=\varepsilon(\alpha_c,\alpha_c)$, so
$\zeta_c^2\varepsilon(\alpha_c,\alpha_c)=\varepsilon^2=1$.  For the isomorphism: injectivity because $b$ is
positive definite and $\Phi$ is an isometry up to $2$ (a vanishing $\C$-combination of images, paired
with images of a dual basis, has zero coefficients); surjectivity because $(V_\Leech^+)_2$ is spanned
by the frame bilinears, each a $Q_C$ of a matrix ($F$ invertible), and the $x_\alpha$, each
$\zeta_c^{-1}\Phi(e_c)$ for the short class $c$ of $\pm\alpha$ (bijectivity of $\varphi_0$).
\end{proof}
\begin{physics}
The untwisted sector of the Monster CFT at weight two \emph{is} AxiomMath's $\chXXISym\oplus X$, with its
OPE.
\end{physics}
\begin{status}
No hypotheses.  All audited in \leanpath{MoonshineCNGriessAudit2.lean} (262 \texttt{\#print axioms}
lines in the file).
\end{status}

%======================================================================
\section{The twisted half: the spin module is \texorpdfstring{$T$}{T}}\label{ch21:sec:twisted}

\subsection{AxiomMath's spin module realises the extraspecial module \texorpdfstring{$T$}{T}}

Recall (Chapter~\ref{ch:06}) that $T=$ \lean{TLeechT} is the unique simple module of the twisted
group algebra $A_\Leech=\langle e_\alpha: e_\alpha e_\beta=\varepsilon(\alpha,\beta)e_{\alpha+\beta}\rangle$
(\lean{alg (cocMat leech)}), of dimension $2^{12}=4096$, defined abstractly through Wedderburn--Artin,
and that the twisted weight-two space is $(T^+)_2=\mathfrak h(-\tfrac12)\otimes T$ (Chapter~\ref{ch:08}).

\begin{definition}[Phased spin matrices]\label{ch21:def:spinOp}
$\mathrm{spinOp}(p)=Z(p)\,R(0,p_a,p_\delta)$ for $p\in\chXXIC\times\chXXIC^*$, a complex $I\times I$ matrix.
\end{definition}
\begin{leancode}
noncomputable def spinOp (p : MonsterGriess.Code.C × Cocode) : Matrix SpinIdx SpinIdx ℂ :=
  signRefinement p • qmatR ℂ 0 p.1 p.2
\end{leancode}
(\leanpath{VOA/\allowbreak Moonshine/\allowbreak MoonshineCNGriessSpinModule.lean}.)
\begin{dictionary}
\lean{signRefinement p} is $Z(p)$ of Definition~\ref{ch21:def:griessShortSign}; \lean{•} is scalar
multiplication of a complex matrix.  The phase converts AxiomMath's sign convention
($\tau$) into our cocycle's ($E$), as Theorem~\ref{ch21:thm:spinIso}(a) shows.
\end{dictionary}
\begin{physics}
The gamma matrices of the twisted sector, normalised so that they obey the same OPE signs as the
lattice vertex operators of our construction of $\Vnat$.
\end{physics}
\begin{status}
Definition; audited in \leanpath{MoonshineCNGriessAudit2.lean}.
\end{status}
\begin{theorem}[$\C^I$ is $T$]\label{ch21:thm:spinIso}
(a) $\mathrm{spinOp}(p)\mathrm{spinOp}(q)=(-1)^{E(p,q)}\mathrm{spinOp}(p+q)$.  (b) Hence
$e_\alpha\mapsto\mathrm{spinOp}(\bar\alpha)$ makes $\C^I$ an $A_\Leech$-module; it is simple, and there is an
$A_\Leech$-linear isomorphism $\Psi:\C^I\to T$ with $\Psi(\mathrm{spinOp}(\bar\alpha)v)=e_\alpha\Psi(v)$.
\end{theorem}
\begin{leancode}
theorem spinOp_mul (p q : MonsterGriess.Code.C × Cocode) :
    spinOp p * spinOp q = Extraspecial.sgn (cocycleExp p q) • spinOp (p + q) := by
\end{leancode}
\begin{leancode}
noncomputable def spinIso : SpinRep ≃ₗ[alg (cocMat leech)] TLeechT :=
  Classical.choice (Extraspecial.nonempty_linearEquiv_rep (cocMat leech)
    (isUnit_det_cocMat leech leech_even leech_isUnit_det_gramMat) SpinRep)
\end{leancode}
\begin{leancode}
theorem spinIso_spinOp (α : leech.carrier) (v : SpinRep) :
    spinIso ((spinOp (coordOfRed (red leech α))).mulVec v : SpinIdx → ℂ)
      = eT leech leech_even leech_isUnit_det_gramMat α (spinIso v) := by
\end{leancode}
(\leanpath{VOA/\allowbreak Moonshine/\allowbreak MoonshineCNGriessSpinModule.lean}.)
\begin{dictionary}
\lean{SpinRep} is a type synonym of $\C^I$ carrying the $A_\Leech$-module structure
$e_x\mapsto\mathrm{spinOp}(\mathrm{coordOfRed}\,x)$, where \lean{coordOfRed} is the inverse of the
isomorphism \lean{latRedEquiv} between AxiomMath's coordinates and our $(\Z/2)^{24}$ (reduction mod $2$ in
our lattice basis).  \lean{red leech α} is $\bar\alpha$.  \lean{eT … α} is the action of $e_\alpha$ on $T$
(Chapter~\ref{ch:06}).  \lean{nonempty_linearEquiv_rep} (Chapter~\ref{ch:06}) is uniqueness of the simple
module; \lean{spinIso} is \emph{a} choice of isomorphism, unique up to a nonzero scalar by Schur.  The
scalar is fixed later by the form (\lean{spinDictScale}).
\end{dictionary}
\begin{proof}[Proof on paper]
(a) By the product law $R(0,p)R(0,q)=(-1)^{\langle q_a,p_\delta\rangle}R(0,p+q)$ and
$Z(p)Z(q)=Z(p+q)(-1)^{E(p,q)+\tau(p,q)}$ (signs are self-inverse).  (b) The algebra map
$A_\Leech\to\End(\C^I)$ is injective because $A_\Leech$ is simple, and surjective by dimension
($2^{24}=4096^2$); so $\C^I$ is a simple $A_\Leech$-module, isomorphic to $T$.
\end{proof}
\begin{physics}
The $4096$ twisted ground states of the Leech orbifold carry the unique irreducible representation of
the ``Clifford algebra'' of lattice exponentials; AxiomMath's explicit signed permutation matrices are
one concrete basis for it.  Schur's lemma says the identification is unique up to an overall
normalisation of the twist fields.
\end{physics}
\begin{status}
No hypotheses; audited in \leanpath{MoonshineCNGriessAudit2.lean}.
\end{status}

\begin{theorem}[The invariant form of $T$ in the spin model]\label{ch21:thm:TForm}
There is $\kappa\ne0$ with $B_T(\Psi v,\Psi w)=\kappa\,(v\cdot w)$.
\end{theorem}
\begin{leancode}
theorem TForm_spinIso (v w : SpinIdx → ℂ) :
    TForm (spinIso v) (spinIso w) = spinFormScale * (v ⬝ᵥ w) := by
\end{leancode}
\begin{leancode}
theorem spinFormScale_ne_zero : spinFormScale ≠ 0 := by
\end{leancode}
(\leanpath{VOA/\allowbreak Moonshine/\allowbreak MoonshineCNGriessSpinForm.lean}.)
\begin{proof}[Proof on paper]
$B_T$ satisfies $B_T(e_\alpha t,t')=(-1)^{\langle\alpha,\alpha\rangle/2}B_T(t,e_\alpha t')$ (Chapter~\ref{ch:06}); the
dot product satisfies the same law for the spin matrices, because
$\mathrm{spinOp}(\bar\alpha)^{\mathsf T}=(-1)^{\langle\alpha,\alpha\rangle/2}\mathrm{spinOp}(\bar\alpha)$
(from \lean{qmatR_transpose} and $q(\bar\alpha)=\langle a,\delta\rangle$).  So the Gram matrix of
$(v,w)\mapsto B_T(\Psi v,\Psi w)$ commutes with every spin matrix, hence with all matrices, hence is
scalar; it is nonzero because $B_T$ is nondegenerate.
\end{proof}
\begin{physics}
The two-point function of twist ground states is, in AxiomMath's basis, a multiple of the identity:
the basis is orthogonal for the CFT inner product.
\end{physics}
\begin{status}
No hypotheses; audited in \leanpath{MoonshineCNGriessAudit3.lean}.
\end{status}

\paragraph{Recall: the frame involution and the frame sectors.}  FLM's triality
(Chapters~\ref{ch:11}--\ref{ch:13}) uses the lattice vector $w=8e_0\in\Leech$ (norm $8$).  On $T$ the
operator $J=c^{-1}e_w$, with $c$ a fixed scalar (\lean{sqrtEps leech frameW}) chosen so that $J^2=1$,
is an involution; its eigenspaces $T_0,T_1$ give the \emph{frame sectors} $T_s^+=(S[\theta]\otimes T_s)^+$ of
$(V_\Leech^T)^+$ (\lean{TPframe s}).  The triality $\sigma$ pairs the odd-frame part $V_3^+$ of $V_\Leech^+$ with
one of them.  To compare with AxiomMath's triality $\tau_B$ in Chapter~\ref{ch:22}, one needs to know
where AxiomMath's two halves $\eta=0,1$ of the spin index go.

\begin{theorem}[Frame halves]\label{ch21:thm:frame}
The frame involution $J=c^{-1}e_w$ ($w=8e_0$) of $T$ that defines the frame sectors $T_0^+,T_1^+$ of
$\Vnat$ (Chapters~\ref{ch:11}--\ref{ch:13}) acts in the spin model as $\kappa_J D$, with
$D=\mathrm{diag}((-1)^\eta)$ and $\kappa_J^2=1$; so $\Psi$ maps the half $\chXXIC_0\times\{t\}$ of $\C^I$
into the $\kappa_J(-1)^t$-eigenspace of $J$.
\end{theorem}
\begin{leancode}
theorem spinIso_frameJ (v : SpinRep) :
    frameJL (spinIso v) = frameSign • spinIso (frameDiag *ᵥ (v : SpinIdx → ℂ) : SpinIdx → ℂ) := by
\end{leancode}
\begin{leancode}
theorem spinIso_frameJ_of_support {t : ZMod 2} {v : SpinRep}
    (hv : ∀ i : SpinIdx, i.1 ≠ t → (v : SpinIdx → ℂ) i = 0) :
    frameJL (spinIso v) = (frameSign * sgnR ℂ t) • spinIso v := by
\end{leancode}
(\leanpath{VOA/\allowbreak Moonshine/\allowbreak MoonshineCNGriessFrameBlock.lean}.)
\begin{dictionary}
\lean{frameJL} is $J$ as a linear map; \lean{frameDiag} $=D$; \lean{frameSign} $=\kappa_J=c^{-1}Z(\Omega,\theta_D\Omega)$
with $c=$ \lean{sqrtEps leech frameW}.  Key facts: the class of $w=8e_0$ has AxiomMath coordinates
$(\Omega,\theta_D(\Omega))$ (\lean{coordOfRed_red_frameW}), and $R(0,\Omega,\theta_D\Omega)=D$ because
$\theta(\Omega,\cdot)=0$ on $\chXXIC$ (\lean{θ_ΩC_eq_zero}, a kernel computation on the $12$ basis
words).
\end{dictionary}
\begin{physics}
The two frame sectors of the twisted sector, which FLM's triality exchanges with the untwisted
sectors, are exactly the two halves $\eta=0,1$ of AxiomMath's spin index; this is why AxiomMath's
triality $\tau_B$ (which permutes those halves) can be compared with FLM's $\sigma$ (Chapter~\ref{ch:22}).
\end{physics}
\begin{status}
No hypotheses; audited in \leanpath{MoonshineCNGriessAudit2.lean}.
\end{status}

\subsection{The twisted weight-two space in spin coordinates}

\begin{definition}[Spin coordinates on $(T^+)_2$]\label{ch21:def:spinTopVec}
For a complex $I\times24$ matrix $N$, $\mathrm{spinTopVec}(N)=\sum_ih_{f_i}(-\tfrac12)\otimes\Psi(N_{\cdot i})$,
where $f_i$ is the frame vector of the axis $8e_i$ and $N_{\cdot i}$ the $i$-th column.
\end{definition}
\begin{leancode}
noncomputable def spinTopVec : Matrix SpinIdx (Fin 24) ℂ →ₗ[ℂ] TP where
  toFun N := ∑ i, spinTopState (conwayVec i) (fun s => N s i)
\end{leancode}
(\leanpath{VOA/\allowbreak Moonshine/\allowbreak MoonshineCNGriessTwistedTop.lean}.)
\begin{dictionary}
\lean{spinTopState v w} $=$ \lean{hP v (spinIso w)}, where \lean{hP v t} $=h_v(-\tfrac12)\mathbb1\otimes t\in$
\lean{TP} (Chapter~\ref{ch:13}); \lean{conwayVec i} is
the frame coordinate vector of $8e_i$; \lean{spinIsoC} is \lean{spinIso} as a $\C$-linear map.
\end{dictionary}
\begin{physics}
The state $\sum_i\partial X^{f_i}\sigma_{N_{\cdot i}}$: a first excited twisted state, with oscillator
index $i$ and ground-state wave function given by the $i$-th column of $N$.
\end{physics}
\begin{status}
Definitions; audited in \leanpath{MoonshineCNGriessAudit3.lean}.
\end{status}
\begin{theorem}[Module channels and form on $(T^+)_2$]\label{ch21:thm:twistedTop}
For symmetric $A$ and minimal $\alpha$ with integer coordinates $y$ and $\Sigma_\alpha=\mathrm{spinOp}(\bar\alpha)$:
(a) $Q_C(A)_{(1)}\mathrm{spinTopVec}(N)=\mathrm{spinTopVec}(NA+\tfrac{\tr A}8N)$;
(b) $(x_\alpha)_{(1)}\mathrm{spinTopVec}(N)=\tfrac18\mathrm{spinTopVec}(\Sigma_\alpha N-\tfrac14\Sigma_\alpha Nyy^{\mathsf T})$;
(c) $\langle\mathrm{spinTopVec}N,\mathrm{spinTopVec}M\rangle=-4\kappa\sum_{s,i}N_{si}M_{si}$.
\end{theorem}
\begin{leancode}
theorem topAct_conwayQuad_spinTopVec {A : Matrix (Fin 24) (Fin 24) ℂ} (hA : Aᵀ = A)
    (N : Matrix SpinIdx (Fin 24) ℂ) :
    topAct (conwayQuad A) (spinTopVec N) = spinTopVec (N * A + (trace A / 8) • N) := by
\end{leancode}
\begin{leancode}
theorem topAct_frameX0P_spinTopVec {α : leech.carrier} (hα : leech.form α α = 4)
    (N : Matrix SpinIdx (Fin 24) ℂ) :
    topAct (frameX0P α) (spinTopVec N)
      = (8 : ℂ)⁻¹ • spinTopVec (spinOp (coordOfRed (red leech α)) * N
          - (4 : ℂ)⁻¹ • (spinOp (coordOfRed (red leech α)) * N
              * vecMulVec (fun i => (LeechMatch.vec α i : ℂ)) (fun i => (LeechMatch.vec α i : ℂ)))) := by
\end{leancode}
\begin{leancode}
theorem TPForm_spinTopVec (N M : Matrix SpinIdx (Fin 24) ℂ) :
    TPForm (spinTopVec N) (spinTopVec M)
      = -4 * spinFormScale * ∑ s, ∑ i, N s i * M s i := by
\end{leancode}
(\leanpath{VOA/\allowbreak Moonshine/\allowbreak MoonshineCNGriessTwistedTop.lean}.)
\begin{dictionary}
\lean{topAct u} $=$ \lean{mode (plusModule.YM u) 1}, the mode $u_{(1)}$ of $u\in V_\Leech^+$ acting on
$(V_\Leech^T)^+$; in $\Vnat$, $(u,0)_{(1)}(0,t)=(0,u_{(1)}t)$ (\lean{vnat_smul_untwisted_twisted}).
Compare AxiomMath's table: (a) is $A\cdot N=\tfrac{\tr A}8N+NA$, (b) is $x\cdot N$ with $R_c$ replaced by
$\Sigma_\alpha$, and (c) is $b$ on $W$ up to $-4\kappa$.
\end{dictionary}
\begin{proof}[Proof on paper]
(a) The twisted vertex operator of $\mathrm{ax}_i(-1)\mathrm{ax}_j(-1)\mathbb1$ acts on
$h(-\tfrac12)\otimes t$ by contracting one oscillator (the twisted Heisenberg commutator) and by the
normal-ordering constant $\tfrac1{16}$ per boson, which produces the $\tfrac{\tr A}8$ term.  (b) The twisted
vertex operator of $x_\alpha$ acts on the ground states through $e_\alpha$ ($\Psi$ intertwines it with $\Sigma_\alpha$)
and on the oscillator through the reflection-type term $-\tfrac14 yy^{\mathsf T}$ coming from the
$\alpha(-\tfrac12)$ part of the twisted exponential.  (c) Twisted Fock form times
Theorem~\ref{ch21:thm:TForm}, with $F^{\mathsf T}F=8$.
\end{proof}
\begin{physics}
These are the OPEs of untwisted spin-two fields with twisted spin-two fields: bilinears act on the
oscillator index (and shift the weight by the twisted-sector vacuum energy $\tfrac{24}{16}=\tfrac32$,
the source of $\tr A/8$), exponentials act on the twist ground states by the gamma matrices.
\end{physics}
\begin{status}
No hypotheses beyond those displayed; audited in \leanpath{MoonshineCNGriessAudit3.lean}.
\end{status}

\begin{warning}[title={Trap: the phases of the full dictionary are $\zeta'=Z^{-1}$, not $\zeta=Z$}]
In the $X\cdot W$ channel, $x_{\alpha_c}$ acts on $T$ through $e_{\alpha_c}=\mathrm{spinOp}(c)=Z(c)R_c$, so
$e_c$ must be sent to $\zeta'_cx_{\alpha_c}$ with $\zeta'_cZ(c)=1$, i.e.\ $\zeta'=Z^{-1}$.  This differs from
the phase $\zeta=Z$ of Theorem~\ref{ch21:thm:griessShortSign} by the sign $Z(c)^2=\varepsilon(\alpha_c,\alpha_c)$
(\lean{fullShortSign_eq}), and that sign is not always $+1$ for the cocycle of this library.  Reason: the
library's cocycle (Chapter~\ref{ch:06}) is $\varepsilon(\alpha,\beta)=(-1)^{\bar\alpha^{\mathsf T}C\bar\beta}$ with
$C$ the strictly upper triangular part of the Gram matrix $G$ in a lattice basis, so with coordinates $a_i$
of $\alpha$, $\varepsilon(\alpha,\alpha)=(-1)^{\sum_{i<j}a_ia_jG_{ij}}$, and
$\sum_{i<j}a_ia_jG_{ij}=\tfrac12\bigl(\langle\alpha,\alpha\rangle-\sum_iG_{ii}a_i^2\bigr)$; for
$\langle\alpha,\alpha\rangle=4$ this gives $\varepsilon(\alpha,\alpha)=+1$ iff $4\mid\sum_iG_{ii}a_i^2$
(\lean{griessTrCocycle_eq_one_iff}).  The familiar FLM normalisation
$\varepsilon(\alpha,\alpha)=(-1)^{\langle\alpha,\alpha\rangle/2}$, which would give $+1$ on norm-$4$ vectors, is a
different choice of cocycle in the same cohomology class; this library's cocycle satisfies only the
commutator relation $\varepsilon(\alpha,\beta)\varepsilon(\beta,\alpha)=(-1)^{\langle\alpha,\beta\rangle}$.  Whether
$\varepsilon(\alpha_c,\alpha_c)=1$ for all $c$ in the chosen basis is not decided, and nothing depends on
it.  Using $\zeta$ instead of $\zeta'$ in the full dictionary would make the $X\cdot W$ channel produce
$Z(c)^2R_c$, which is wrong where $\varepsilon(\alpha_c,\alpha_c)=-1$; formulas stated with $\zeta$ are
converted in \S\ref{ch21:sec:transfer}.
\end{warning}

\begin{definition}[Full phases and the $W$-dictionary]\label{ch21:def:spinDict}
$\zeta'_c=\zeta_c^{-1}$.  $N^\pi=(N_{s,\pi i})_{s,i}$ (columns moved to our coordinates, read in $\C$);
$\lambda$ a fixed square root of $-1/(2\kappa)$; $\mathrm{spinDict}(N)=\lambda\,\mathrm{spinTopVec}(N^\pi)$.
\end{definition}
\begin{leancode}
noncomputable def fullShortSign (c : ShortIdx) : ℂ := (griessShortSign c)⁻¹
\end{leancode}
\begin{leancode}
theorem fullShortSign_eq (c : ShortIdx) :
    fullShortSign c = griessShortSign c * latticeCocycle leech (shortVec c) (shortVec c) := by
\end{leancode}
\begin{leancode}
theorem fullShortSign_smul_spinOp (c : ShortIdx) :
    fullShortSign c • spinOp (coordOfRed (red leech (shortVec c))) = (Rmat c).map (Rat.castHom ℂ) := by
\end{leancode}
\begin{leancode}
noncomputable def spinDictScale : ℂ := (-(2 * spinFormScale)⁻¹) ^ (((2 : ℕ) : ℂ)⁻¹)
\end{leancode}
\begin{leancode}
noncomputable def spinDict : SpaceW →+ TP where
  toFun N := spinDictScale • spinTopVec (spinToConway N)
\end{leancode}
(\leanpath{VOA/\allowbreak Moonshine/\allowbreak MoonshineCNGriessSpinDict.lean}.)
\begin{dictionary}
\lean{(Rmat c).map (Rat.castHom ℂ)} is $R_c$ read over $\C$.  \lean{spinToConway N} $=N^\pi$.
\lean{spinDictScale} is the principal square root (complex power $1/2$); \lean{spinDictScale_sq} records
$\lambda^2=-1/(2\kappa)$.  \lean{shortSignDictionary_fullShortSign}: $\zeta'$ is again a sign dictionary, so
Theorem~\ref{ch21:thm:untwisted} applies to it.
\end{dictionary}
\begin{physics}
The first excited twisted states $\partial X^i\sigma_s$, normalised (by $\lambda$) so that their two-point
function is AxiomMath's $2b$, with the bosons relabelled by $\pi$.
\end{physics}
\begin{status}
Definitions; audited in \leanpath{MoonshineCNGriessAudit3.lean}.
\end{status}
\begin{theorem}[The $W$-block: form, module channels, and AxiomMath's invariance]\label{ch21:thm:spinDict}
(a) $\langle\mathrm{spinDict}N,\mathrm{spinDict}M\rangle=2\sum N_{ji}M_{ji}$.
(b) $\mathrm{matDict}(A)_{(1)}\mathrm{spinDict}(N)=\mathrm{spinDict}(\tfrac{\tr A}8N+NA)$ and
$\mathrm{shortDict}_{\zeta'}(x)_{(1)}\mathrm{spinDict}(N)=\mathrm{spinDict}(\tfrac18\sum_cx_c(R_cN-\tfrac14R_cNU_c))$.
(c) On AxiomMath's side, for $u$ with zero $W$-part, $b(N\cdot M,u)=b(M,u\cdot N)$.
\end{theorem}
\begin{leancode}
theorem TPForm_spinDict (N M : SpaceW) :
    TPForm (spinDict N) (spinDict M) = 2 * ((∑ j, ∑ i, N j i * M j i : ℚ) : ℂ) := by
\end{leancode}
\begin{leancode}
theorem topAct_matDict_spinDict (A : symMat) (N : SpaceW) :
    topAct (matDict A) (spinDict N)
      = spinDict ((trace (A : Matrix (Fin 24) (Fin 24) ℚ) / 8) • N + N * (A : Matrix (Fin 24) (Fin 24) ℚ)) := by
\end{leancode}
\begin{leancode}
theorem topAct_shortDict_spinDict (x : SpaceX) (N : SpaceW) :
    topAct (shortDict fullShortSign x) (spinDict N)
      = spinDict ((8 : ℚ)⁻¹ • ∑ c, x c • (Rmat c * N - (4 : ℚ)⁻¹ • (Rmat c * N * Umat c))) := by
\end{leancode}
\begin{leancode}
theorem bform_gmul_spinBlock_spinBlock (N M : SpaceW) {u : SpaceB} (hu : u.2.2 = 0) :
    bform (gmul (spinBlock N) (spinBlock M)) u
      = bform (spinBlock M) (spinBlock (spinBlockAction u N)) := by
\end{leancode}
(\leanpath{VOA/\allowbreak Moonshine/\allowbreak MoonshineCNGriessSpinDict.lean}.)
\begin{dictionary}
\lean{spinBlockAction u N} is the $W$-part of $u\cdot(0,0,N)$, i.e.\
$\tfrac{\tr A}8N+NA+\tfrac18\sum_cx_c(R_cN-\tfrac14R_cNU_c)$.  (c) is proved directly from the formulas
of Definition~\ref{ch21:def:gmul} (using $R_c^{\mathsf T}=R_c$ for short $c$, \lean{Rmat_transpose_of_short},
and $q_c(M)=\tr(MU_c)$); it is a special case of the associativity of $b$ but is proved independently, since
the port does not contain \lean{bform_gmul_assoc}.
\end{dictionary}
\begin{physics}
(a) fixes the normalisation of the twisted fields so that their two-point function matches
AxiomMath's; (b) is the untwisted$\times$twisted OPE in AxiomMath's coordinates; (c) is the
cyclic symmetry of the three-point function $\langle\sigma\sigma\phi\rangle$ of two twist fields and one
untwisted field.
\end{physics}
\begin{status}
No hypotheses beyond those displayed; all audited in \leanpath{MoonshineCNGriessAudit3.lean}.
\end{status}

%======================================================================
\section{The Griess isomorphism \texorpdfstring{$\mathbb C\otimes\mathcal B\cong V^\natural_2$}{C (x) B = V-natural\_2}}\label{ch21:sec:iso}

\subsection{The rational dictionary on all of \texorpdfstring{$\mathcal B$}{B}}

\begin{definition}[The Griess dictionary]\label{ch21:def:griessNat}
$\mathrm{griessNat}(A,x,N)=\bigl(Q_C(A^\pi)+\sum_cx_c\zeta'_cx_{\alpha_c},\ \lambda\sum_ih_{f_i}(-\tfrac12)\otimes\Psi(N_{\cdot\pi i})\bigr)\in\Vnat_2$.
\end{definition}
\begin{leancode}
noncomputable def griessNat : SpaceB →+ VnatVOA.grading 2 where
  toFun u := ⟨(untwistedDict fullShortSign u, spinDict u.2.2), griessNat_mem u⟩
\end{leancode}
\begin{leancode}
theorem griessNat_coe (u : SpaceB) :
    (griessNat u : Vnat) = (untwistedDict fullShortSign u, spinDict u.2.2) := rfl
\end{leancode}
(\leanpath{VOA/\allowbreak Moonshine/\allowbreak MoonshineCNGriessIso.lean}, namespace \lean{VOA.Moonshine.GriessDictionary}.)
\begin{dictionary}
An element of \lean{VnatVOA.grading 2} is a pair $(v,t)\in V_\Leech^+\times(V_\Leech^T)^+$ with a proof
that it has weight $2$ (\lean{griessNat_mem}).  The first component is the untwisted dictionary with
$\zeta'$, the second the $W$-dictionary.  \lean{griessNat_smul}: $\Q$-homogeneous.
\end{dictionary}
\begin{physics}
Every element of AxiomMath's $\chXXIB$ is assigned a definite spin-two field of the Monster CFT, with
rational coefficients in a fixed basis of vertex operators.
\end{physics}
\begin{status}
Definition; audited in \leanpath{MoonshineCNGriessAudit3.lean}.
\end{status}

\begin{theorem}[The dictionary is multiplicative and isometric up to $2$]\label{ch21:thm:griessNat}
For all $u,v\in\chXXIB$: $\mathrm{griessNat}(u\cdot v)=\mathrm{griessNat}(u)\times\mathrm{griessNat}(v)$
and $\langle\mathrm{griessNat}\,u,\mathrm{griessNat}\,v\rangle=2\,b(u,v)$.
\end{theorem}
\begin{leancode}
theorem griessNat_gmul (u v : SpaceB) :
    griessNat (gmul u v) = VnatVOA.griessMul (griessNat u) (griessNat v) := by
\end{leancode}
\begin{leancode}
theorem griessNat_form (u v : SpaceB) :
    SimpleNat.natForm (griessNat u : Vnat) (griessNat v : Vnat) = 2 * (bform u v : ℂ) := by
\end{leancode}
(\leanpath{VOA/\allowbreak Moonshine/\allowbreak MoonshineCNGriessIso.lean}.)
\begin{dictionary}
\lean{VnatVOA.griessMul} is $u_{(1)}v$ on $\Vnat_2$ (\S\ref{ch21:sec:vnatside}); \lean{natForm} the
invariant form of $\Vnat$.  No hypotheses: $u,v$ range over all of $\chXXIB$.
\end{dictionary}
\begin{proof}[Proof on paper]
Write $u=u_0+(0,0,N)$, $v=v_0+(0,0,M)$ with $u_0,v_0$ untwisted.  Expand both sides bilinearly.
\begin{itemize}
\item Untwisted$\times$untwisted: Theorem~\ref{ch21:thm:untwisted} with $\zeta'$.
\item Untwisted$\times$twisted (both orders): in $\Vnat$, $(a,0)_{(1)}(0,t)=(0,a_{(1)}t)$; the
  channels of Theorem~\ref{ch21:thm:spinDict}(b) give exactly $\mathrm{spinDict}(u_0\cdot N)$; the other
  order follows from commutativity of $u_{(1)}v$ on $\Vnat_2$ (\lean{vnat_smul_one_comm}, from $\Vnat_1=0$)
  and of $\cdot$ on the $W$-channels.
\item Twisted$\times$twisted: $t_{(1)}t'$ lies in $V_\Leech^+$ (\lean{vnat_smul_twisted_twisted_snd}).
  It is determined by its pairings with $(V_\Leech^+)_2$, which is spanned by the untwisted images
  (Theorem~\ref{ch21:thm:griessShortSign} with $\zeta'$; \lean{span_fullUntwistedImage}).  For an
  untwisted $w$: by self-adjointness of the Griess product for \lean{natForm}
  (\lean{natForm_smul_one_adjoint}: $\langle x_{(1)}a,b\rangle=\langle a,x_{(1)}b\rangle$ for $x\in\Vnat_2$,
  from the FHL invariance law and $L_1\Vnat_2\subset\Vnat_1=0$) and commutativity,
  $\langle\mathrm{spinDict}N_{(1)}\mathrm{spinDict}M,\Phi w\rangle=\langle\mathrm{spinDict}M,(\Phi w)_{(1)}\mathrm{spinDict}N\rangle
  =\langle\mathrm{spinDict}M,\mathrm{spinDict}(w\cdot N)\rangle=2b(M,w\cdot N)=2b(N\cdot M,w)$,
  the last step by Theorem~\ref{ch21:thm:spinDict}(c); and $2b(N\cdot M,w)=\langle\Phi(N\cdot M),\Phi w\rangle$
  since $N\cdot M$ is untwisted.  Nondegeneracy of the form on $(V_\Leech^+)_2$ gives
  $\mathrm{spinDict}N_{(1)}\mathrm{spinDict}M=\Phi(N\cdot M)$.
\end{itemize}
The form: \lean{natForm} is the sum of the untwisted and twisted forms, and the two pieces are
Theorem~\ref{ch21:thm:untwisted} and Theorem~\ref{ch21:thm:spinDict}(a); mixed pairings vanish on both
sides.
\end{proof}
\begin{physics}
The twist$\times$twist OPE is not computed directly (that would need the twisted-twisted vertex
operators, i.e.\ the hardest part of the construction of $\Vnat$, Chapters~\ref{ch:09}--\ref{ch:14});
instead it is read off from the three-point function $\langle\sigma\sigma\phi\rangle$ by crossing
symmetry, $\langle(\sigma\times\sigma'),\phi\rangle=\langle\sigma',\phi\times\sigma\rangle$, which is
how a physicist would determine it.
\end{physics}
\begin{status}
No hypotheses; audited in \leanpath{MoonshineCNGriessAudit3.lean}.
\end{status}

\subsection{The isomorphism}

\begin{theorem}[\lean{griessIso}: the Griess algebra of $\Vnat$ is AxiomMath's, over $\C$]\label{ch21:thm:griessIso}
There is a $\C$-linear isomorphism $\Theta:\C\otimes_\Q\chXXIB\to\Vnat_2$ with
$\Theta(1\otimes u)=\mathrm{griessNat}(u)$, such that $\Theta$ carries the $\C$-bilinear extension of the
product of $\chXXIB$ to the Griess product of $\Vnat$ and the $\C$-bilinear extension of $2b$ to the
invariant form of $\Vnat$.
\end{theorem}
\begin{leancode}
noncomputable def griessIso : ℂ ⊗[ℚ] SpaceB ≃ₗ[ℂ] VnatVOA.grading 2 :=
  (Algebra.TensorProduct.basis ℂ griessBasisQ).equiv griessBasisC (Equiv.refl _)
\end{leancode}
\begin{leancode}
theorem griessIso_tmul (u : SpaceB) : griessIso (1 ⊗ₜ u) = griessNat u := by
\end{leancode}
\begin{leancode}
theorem griessIso_mul (x y : ℂ ⊗[ℚ] SpaceB) :
    griessIso (LinearMap.BilinMap.baseChange ℂ gmul x y)
      = VnatVOA.griessMul (griessIso x) (griessIso y) := by
\end{leancode}
\begin{leancode}
theorem griessIso_form (x y : ℂ ⊗[ℚ] SpaceB) :
    SimpleNat.natForm (griessIso x : Vnat) (griessIso y : Vnat)
      = 2 * LinearMap.BilinForm.baseChange ℂ bform x y := by
\end{leancode}
(\leanpath{VOA/\allowbreak Moonshine/\allowbreak MoonshineCNGriessIso.lean}.)

\begin{dictionary}
\textbf{Symbols.}  \lean{ℂ ⊗[ℚ] SpaceB} is the tensor product of $\Q$-vector spaces with its
$\C$-structure on the left factor (extension of scalars); \lean{1 ⊗ₜ u} is $1\otimes u$.
\lean{≃ₗ[ℂ]} is a $\C$-linear bijection.  \lean{VnatVOA.grading 2} is $\Vnat_2$, a $\C$-subspace of
\lean{Vnat}; its elements are coerced to \lean{Vnat} by \lean{(· : Vnat)}.
\lean{LinearMap.BilinMap.baseChange ℂ gmul} is the unique $\C$-bilinear map on $\C\otimes\chXXIB$ with
$(a\otimes u,b\otimes v)\mapsto ab\otimes(u\cdot v)$; \lean{LinearMap.BilinForm.baseChange ℂ bform}
likewise for $b$ (a \emph{bilinear}, not Hermitian, extension).

\textbf{How \lean{griessIso} is built.}  \lean{griessBasisQ} is a $\Q$-basis of $\chXXIB$ (Mathlib's
\lean{Module.finBasis}); \lean{griessBasisC} is the family of its \lean{griessNat}-images, shown to be a
$\C$-basis of $\Vnat_2$ (linear independence from Theorem~\ref{ch21:thm:griessNat} and positive
definiteness of $b$, by pairing against a dual basis; spanning by the dimension count
\lean{finrank_spaceB : finrank ℚ SpaceB = finrank ℂ (VnatVOA.grading 2)}, both $196884$).
\lean{Basis.equiv} sends basis to basis.  So \lean{griessIso} is \emph{characterised} by
\lean{griessIso_tmul}: it is the $\C$-linear extension of \lean{griessNat}.  The choice of
\lean{griessBasisQ} does not matter.

\textbf{Hypotheses.}  None.  \lean{griessIso_mul} and \lean{griessIso_form} hold for all
$x,y\in\C\otimes\chXXIB$, not only for pure tensors.

\textbf{Same as the textbook statement?}  The classical statement is ``$\Vnat_2$ with $u_{(1)}v$ and its
invariant form is isomorphic to the complexified Griess algebra''.  The Lean statement is exactly that,
with ``the Griess algebra'' being AxiomMath's rational model and with a specific normalisation:
the form of $\Vnat$ is $2b$.  The factor $2$ is a normalisation of $b$ against
$\langle\mathbb1,\mathbb1\rangle=1$: AxiomMath's identity element $e=(\tfrac14I,0,0)$ goes to
$\omega/2$, and $2b(e,e)=2\cdot\tfrac{24}{16}=3=\langle\omega/2,\omega/2\rangle=\tfrac14\cdot\tfrac c2$.
It is harmless: automorphisms preserve $b$ iff they preserve $2b$.

\textbf{Could it be vacuous?}  The left side is a $196884$-dimensional space with a genuine product
(\lean{gmul} is AxiomMath's definition, \S\ref{ch21:sec:spaceB}); the right side is the weight-two space
of the constructed moonshine module \lean{VnatVOA}, with \lean{griessMul} defined from its vertex
operators (not from the dictionary).  A linear bijection intertwining two products is an algebra
isomorphism; there is no trivial object on either side.  Two ways it could have been weaker were
avoided: the statement is for all of $\C\otimes\chXXIB$ (not only the untwisted part, which was the first
milestone \lean{untwistedIsoC}), and it involves the product of \lean{VnatVOA} itself (not of a
model of $\Vnat_2$ built for the purpose).

\textbf{Which product on $\Vnat$?}  \lean{VnatVOA} is the VOA of Chapter~\ref{ch:14}
(\lean{VnatVOA_Y : VnatVOA.Y = Ynat wbarL}); its twisted$\times$twisted block is the form-adjoint of the
skew intertwiner.  Any other VOA structure on the same space would be a different object; the theorem is
about this one.
\end{dictionary}
\begin{proof}[Proof on paper]
Bilinearity reduces \lean{griessIso_mul} and \lean{griessIso_form} to pure tensors
$a\otimes u$, $b\otimes v$, where they are $ab$ times Theorem~\ref{ch21:thm:griessNat}.
\end{proof}
\begin{physics}
``The $196884$ spin-two fields of the Monster CFT, with their OPE, are Griess's algebra.''  This is the
foundational observation of Frenkel--Lepowsky--Meurman that made the moonshine module a construction
\emph{of the Monster}: once the Griess algebra sits inside $\Vnat$, every symmetry of the Griess algebra
is a candidate symmetry of the CFT, and every symmetry of the CFT acts on the Griess algebra.
\end{physics}
\begin{status}
No hypotheses.  Audited in \leanpath{MoonshineCNGriessAudit3.lean} and again in
\leanpath{MoonshineCNMonsterFinalAudit.lean}; recorded \texttt{\#print axioms}:
\lean{[propext, Classical.choice, Quot.sound]}.
\end{status}

\subsection{Block statements}

\begin{theorem}[Blocks]\label{ch21:thm:blocks}
(a) $\Theta(1\otimes(0,0,N))=(0,\mathrm{spinDict}N)$ and $\Theta(1\otimes u)=(\Phi_{\zeta'}u,0)$ for untwisted $u$.
(b) The $W$-images span $(T^+)_2$.
(c) There is $s_0\in\chXXIF$ such that $W$-vectors supported on the rows $\chXXIC_0\times\{t\}$ go into the
frame sector $T^+_{s_0+t}$, for $t=0,1$.
\end{theorem}
\begin{leancode}
theorem griessIso_tmul_spinBlock (N : SpaceW) :
    (griessIso (1 ⊗ₜ spinBlock N) : Vnat) = (0, spinDict N) := by
\end{leancode}
\begin{leancode}
theorem griessIso_tmul_untwisted {u : SpaceB} (hu : u.2.2 = 0) :
    (griessIso (1 ⊗ₜ u) : Vnat) = (untwistedDict fullShortSign u, 0) := by
\end{leancode}
\begin{leancode}
theorem span_range_spinDict :
    Submodule.span ℂ (Set.range spinDict) = plusModule.weight (2 : ℂ) := by
\end{leancode}
\begin{leancode}
theorem spinDict_mem_TPframe_shift :
    ∃ s₀ : ZMod 2, ∀ (N : SpaceW) (t : ZMod 2), (∀ j i, j.1 ≠ t → N j i = 0) →
      spinDict N ∈ TPframe (s₀ + t) := by
\end{leancode}
(\leanpath{VOA/\allowbreak Moonshine/\allowbreak MoonshineCNGriessIso.lean}.)
\begin{dictionary}
\lean{plusModule.weight (2 : ℂ)} is $(T^+)_2$ (weight-two subspace of the $V_\Leech^+$-module
\lean{TP}); \lean{TPframe s} is the frame sector $T_s^+$ (Chapter~\ref{ch:11}).  The existential
$s_0$ is the exponent with $(-1)^{s_0}=\kappa_J$, the sign of Theorem~\ref{ch21:thm:frame}
(\lean{frameSign_eq_sgn}); it depends on the choices of the scalar $c$ in $J$ and of the refinement $Z$,
and the library does not compute it.  The content of (c) is therefore not a triviality: each half
$\chXXIC_0\times\{t\}$ of $W$ lands in a \emph{single} frame sector, and the two halves land in
\emph{different} sectors ($s_0$ and $s_0+1$), in one of the two possible orders.  This is what
Chapter~\ref{ch:22} uses to compare AxiomMath's $\tau_B$ (which permutes the two halves) with FLM's $\sigma$
(which moves the frame sectors).  So: AxiomMath's $x_B$ ($-1$ on $W$) corresponds to the sector sign $z$ of
$\Vnat$ ($+1$ on $V_\Leech^+$, $-1$ on $(V_\Leech^T)^+$), and the two halves of $W$ permuted by $\tau_B$
are the two frame sectors.
\end{dictionary}
\begin{physics}
The orbifold structure of $\Vnat$ (untwisted/twisted, and the finer frame decomposition used by FLM's
triality) is visible in AxiomMath's coordinates: AxiomMath's parity involution is the orbifold
involution.
\end{physics}
\begin{status}
No hypotheses; audited in \leanpath{MoonshineCNGriessAudit3.lean}.
\end{status}

%======================================================================
\section{The Monster acts on the Griess algebra of \texorpdfstring{$V^\natural$}{V-natural}}\label{ch21:sec:action}

Part~C begins; its Lean names are in \lean{VOA.Moonshine.MonsterAction}.  The automorphism group of
$\Vnat$ is \lean{AutVnat := Aut VnatVOA}; \lean{GriessAutVnat} is the group of $\C$-linear automorphisms
of $\Vnat_2$ preserving \lean{griessMul} and the intrinsic form $u_{(3)}v=\langle u,v\rangle_3\mathbb1$;
\lean{vnatGriessRestrict : AutVnat →* GriessAutVnat} restricts an automorphism to $\Vnat_2$; and
\lean{vnatGriessRestrict_injective} (no hypotheses) says this restriction is injective.  All four are
documented in Chapter~\ref{ch:20} (\leanpath{VOA/\allowbreak Moonshine/\allowbreak AutVnat.lean},
\leanpath{VOA/\allowbreak Moonshine/\allowbreak AutFaithfulAssembly.lean}, \leanpath{VOA/\allowbreak Core/\allowbreak GriessAutomorphism.lean}).

\subsection{The Griess conjugate}

\begin{definition}[Griess conjugate]\label{ch21:def:griessConj}
For $g\in GL_\Q(\chXXIB)$, $\mathrm{griessConj}(g)=\Theta\circ(1_\C\otimes g)\circ\Theta^{-1}\in GL_\C(\Vnat_2)$.
\end{definition}
\begin{leancode}
noncomputable def griessBaseChange (g : SpaceB ≃ₗ[ℚ] SpaceB) : ℂ ⊗[ℚ] SpaceB ≃ₗ[ℂ] ℂ ⊗[ℚ] SpaceB :=
  g.baseChange ℚ ℂ SpaceB SpaceB
\end{leancode}
\begin{leancode}
noncomputable def griessConj (g : SpaceB ≃ₗ[ℚ] SpaceB) :
    VnatVOA.grading 2 ≃ₗ[ℂ] VnatVOA.grading 2 :=
  griessIso.symm ≪≫ₗ griessBaseChange g ≪≫ₗ griessIso
\end{leancode}
\begin{leancode}
theorem griessConj_griessNat (g : SpaceB ≃ₗ[ℚ] SpaceB) (u : SpaceB) :
    griessConj g (griessNat u) = griessNat (g u) := by
\end{leancode}
(\leanpath{VOA/\allowbreak Moonshine/\allowbreak MoonshineCNMonsterActionConj.lean}.)
\begin{dictionary}
\lean{e₁ ≪≫ₗ e₂} is composition of linear equivalences, first $e_1$ then $e_2$, so \lean{griessConj g}
$=\Theta\circ(1\otimes g)\circ\Theta^{-1}$.  \lean{LinearEquiv.baseChange} is $1_\C\otimes g$
(\lean{griessBaseChange_tmul : griessBaseChange g (a ⊗ₜ u) = a ⊗ₜ g u}).  \lean{griessConj_griessNat} says
that on the rational lattice $\mathrm{griessNat}(\chXXIB)\subset\Vnat_2$, \lean{griessConj g} \emph{is} $g$.
\lean{griessConjHom} packages it as a group homomorphism on all of $GL_\Q(\chXXIB)$, injective
(\lean{griessConjHom_injective}).
\end{dictionary}
\begin{physics}
Transport of a symmetry of AxiomMath's algebra to a linear map of the $196884$ spin-two fields of the CFT.
\end{physics}
\begin{status}
Definition; audited in \leanpath{MoonshineCNMonsterActionAudit.lean}.
\end{status}

\begin{theorem}[\lean{AutO} acts on the Griess algebra of $\Vnat$]\label{ch21:thm:griessFormNat}
(a) If $g$ preserves the product (resp.\ $b$), \lean{griessConj g} preserves \lean{griessMul}
(resp.\ \lean{natForm}) on $\Vnat_2$.
(b) On $\Vnat_2$ the intrinsic form equals \lean{natForm}.
(c) Hence $g\mapsto$ \lean{griessConj g} is an injective homomorphism \lean{AutO} $\to$ \lean{GriessAutVnat}.
(d) Restriction $\Aut(\Vnat)\to GL_\C(\Vnat_2)$ is an injective homomorphism.
\end{theorem}
\begin{leancode}
theorem griessConj_griessMul {g : SpaceB ≃ₗ[ℚ] SpaceB} (hg : IsAlgAut g) (u v : VnatVOA.grading 2) :
    griessConj g (VnatVOA.griessMul u v) = VnatVOA.griessMul (griessConj g u) (griessConj g v) := by
\end{leancode}
\begin{leancode}
theorem griessForm_eq_natForm (u v : VnatVOA.grading 2) :
    GriessAut.griessForm VnatVOA_vacuum_ne_zero u v = SimpleNat.natForm (u : Vnat) (v : Vnat) := by
\end{leancode}
\begin{leancode}
noncomputable def griessConjAut : AutO →* GriessAutVnat :=
  MonoidHom.codRestrict (griessConjHom.comp AutO.subtype) GriessAutVnat
    fun g => griessConj_mem_griessAutVnat g.2
\end{leancode}
\begin{leancode}
theorem griessConjAut_injective : Function.Injective griessConjAut := by
\end{leancode}
\begin{leancode}
noncomputable def vnatWeightTwoRestrict : AutVnat →* (VnatVOA.grading 2 ≃ₗ[ℂ] VnatVOA.grading 2) :=
  GriessAutVnat.subtype.comp vnatGriessRestrict
\end{leancode}
\begin{leancode}
theorem vnatWeightTwoRestrict_injective : Function.Injective vnatWeightTwoRestrict :=
\end{leancode}
(\leanpath{VOA/\allowbreak Moonshine/\allowbreak MoonshineCNMonsterActionConj.lean}.)
\begin{dictionary}
\lean{AutO} is used as a type (its elements are pairs of $g$ and a membership proof); \lean{AutO.subtype}
forgets the proof.  \lean{MonoidHom.codRestrict} restricts the codomain to \lean{GriessAutVnat}, using
(a), (b).  \lean{GriessAut.griessForm hvac u v} is the scalar $\langle u,v\rangle_3$ with
$u_{(3)}v=\langle u,v\rangle_3\mathbb1$; \lean{VnatVOA_vacuum_ne_zero} is the proof that $\mathbb1\ne0$
needed to extract it.  (d) composes \lean{vnatGriessRestrict} with the inclusion
\lean{GriessAutVnat.subtype}, and its injectivity is Chapter~\ref{ch:20}'s
\lean{vnatGriessRestrict_injective}.
\end{dictionary}
\begin{proof}[Proof on paper]
(a) Transport by $\Theta$ (\lean{griessIso_mul}, \lean{griessIso_form}).  (b) The FHL invariance law of
\lean{natForm} at weight two, $\langle u_{(3)}v,\mathbb1\rangle=\langle v,u'_{(-1)}\mathbb1\rangle$ with
$u'=(-1)^2e^{L_1}u=u$ (since $L_1\Vnat_2\subset\Vnat_1=0$), together with $u_{(-1)}\mathbb1=u$ and
$\langle\mathbb1,\mathbb1\rangle=1$.  (c) From (a), (b) and injectivity of $g\mapsto1\otimes g$.
\end{proof}
\begin{physics}
Every element of AxiomMath's Monster gives a symmetry of the spin-two sector of the Monster CFT that
preserves its OPE and its two-point function.  Whether it extends to a symmetry of the whole CFT is the
lifting question of the next subsection.
\end{physics}
\begin{status}
No hypotheses; audited in \leanpath{MoonshineCNMonsterActionAudit.lean}.
\end{status}

\subsection{Liftable elements and the lift}

\begin{definition}[Liftable elements]\label{ch21:def:liftable}
$g\in$ \lean{AutO} is \emph{liftable} if some automorphism $h$ of $\Vnat$ restricts on $\Vnat_2$ to
\lean{griessConj g}.  The liftable elements form a subgroup.
\end{definition}
\begin{leancode}
def IsVnatLiftable (g : AutO) : Prop :=
  ∃ h : AutVnat, vnatWeightTwoRestrict h = griessConj (g : SpaceB ≃ₗ[ℚ] SpaceB)
\end{leancode}
\begin{leancode}
noncomputable def liftableSubgroup : Subgroup AutO :=
  vnatGriessRestrict.range.comap griessConjAut
\end{leancode}
\begin{leancode}
theorem mem_liftableSubgroup {g : AutO} : g ∈ liftableSubgroup ↔ IsVnatLiftable g := by
\end{leancode}
(\leanpath{VOA/\allowbreak Moonshine/\allowbreak MoonshineCNMonsterActionLift.lean}.)
\begin{dictionary}
\lean{(g : SpaceB ≃ₗ[ℚ] SpaceB)} is the underlying linear map of $g\in$ \lean{AutO}.
\lean{H.comap f} is the preimage of a subgroup $H$ under a homomorphism $f$; \lean{f.range} the image.
So \lean{liftableSubgroup} $=\{g\in\mathrm{AutO}:\mathrm{griessConjAut}\,g\in\mathrm{image\ of}\ \Aut(\Vnat)\}$,
and \lean{mem_liftableSubgroup} identifies it with \lean{IsVnatLiftable}.  The equation in
\lean{IsVnatLiftable} is equality of $\C$-linear automorphisms of $\Vnat_2$, i.e.\ $h(v)=$
\lean{griessConj g v} for every $v\in\Vnat_2$.
\end{dictionary}
\begin{physics}
$g$ is liftable when the symmetry it defines on the spin-two fields extends to a symmetry of the whole
chiral algebra.  Such an extension, if it exists, is unique, because a symmetry of $\Vnat$ is determined by
its action on $\Vnat_2$; the Lean proof of this (Chapter~\ref{ch:20}) does not assume that $\Vnat$ is
generated by $\Vnat_2$.
\end{physics}
\begin{status}
Definitions; audited in \leanpath{MoonshineCNMonsterActionAudit.lean}.
\end{status}

\begin{theorem}[The Monster lift]\label{ch21:thm:monsterLift}
There is an injective homomorphism $\mathrm{monsterLift}$ from the liftable subgroup to $\Aut(\Vnat)$
such that $\mathrm{monsterLift}(g)$ acts on $\Vnat_2$ as \lean{griessConj g}, in particular as $g$ on
$\mathrm{griessNat}(\chXXIB)$; it is the unique automorphism with this property.  If every element of
a generating set of a subgroup is liftable, so is the subgroup.
\end{theorem}
\begin{leancode}
noncomputable def monsterLift : liftableSubgroup →* AutVnat :=
\end{leancode}
\begin{leancode}
theorem monsterLift_restrict (g : liftableSubgroup) :
    vnatWeightTwoRestrict (monsterLift g) = griessConj ((g : AutO) : SpaceB ≃ₗ[ℚ] SpaceB) :=
\end{leancode}
\begin{leancode}
theorem monsterLift_apply_griessNat (g : liftableSubgroup) (u : SpaceB) :
    monsterLift g (griessNat u : Vnat) = griessNat (((g : AutO) : SpaceB ≃ₗ[ℚ] SpaceB) u) := by
\end{leancode}
\begin{leancode}
theorem monsterLift_unique (g : liftableSubgroup) (h : AutVnat)
    (hh : vnatWeightTwoRestrict h = griessConj ((g : AutO) : SpaceB ≃ₗ[ℚ] SpaceB)) :
    h = monsterLift g :=
\end{leancode}
\begin{leancode}
theorem monsterLift_injective : Function.Injective monsterLift := by
\end{leancode}
\begin{leancode}
theorem subgroupClosure_le_liftable {S : Set AutO} (hS : ∀ g ∈ S, IsVnatLiftable g) :
    Subgroup.closure S ≤ liftableSubgroup :=
\end{leancode}
\begin{leancode}
noncomputable def monsterLiftOfAll (hlift : ∀ g : AutO, IsVnatLiftable g) : AutO →* AutVnat :=
\end{leancode}
(\leanpath{VOA/\allowbreak Moonshine/\allowbreak MoonshineCNMonsterActionLift.lean}.)
\begin{dictionary}
\lean{monsterLift} is defined as the composite of \lean{griessConjAut} (restricted to the liftable subgroup,
whose image lies in the range of \lean{vnatGriessRestrict}) with the inverse of
\lean{vnatGriessRestrict} on its range (\lean{MonoidHom.ofInjective}); the omitted body says exactly
this.  \lean{monsterLiftOfAll hlift} is the same map on all of \lean{AutO} \emph{under the hypothesis}
\lean{hlift} that everything lifts (it is the input to Chapter~\ref{ch:22}'s \lean{monsterEquivAutVnat}).
\lean{subgroupClosure_insert_le_liftable} is the variant for $\langle S\cup\{\tau\}\rangle$, the shape of
$G_{\mathrm{mat}}$.
\end{dictionary}
\begin{proof}[Proof on paper]
Faithfulness of $\Aut(\Vnat)$ on $\Vnat_2$ (Chapter~\ref{ch:20}) makes the restriction an isomorphism onto
its image; \lean{monsterLift} is its inverse composed with $g\mapsto$ \lean{griessConj g}.  Composites of
injective homomorphisms are injective homomorphisms; uniqueness is injectivity of the restriction.
\end{proof}
\begin{physics}
Each liftable Monster element is a well-defined global symmetry of the Monster CFT, uniquely
determined by its action on the spin-two fields.
\end{physics}
\begin{status}
No hypotheses (except \lean{hlift} where displayed); audited in
\leanpath{MoonshineCNMonsterActionAudit.lean}.
\end{status}

\begin{theorem}[Rationality criterion]\label{ch21:thm:rational}
An automorphism $h$ of $\Vnat$ is $\mathrm{monsterLift}(g)$ for some liftable $g$ iff $h$ maps the rational
lattice $\mathrm{griessNat}(\chXXIB)$ into itself.  If every element of \lean{AutO} lifts and every
automorphism of $\Vnat$ is rational in this sense, then $\Aut(\Vnat)\cong$ \lean{AutO}.
\end{theorem}
\begin{leancode}
theorem isVnatLiftable_of_rational (h : AutVnat)
    (hrat : ∀ u : SpaceB, ∃ u' : SpaceB, h (griessNat u : Vnat) = griessNat u') :
    ∃ g : liftableSubgroup, monsterLift g = h := by
\end{leancode}
\begin{leancode}
theorem mem_range_monsterLift_iff (h : AutVnat) :
    h ∈ monsterLift.range ↔ ∀ u : SpaceB, ∃ u' : SpaceB, h (griessNat u : Vnat) = griessNat u' := by
\end{leancode}
\begin{leancode}
theorem autVnat_mulEquiv_autO_of_liftable_of_rational (hlift : ∀ g : AutO, IsVnatLiftable g)
    (hrat : ∀ h : AutVnat, ∀ u : SpaceB, ∃ u' : SpaceB, h (griessNat u : Vnat) = griessNat u') :
    Nonempty (AutVnat ≃* AutO) := by
\end{leancode}
(\leanpath{VOA/\allowbreak Moonshine/\allowbreak MoonshineCNMonsterActionLift.lean}.)
\begin{dictionary}
\lean{hrat} quantifies over all $u$: $h$ maps each rational vector to a rational vector.
\lean{Nonempty (AutVnat ≃* AutO)}: there exists a group isomorphism (the statement does not name it).
Both hypotheses of the last theorem are genuine inputs at this point; Chapter~\ref{ch:22} discharges
\lean{hrat} without hypotheses (\lean{autVnat_rational}) and replaces \lean{hlift} by
\lean{GmatEqAutO}.
\end{dictionary}
\begin{proof}[Proof on paper]
If $h$ preserves $\mathrm{griessNat}(\chXXIB)$, define $r:\chXXIB\to\chXXIB$ by
$\mathrm{griessNat}(r u)=h(\mathrm{griessNat}\,u)$ (\lean{rationalPart}; \lean{griessNat} is injective).
$r$ is $\Q$-linear and injective, hence bijective; it preserves the product and $b$ because $h$
preserves \lean{griessMul} and \lean{natForm} and \lean{griessNat} is an isometric algebra map
(Theorem~\ref{ch21:thm:griessNat}); so $r\in$ \lean{AutO}.  The restriction of $h$ and
\lean{griessConj r} agree on the spanning set $\mathrm{griessNat}(\chXXIB)$, hence are equal.  The
converse is \lean{monsterLift_apply_griessNat}.  With \lean{hlift}, \lean{monsterLiftOfAll} is defined;
\lean{hrat} makes it surjective.
\end{proof}
\begin{physics}
A symmetry of the CFT comes from AxiomMath's Monster exactly when it preserves the rational structure of
the spin-two fields in the vertex-operator basis.  Since the Monster's character values on $196884$ are
rational, this is expected for all symmetries; proving it is the converse direction of Chapter~\ref{ch:22}.
\end{physics}
\begin{status}
Conditional only on the displayed hypotheses; audited in \leanpath{MoonshineCNMonsterActionAudit.lean}.
\end{status}

\subsection{The extraspecial elements lift, explicitly}

\begin{theorem}[Extraspecial lifts]\label{ch21:thm:extra}
For every $\mu\in\Leech$ there are a sign $\sigma\in\{\pm1\}$ and a liftable $g\in$ \lean{AutO} whose lift is
the extraspecial automorphism \lean{extraVnatAut μ} of $\Vnat$, and $g(A,x,N)=(A,\ \varepsilon_\mu x,\ \sigma R(0,\bar\mu)N)$
with $(\varepsilon_\mu x)_c=(-1)^{\langle\mu,\alpha_c\rangle}x_c$.
\end{theorem}
\begin{leancode}
theorem exists_monsterLift_eq_extraVnatAut (μ : leech.carrier) :
    ∃ σ : ℚ, (σ = 1 ∨ σ = -1) ∧ ∃ g : liftableSubgroup, monsterLift g = extraVnatAut μ ∧
      ∀ u : SpaceB, ((g : AutO) : SpaceB ≃ₗ[ℚ] SpaceB) u
        = (u.1, liftBlockExtraX μ u.2.1,
            (σ • qmatR ℚ 0 (coordOfRed (red leech μ)).1 (coordOfRed (red leech μ)).2) * u.2.2) := by
\end{leancode}
\begin{leancode}
theorem extraVnatAut_griessNat (μ : leech.carrier) (σ : ℚ)
    (hσ : (σ : ℂ) = extraPlusFormScalar μ * signRefinement (coordOfRed (red leech μ)))
    (u : SpaceB) :
    extraVnatAut μ (griessNat u : Vnat)
      = griessNat (u.1, liftBlockExtraX μ u.2.1,
          (σ • qmatR ℚ 0 (coordOfRed (red leech μ)).1 (coordOfRed (red leech μ)).2) * u.2.2) := by
\end{leancode}
(\leanpath{VOA/\allowbreak Moonshine/\allowbreak MoonshineCNMonsterActionExtra.lean}.)
\begin{dictionary}
\lean{extraVnatAut μ} (\leanpath{MoonshineLRLiftExtraAut.lean}, Chapter~\ref{ch:22}) acts on $V_\Leech^+$
by $x_\alpha\mapsto(-1)^{\langle\mu,\alpha\rangle}x_\alpha$, oscillators fixed, and on $(V_\Leech^T)^+$ by
$c_\mu e_\mu$ with $c_\mu^2=(-1)^{\langle\mu,\mu\rangle/2}\varepsilon(\mu,\mu)$
(\lean{extraPlusFormScalar μ} $=c_\mu$).  \lean{liftBlockExtraX μ x} $=(c\mapsto(-1)^{\langle\mu,\alpha_c\rangle}x_c)$.
In spin coordinates $e_\mu=\mathrm{spinOp}(\bar\mu)=Z(\bar\mu)R(0,\bar\mu)$, so the $W$-block is multiplied
by $\sigma=c_\mu Z(\bar\mu)$, and $\sigma^2=1$ (\lean{extraSpinPhase_mul_self}): the action is rational.
\lean{extraLiftable μ} names the element and \lean{extraLiftable_apply} gives its action with
\lean{extraSign μ}.
\end{dictionary}
\begin{proof}[Proof on paper]
Compute \lean{extraVnatAut μ} on \lean{griessNat u} block by block (the formula displayed), observe that the
result is again \lean{griessNat} of a rational vector, and apply the rationality criterion.
\end{proof}
\begin{physics}
The extraspecial elements are the ``momentum-shift'' symmetries $e^{\pi i\langle\mu,p\rangle}$ of the
lattice CFT, extended to the twisted sector by the action of $e_\mu$ on the twist ground states.  They
generate the normal subgroup $2^{1+24}$ of $C_{\mathbb M}(z)$; this theorem shows each is a lift of an
explicit element of AxiomMath's Monster.
\end{physics}
\begin{status}
No hypotheses (\lean{extraVnatAut_griessNat} has the displayed hypothesis \lean{hσ}, satisfied by
\lean{extraSign}); audited in \leanpath{MoonshineCNMonsterActionAudit.lean}.
\end{status}

\subsection{Conway--Norton for liftable elements}

\begin{theorem}[Conway--Norton in Monster language]\label{ch21:thm:cnLiftable}
(a) For every liftable $g$, the McKay--Thompson series of $\mathrm{monsterLift}(g)$ is a normalised
Hauptmodul of a genus-zero group containing some $\Gamma_0(N)$, and its $q$-coefficient is
$\tr(g\mid\chXXIB)\in\Q$.
(b) Let $S\subseteq$ \lean{AutO} consist of liftable elements and $\tau$ be liftable; if $S\cup\{\tau\}$
generates \lean{AutO}, then every $g\in$ \lean{AutO} acts on $\Vnat$ by some $h$ restricting to
\lean{griessConj g}, and the McKay--Thompson series of $h$ is such a Hauptmodul.
\end{theorem}
\begin{leancode}
theorem conwayNorton_liftable (g : liftableSubgroup) : ConwayNortonFor (monsterLift g) :=
  conwayNorton (monsterLift g)
\end{leancode}
\begin{leancode}
theorem mckayThompson_monsterLift_coeff_one (g : liftableSubgroup) :
    PowerSeries.coeff 1 (mckayThompson (monsterLift g))
      = ((LinearMap.trace ℚ SpaceB ((g : AutO) : SpaceB ≃ₗ[ℚ] SpaceB).toLinearMap : ℚ) : ℂ) := by
\end{leancode}
\begin{leancode}
theorem conwayNorton_monster_of_isVnatLiftable (g : AutO) (hg : IsVnatLiftable g) :
    ∃ h : AutVnat, vnatWeightTwoRestrict h = griessConj (g : SpaceB ≃ₗ[ℚ] SpaceB) ∧
      (∀ h' : AutVnat, vnatWeightTwoRestrict h' = griessConj (g : SpaceB ≃ₗ[ℚ] SpaceB) → h' = h) ∧
      PowerSeries.coeff 1 (mckayThompson h)
        = ((LinearMap.trace ℚ SpaceB (g : SpaceB ≃ₗ[ℚ] SpaceB).toLinearMap : ℚ) : ℂ) ∧
      ∃ (Γ : Subgroup (GL (Fin 2) ℝ)) (N : ℕ), IsMoonshineGroup Γ N ∧
        IsNormalizedHauptmodul Γ (mckayThompsonFun h) := by
\end{leancode}
\begin{leancode}
theorem conwayNorton_autO_of_generators {Gvis : Set AutO} {τ : AutO}
    (hvis : ∀ g ∈ Gvis, IsVnatLiftable g) (hτ : IsVnatLiftable τ)
    (hGmat : Subgroup.closure (insert τ Gvis) = ⊤) (g : AutO) :
    ∃ h : AutVnat, vnatWeightTwoRestrict h = griessConj (g : SpaceB ≃ₗ[ℚ] SpaceB) ∧
      ∃ (Γ : Subgroup (GL (Fin 2) ℝ)) (N : ℕ), IsMoonshineGroup Γ N ∧
        IsNormalizedHauptmodul Γ (mckayThompsonFun h) := by
\end{leancode}
(\leanpath{VOA/\allowbreak Moonshine/\allowbreak MoonshineCNMonsterActionMoonshine.lean}.)
\begin{dictionary}
\lean{ConwayNortonFor}, \lean{IsMoonshineGroup}, \lean{IsNormalizedHauptmodul}, \lean{mckayThompson},
\lean{mckayThompsonFun} and the theorem \lean{conwayNorton : ConwayNortonConjecture} are those of
Chapters~\ref{ch:16} and~\ref{ch:20} (\leanpath{MoonshineCNStatement.lean}, \leanpath{MoonshineCNTrace.lean}):
\lean{PowerSeries.coeff n (mckayThompson g)} is $\tr(g\mid\Vnat_{n+1})$, so \lean{coeff 1} is the trace on
$\Vnat_2$; \lean{IsMoonshineGroup Γ N} says $\Gamma\le SL_2(\mathbb{R})$ contains $\Gamma_0(N)$ and is commensurable
with $SL_2(\Z)$.  \lean{LinearMap.trace ℚ SpaceB} is the trace over $\Q$; the cast to $\C$ is explicit.

\textbf{A naming hazard.}  In \lean{conwayNorton_autO_of_generators} the bound variable is called
\lean{Gvis}, but it is an \emph{arbitrary} set of elements of \lean{AutO}, not AxiomMath's
\lean{Visible.Gvis} (which is a subgroup of $GL_\Q(\chXXIB)$, not a set of elements of \lean{AutO}).  The
theorem is a general statement about any generating set; Chapter~\ref{ch:22} applies it to the image of
$G_{\mathrm{vis}}$.  The name is suggestive, not a hypothesis.

\textbf{What is unconditional.}  (a) has the single hypothesis that $g$ lifts; (b) has \lean{hvis},
\lean{hτ}, \lean{hGmat}.  The $q$-coefficient statement makes the series an invariant of $g$ itself:
\lean{mckayThompsonFun_eq_of_restrict} shows the series does not depend on which lift is used (there is
only one).
\end{dictionary}
\begin{proof}[Proof on paper]
(a) \lean{conwayNorton} holds for every automorphism of $\Vnat$ (Chapter~\ref{ch:20}), in particular for
$\mathrm{monsterLift}(g)$.  The trace on $\Vnat_2$ is the trace of \lean{griessConj g}, which is the trace of
$1\otimes g$ (conjugation), which is the trace of $g$ (base change preserves traces):
\lean{trace_griessConj}.  (b) By \lean{subgroupClosure_insert_le_liftable}, every $g$ is liftable.
\end{proof}
\begin{physics}
For each liftable Monster element, the twisted partition function $\tr(g\,q^{L_0-1})$ of the Monster CFT
is a genus-zero Hauptmodul, and its first nontrivial coefficient is the character of $g$ on $196884$
computed in AxiomMath's rational model.  This is Conway and Norton's original formulation, which starts
from the Monster's character table.
\end{physics}
\begin{status}
Conditional only on the displayed hypotheses; audited in \leanpath{MoonshineCNMonsterActionAudit.lean};
standard axioms.  The hypotheses are discharged in Chapter~\ref{ch:22}: every element of $G_{\mathrm{mat}}$
lifts (no hypotheses), and every element of \lean{AutO} lifts given \lean{GmatEqAutO}.
\end{status}

%======================================================================
\section{Transfer of the block formulas to the full dictionary}\label{ch21:sec:transfer}

The block formulas for the $\mathrm{Co}_0$-lifts, the extraspecial automorphisms and FLM's triality on
$\Vnat_2$ (Chapter~\ref{ch:22}) are stated against the untwisted dictionary with the phase $\zeta$; the full
dictionary uses $\zeta'$ (see the trap in \S\ref{ch21:sec:twisted}).  This short section records the
conversion; its mathematical content is the diagonal sign change $D$.  \leanpath{VOA/\allowbreak Moonshine/\allowbreak MoonshineLRGriessTransfer.lean} (namespace
\lean{VOA.Moonshine.ConwayNorton.M8}) converts them.

\begin{definition}[The diagonal change of short coordinates]\label{ch21:def:griessTrX}
$d_c=\varepsilon(\alpha_c,\alpha_c)\in\{\pm1\}$, as a rational number; $D(x)_c=d_cx_c$.
\end{definition}
\begin{leancode}
def griessTrSign (c : ShortIdx) : ℚ := MonsterGriess.sgn (cocycleExp c.1 c.1)
\end{leancode}
\begin{leancode}
def griessTrX (x : SpaceX) : SpaceX := fun c => griessTrSign c * x c
\end{leancode}
(\leanpath{VOA/\allowbreak Moonshine/\allowbreak MoonshineLRGriessTransfer.lean}.)
\begin{physics}
A sign flip of those vertex operators $\cos(\alpha_cX)$ whose cocycle value $\varepsilon(\alpha_c,\alpha_c)$
is $-1$: a change of basis of the spin-two fields, not a change of the theory.
\end{physics}
\begin{status}
Definitions; audited in \leanpath{MoonshineLRLiftAudit.lean}.
\end{status}
\begin{theorem}[Transfer]\label{ch21:thm:transfer}
(a) $D$ is an involution and $\mathrm{shortDict}_{\zeta'}(x)=\mathrm{shortDict}_\zeta(Dx)$; hence
$\mathrm{griessNat}(A,x,0)=\Phi_\zeta(A,Dx,0)$.
(b) $\varepsilon(\alpha_c,\alpha_c)=1$ iff $4$ divides the diagonal Gram contribution
$\sum_i\langle B_i,B_i\rangle(\alpha_c)_i^2$ (coordinates in the chosen lattice basis $B$).
(c) The $\mathrm{Co}_0$-lifts act on $\Theta(1\otimes(A,x,0))$ by $(A,x)\mapsto(gAg^{\mathsf T},D\,P\,D\,x)$ with $P$ the
signed permutation computed with $\zeta$; the extraspecial lifts act on $X$ by
$x_c\mapsto(-1)^{\langle\mu,\alpha_c\rangle}x_c$ unchanged; FLM's $\sigma$ on the even part is the old formula
conjugated by $D$.
\end{theorem}
\begin{leancode}
theorem griessTrNat (A : symMat) (x : SpaceX) :
    griessNat (A, x, 0) = untwistedDictNat griessShortSign (A, griessTrX x, 0) := by
\end{leancode}
\begin{leancode}
theorem griessTrCocycle_eq_one_iff (c : ShortIdx) :
    latticeCocycle leech (shortVec c) (shortVec c) = 1 ↔ 4 ∣ griessTrDiagonal c := by
\end{leancode}
\begin{leancode}
theorem griessTrLiftSignNat_untwisted (A : symMat) (x : SpaceX) :
    liftTwVnatAut ι hι hsurj (griessIso (1 ⊗ₜ[ℚ] (A, x, 0)) : Vnat) =
      (griessIso (1 ⊗ₜ[ℚ]
        (liftBlockSymMat ι A, griessTrX (liftSignX ι hι (griessTrX x)), 0)) : Vnat) := by
\end{leancode}
(\leanpath{VOA/\allowbreak Moonshine/\allowbreak MoonshineLRGriessTransfer.lean}.)
\begin{dictionary}
\lean{cocycleExp c.1 c.1} is the exponent $E(c,c)$, so \lean{griessTrSign c} $=\varepsilon(\alpha_c,\alpha_c)$
(\lean{griessTrSign_cast}).  In the isometry section, \lean{ι : leech.carrier →ₗ[ℤ] leech.carrier} with
\lean{hι} (isometry) and \lean{hsurj} (surjective) are section variables: $\iota\in\mathrm{Co}_0$.
\lean{liftTwVnatAut}, \lean{liftBlockSymMat}, \lean{liftSignX}, \lean{sigmaAut}, \lean{sigmaBlockOFull},
\lean{sigmaBlockParity} are constructions of Chapter~\ref{ch:22}; here only their transport matters.
The extraspecial and triality parts of (c) are \lean{griessTrExtraNat_shortBlock} and
\lean{griessTrSigmaBlockOFull_nat} (the latter for $x$ supported on the even-parity short classes, those
in the frame sector $V_0^+$); (a) also contains \lean{griessTrShortDict}.  The diagonal entries of $D$ are
exactly the diagonal cocycle values (\lean{griessTrSign_cast}), not an assumed normalisation.
\end{dictionary}
\begin{proof}[Proof on paper]
(a) $\zeta'_c=\zeta_c\varepsilon(\alpha_c,\alpha_c)$ (\lean{fullShortSign_eq}); $d_c^2=1$.  (b) The cocycle
exponent satisfies $2E(\alpha,\alpha)+\sum_i\langle B_i,B_i\rangle a_i^2=\langle\alpha,\alpha\rangle=4$
(\lean{cocExp_add_swap} with the diagonal written out), so $E(\alpha,\alpha)$ is even iff $4$ divides the
diagonal sum.  (c) Substitute (a) into the formulas of Chapter~\ref{ch:22}.
\end{proof}
\begin{physics}
A change of normalisation of the $98280$ vertex operators $\cos(\alpha X)$ by signs; all symmetry formulas
are conjugated by this diagonal change.
\end{physics}
\begin{status}
No hypotheses beyond those displayed; audited in \leanpath{MoonshineLRLiftAudit.lean}.  Open but
unneeded: whether $\varepsilon(\alpha_c,\alpha_c)=1$ for all $c$.
\end{status}

%======================================================================
\section{The assigned files and how each is covered}\label{ch21:sec:filemap}

``Individually'' means that the file's headline declarations are presented with the Fidelity Protocol;
``collectively'' that the file's role is described in the text indicated.

\begin{longtable}{p{0.36\textwidth}p{0.18\textwidth}p{0.38\textwidth}}
\toprule
file (\leanpath{VOA/\allowbreak Moonshine/\allowbreak …}) & coverage & where / role\\\midrule
\endhead
\leanpath{MonsterGriess/\allowbreak Golay.lean} & individually & \S\ref{ch21:sec:golay}: $\chXXIC$, $\theta$\\
\leanpath{MonsterGriess/\allowbreak CodeCocycle.lean} & individually & \S\ref{ch21:sec:golay}: cocycle identities\\
\leanpath{MonsterGriess/\allowbreak Spin.lean} & individually & \S\ref{ch21:sec:golay}, \S\ref{ch21:sec:spin}: cocode, pairing, $R(s,a,\delta)$\\
\leanpath{MonsterGriess/\allowbreak SpinMatrixLaws.lean} & individually & Thm~\ref{ch21:thm:qmatRmul}\\
\leanpath{MonsterGriess/\allowbreak SpinMatrixMap.lean} & collectively & \S\ref{ch21:sec:spin}: change of ring (\lean{map_qmatR})\\
\leanpath{MonsterGriess/\allowbreak Lattice.lean} & individually & Def~\ref{ch21:def:codeLattice}\\
\leanpath{MonsterGriess/\allowbreak LeechMatch.lean} & individually & Thms~\ref{ch21:thm:golaymatch}, \ref{ch21:thm:leechmatch}\\
\leanpath{MonsterGriess/\allowbreak MinimalVectors.lean} & individually & Thm~\ref{ch21:thm:minvecs}(a,b)\\
\leanpath{MonsterGriess/\allowbreak LatticeFamilies.lean} & individually & Thm~\ref{ch21:thm:minvecs}(c,d)\\
\leanpath{MonsterGriess/\allowbreak Short.lean} & individually & Def~\ref{ch21:def:short}\\
\leanpath{MonsterGriess/\allowbreak LatticeQuadratic.lean} & individually & Thm~\ref{ch21:thm:phi0}\\
\leanpath{MonsterGriess/\allowbreak SpinGroupQ.lean} & individually & Def~\ref{ch21:def:Q}\\
\leanpath{MonsterGriess/\allowbreak SpinExtraspecial.lean} & individually & \S\ref{ch21:sec:spin}: $|Q|$, Frattini, $\phi$, dyadic normalisers\\
\leanpath{MonsterGriess/\allowbreak Algebra.lean} & individually & Defs~\ref{ch21:def:spaceB}, \ref{ch21:def:gmul}\\
\leanpath{MonsterGriess/\allowbreak AutO.lean} & individually & Def~\ref{ch21:def:AutO}\\
\leanpath{MonsterGriess/\allowbreak GriessTrilinear.lean} & individually & Def~\ref{ch21:def:T}, Thm~\ref{ch21:thm:Tunique}, $x_B$\\
\leanpath{MonsterGriess/\allowbreak ScalarKernel.lean} & individually & Def~\ref{ch21:def:latticeO}\\
\leanpath{MonsterGriess/\allowbreak VisiblePair.lean} & individually & Defs~\ref{ch21:def:Pvis}, \ref{ch21:def:Gvis}\\
\leanpath{MonsterGriess/\allowbreak VisibleAction.lean} & individually & Thm~\ref{ch21:thm:gvis}\\
\leanpath{MonsterGriess/\allowbreak ParityTraceSplitting.lean} & collectively & proof of Thm~\ref{ch21:thm:parity}(a)\\
\leanpath{MonsterGriess/\allowbreak ParityRealExtension.lean} & collectively & proof of Thm~\ref{ch21:thm:parity}(b)\\
\leanpath{MonsterGriess/\allowbreak ParityShortLines.lean} & collectively & proof of Thm~\ref{ch21:thm:parity}(c)\\
\leanpath{MonsterGriess/\allowbreak ParityEqualsVisible.lean} & individually & Thm~\ref{ch21:thm:parity}\\
\leanpath{MonsterGriess/\allowbreak OctadEvenSubset.lean} & collectively & \S\ref{ch21:sec:triality}: octad quadratic spaces, $H$, $D$\\
\leanpath{MonsterGriess/\allowbreak OctadThreeOctad.lean} & collectively & proof of Thm~\ref{ch21:thm:tauBgmul} (octad${}^3$)\\
\leanpath{MonsterGriess/\allowbreak TrialityLabels.lean} & collectively & Def~\ref{ch21:def:tbasis}: labels \lean{rMinus}, \dots\\
\leanpath{MonsterGriess/\allowbreak TrialityClasses.lean} & collectively & Def~\ref{ch21:def:tbasis}: \lean{labelEquiv}\\
\leanpath{MonsterGriess/\allowbreak TrialityBasis.lean} & individually & Def~\ref{ch21:def:tbasis}, dimensions\\
\leanpath{MonsterGriess/\allowbreak TrialityMatrix.lean} & individually & Def~\ref{ch21:def:tauB}\\
\leanpath{MonsterGriess/\allowbreak TrialityFourGroup.lean} & individually & Def~\ref{ch21:def:Gmat}, Thm~\ref{ch21:thm:fourgroup}\\
\leanpath{MonsterGriess/\allowbreak GriessGrading.lean} & collectively & proof of Thm~\ref{ch21:thm:tauBgmul}(ii)\\
\leanpath{MonsterGriess/\allowbreak InvarianceDiagonal.lean} & collectively & proof of Thm~\ref{ch21:thm:tauBgmul}(i)\\
\leanpath{MonsterGriess/\allowbreak InvarianceTwoCoord.lean} & collectively & proof of Thm~\ref{ch21:thm:tauBgmul}(iii)\\
\leanpath{MonsterGriess/\allowbreak InvarianceXYZ.lean} & collectively & proof of Thm~\ref{ch21:thm:tauBgmul}(iii)\\
\leanpath{MonsterGriess/\allowbreak InvarianceOddTwoCoord.lean} & collectively & proof of Thm~\ref{ch21:thm:tauBgmul}(iii)\\
\leanpath{MonsterGriess/\allowbreak InvarianceOctadBasic.lean} & collectively & proof of Thm~\ref{ch21:thm:tauBgmul}(iii)\\
\leanpath{MonsterGriess/\allowbreak InvarianceOctadTwoCoord.lean} & collectively & proof of Thm~\ref{ch21:thm:tauBgmul}(iii)\\
\leanpath{MonsterGriess/\allowbreak InvarianceThreeOctad.lean} & collectively & proof of Thm~\ref{ch21:thm:tauBgmul}(iii)\\
\leanpath{MonsterGriess/\allowbreak InvarianceOddValues.lean} & collectively & proof of Thm~\ref{ch21:thm:tauBgmul}(iii)\\
\leanpath{MonsterGriess/\allowbreak InvarianceOddOctad.lean} & collectively & proof of Thm~\ref{ch21:thm:tauBgmul}(iii)\\
\leanpath{MonsterGriess/\allowbreak ProductInvariance.lean} & individually & Thm~\ref{ch21:thm:tauBgmul}\\
\leanpath{MonsterGriess/\allowbreak ParityAudit.lean} & audit & \S\ref{ch21:sec:audits}\\
\leanpath{MonsterGriess/\allowbreak TrialityAudit.lean} & audit & \S\ref{ch21:sec:audits}\\
\leanpath{MonsterGriess/\allowbreak ProductInvarianceAudit.lean} & audit & \S\ref{ch21:sec:audits}\\
\leanpath{MoonshineCNGriessTable.lean} & individually & \S\ref{ch21:sec:spaceB} (table)\\
\leanpath{MoonshineCNGriessConway.lean} & individually & Def~\ref{ch21:def:conwayQuad}, Thm~\ref{ch21:thm:conway}\\
\leanpath{MoonshineCNGriessUntwisted.lean} & individually & Defs~\ref{ch21:def:dict}, \ref{ch21:def:ssd}, Thm~\ref{ch21:thm:untwisted}\\
\leanpath{MoonshineCNGriessShortClasses.lean} & individually & Thm~\ref{ch21:thm:ofphi0}\\
\leanpath{MoonshineCNGriessSign.lean} & individually & Thm~\ref{ch21:thm:signref}, Def~\ref{ch21:def:griessShortSign}, Thm~\ref{ch21:thm:griessShortSign}\\
\leanpath{MoonshineCNGriessUntwistedBasis.lean} & individually & Thm~\ref{ch21:thm:griessShortSign} (\lean{untwistedIsoC})\\
\leanpath{MoonshineCNGriessSpinModule.lean} & individually & Thm~\ref{ch21:thm:spinIso}\\
\leanpath{MoonshineCNGriessSpinForm.lean} & individually & Thm~\ref{ch21:thm:TForm}\\
\leanpath{MoonshineCNGriessFrameBlock.lean} & individually & Thm~\ref{ch21:thm:frame}\\
\leanpath{MoonshineCNGriessTwistedTop.lean} & individually & Def~\ref{ch21:def:spinTopVec}, Thm~\ref{ch21:thm:twistedTop}\\
\leanpath{MoonshineCNGriessSpinDict.lean} & individually & Def~\ref{ch21:def:spinDict}, Thm~\ref{ch21:thm:spinDict}\\
\leanpath{MoonshineCNGriessIso.lean} & individually & \S\ref{ch21:sec:iso}\\
\leanpath{MoonshineCNGriessAudit.lean} & audit & \S\ref{ch21:sec:audits}\\
\leanpath{MoonshineCNGriessAudit2.lean} & audit & \S\ref{ch21:sec:audits}\\
\leanpath{MoonshineCNGriessAudit3.lean} & audit & \S\ref{ch21:sec:audits}\\
\leanpath{MoonshineCNMonsterActionConj.lean} & individually & Def~\ref{ch21:def:griessConj}, Thm~\ref{ch21:thm:griessFormNat}\\
\leanpath{MoonshineCNMonsterActionLift.lean} & individually & Def~\ref{ch21:def:liftable}, Thms~\ref{ch21:thm:monsterLift}, \ref{ch21:thm:rational}\\
\leanpath{MoonshineCNMonsterActionExtra.lean} & individually & Thm~\ref{ch21:thm:extra}\\
\leanpath{MoonshineCNMonsterActionMoonshine.lean} & individually & Thm~\ref{ch21:thm:cnLiftable}\\
\leanpath{MoonshineCNMonsterActionAudit.lean} & audit & \S\ref{ch21:sec:audits}\\
\leanpath{MoonshineLRGriessTransfer.lean} & individually & \S\ref{ch21:sec:transfer}\\
\bottomrule
\end{longtable}

%======================================================================
\section{Audit files}\label{ch21:sec:audits}

Each audit file imports the files it audits and contains one \texttt{\#print axioms} command per public
declaration; the expected output of every line is \lean{[propext, Classical.choice, Quot.sound]} (or a
subset), with no \lean{sorryAx}.  The counts below are the numbers of \texttt{\#print axioms} lines in the
file at the release.
\begin{itemize}
\item \leanpath{VOA/\allowbreak Moonshine/\allowbreak MonsterGriess/\allowbreak ParityAudit.lean} (795): the parity-centralizer port
  (\leanpath{MinimalVectors} through \leanpath{ParityEqualsVisible}) and the triality prerequisites
  \leanpath{LatticeFamilies}, \leanpath{OctadEvenSubset}.
\item \leanpath{VOA/\allowbreak Moonshine/\allowbreak MonsterGriess/\allowbreak TrialityAudit.lean} (372): \leanpath{TrialityLabels} through
  \leanpath{TrialityFourGroup}, with \lean{tauB}, \lean{yB}, \lean{zB}, \lean{Generated.Gmat}.
\item \leanpath{VOA/\allowbreak Moonshine/\allowbreak MonsterGriess/\allowbreak ProductInvarianceAudit.lean} (659): \leanpath{GriessGrading},
  the \leanpath{Invariance*} files, \leanpath{OctadThreeOctad}, \leanpath{ProductInvariance}; headlines
  \lean{tauB_gmul}, \lean{isAlgAut_tauB}, \lean{tauB_mem_AutO}, \lean{Gmat_le_AutO}.
\item \leanpath{VOA/\allowbreak Moonshine/\allowbreak MoonshineCNGriessAudit.lean} (176): the table, the Conway coordinates, the
  untwisted dictionary, the short-class combinatorics.
\item \leanpath{VOA/\allowbreak Moonshine/\allowbreak MoonshineCNGriessAudit2.lean} (262): the ported cocycle and
  $\Leech/2\Leech$ identities, the sign map, \lean{untwistedIsoC}, the spin module, the frame blocks.
\item \leanpath{VOA/\allowbreak Moonshine/\allowbreak MoonshineCNGriessAudit3.lean} (154): \lean{SpinMatrixMap}, the spin form,
  the twisted top, the $W$-dictionary and \lean{griessIso}.
\item \leanpath{VOA/\allowbreak Moonshine/\allowbreak MoonshineCNMonsterActionAudit.lean} (84): the lift framework.
\end{itemize}
\leanpath{MoonshineLRGriessTransfer.lean} is audited in \leanpath{MoonshineLRLiftAudit.lean} (Chapter~\ref{ch:22}).
The axiom lists stated in the status boxes are the recorded results of building these audit files; the
author of this chapter did not re-run any build, and at the time of those records the aggregate build of
the whole library had not been re-run.

%======================================================================
\section{Referee's checklist}\label{ch21:sec:checklist}

\begin{longtable}{p{0.27\textwidth}p{0.33\textwidth}p{0.2\textwidth}p{0.1\textwidth}}
\toprule
declaration & intended meaning & hypotheses & check \\ \midrule
\endhead
\lean{Code.C} & AxiomMath's Golay code (MOG/hexacode) & — (definition) & Def~\ref{ch21:def:C}\\
\lean{golay_map_eq_C} & our Golay code is AxiomMath's up to $\pi$ & none & Thm~\ref{ch21:thm:golaymatch}\\
\lean{leechEquiv}, \lean{codeForm_leechEquiv} & our Leech lattice is AxiomMath's, isometrically & none & Thm~\ref{ch21:thm:leechmatch}\\
\lean{classOf_eq_iff}, \lean{card_shortIdx} & minimal classes are $\pm$ pairs; $98280$ of them & none & Thm~\ref{ch21:thm:minvecs}\\
\lean{phi0}, \lean{ShortIdx}, \lean{uvec} & coordinates on $\Leech/2\Leech$; short classes & — (definitions) & Def~\ref{ch21:def:short}\\
\lean{phi0Equiv}, \lean{qL_phi0} & $\varphi_0$ isomorphism; $q=$ pairing & none & Thm~\ref{ch21:thm:phi0}\\
\lean{qmatR}, \lean{qmatR_mul}, \lean{Q} & spin matrices; extraspecial $2^{1+24}$ & — / none & Def~\ref{ch21:def:qmatR}, Thm~\ref{ch21:thm:qmatRmul}, Def~\ref{ch21:def:Q}\\
\lean{SpaceB}, \lean{bform}, \lean{gmul} & AxiomMath's rational Griess algebra & — (definitions) & Defs~\ref{ch21:def:spaceB}, \ref{ch21:def:gmul}\\
\lean{finrank_B} & $\dim\chXXIB=196884$ & none & \S\ref{ch21:sec:spaceB}\\
\lean{T_unique} & Griess trilinear form determined by block values & symmetry, block values & Thm~\ref{ch21:thm:Tunique}\\
\lean{AutO} & AxiomMath's Monster: $\Aut(\chXXIB,\cdot,b)$ & — (definition) & Def~\ref{ch21:def:AutO}\\
\lean{Pvis}, \lean{Gvis}, \lean{thetaEquiv} & $2^{1+24}.\mathrm{Co}_1$ acting on $\chXXIB$ & — (definitions) & Defs~\ref{ch21:def:Pvis}, \ref{ch21:def:Gvis}\\
\lean{Gvis_le_AutO}, \lean{rhoEquiv_zQ} & $G_{\mathrm{vis}}\le$ AutO; $z\mapsto x_B$ & none & Thm~\ref{ch21:thm:gvis}\\
\lean{tbasis}, \lean{tauB}, \lean{trace_tauB} & triality basis; order-3 element, trace $783$ & none & Defs~\ref{ch21:def:tbasis}, \ref{ch21:def:tauB}\\
\lean{yB}, \lean{zB}, \lean{Generated.Gmat} & four-group; $G_{\mathrm{mat}}=\langle G_{\mathrm{vis}},\tau_B\rangle$ & — & Def~\ref{ch21:def:Gmat}, Thm~\ref{ch21:thm:fourgroup}\\
\lean{tauB_gmul}, \lean{tauB_mem_AutO}, \lean{Gmat_le_AutO} & $\tau_B$ is an automorphism; $G_{\mathrm{mat}}\le$ AutO & none & Thm~\ref{ch21:thm:tauBgmul}\\
\lean{rational_parity_centralizer_equals_visible}, \lean{AutO_inf_centralizer_xB_eq_Gvis} & AutO${}\cap C(x_B)=G_{\mathrm{vis}}$ & none & Thm~\ref{ch21:thm:parity}\\
(not ported) \lean{gmat_eq_griess} & $G_{\mathrm{mat}}=$ AutO & — & \S\ref{ch21:sec:notported}; named hypothesis in Ch.~\ref{ch:22}\\
\lean{conwayQuad_one_conwayQuad_of_symm} etc. & untwisted products in Conway coordinates & displayed & Thm~\ref{ch21:thm:conway}\\
\lean{ShortSignDictionary} & coherent phases for $\cos(\alpha X)$ & — (structure) & Def~\ref{ch21:def:ssd}\\
\lean{untwistedDict_gmul}, \lean{untwistedDict_form} & untwisted dictionary is isometric algebra map & sign dictionary, $W=0$ & Thm~\ref{ch21:thm:untwisted}\\
\lean{exists_sign_refinement} & quadratic refinements exist & symmetric form & Thm~\ref{ch21:thm:signref}\\
\lean{shortSignDictionary_griessShortSign}, \lean{untwistedIsoC} & untwisted dictionary unconditional; $\C\otimes(\chXXISym\oplus X)\cong(V_\Leech^+)_2$ & none & Thm~\ref{ch21:thm:griessShortSign}\\
\lean{spinIso}, \lean{spinIso_spinOp} & AxiomMath's spin module is $T$ & none & Thm~\ref{ch21:thm:spinIso}\\
\lean{TForm_spinIso} & $B_T=\kappa\cdot$ dot product & none & Thm~\ref{ch21:thm:TForm}\\
\lean{spinIso_frameJ_of_support} & spin halves $=$ frame sectors & support & Thm~\ref{ch21:thm:frame}\\
\lean{topAct_conwayQuad_spinTopVec} etc. & untwisted$\times$twisted channels & displayed & Thm~\ref{ch21:thm:twistedTop}\\
\lean{fullShortSign}, \lean{spinDict}, \lean{TPForm_spinDict} & phases $\zeta'$; $W$-dictionary; form & none & Def~\ref{ch21:def:spinDict}, Thm~\ref{ch21:thm:spinDict}\\
\lean{griessNat_gmul}, \lean{griessNat_form} & dictionary multiplicative, isometric up to $2$ & none & Thm~\ref{ch21:thm:griessNat}\\
\lean{griessIso}, \lean{griessIso_mul}, \lean{griessIso_form} & $\C\otimes\chXXIB\cong\Vnat_2$ as algebras with forms & \textbf{none} & Thm~\ref{ch21:thm:griessIso}\\
\lean{spinDict_mem_TPframe_shift} & $W$ halves $\to$ frame sectors & none & Thm~\ref{ch21:thm:blocks}\\
\lean{griessConj}, \lean{griessConjAut} & AutO acts on the Griess algebra of $\Vnat$ & none & Def~\ref{ch21:def:griessConj}, Thm~\ref{ch21:thm:griessFormNat}\\
\lean{griessForm_eq_natForm} & intrinsic form $=$ \lean{natForm} on $\Vnat_2$ & none & Thm~\ref{ch21:thm:griessFormNat}\\
\lean{IsVnatLiftable}, \lean{liftableSubgroup} & elements extending to $\Aut(\Vnat)$ & — (definitions) & Def~\ref{ch21:def:liftable}\\
\lean{monsterLift}, \lean{monsterLift_injective}, \lean{monsterLift_unique} & unique injective lift & none & Thm~\ref{ch21:thm:monsterLift}\\
\lean{isVnatLiftable_of_rational}, \lean{autVnat_mulEquiv_autO_of_liftable_of_rational} & rational automorphisms come from AutO & \lean{hrat}; \lean{hlift}, \lean{hrat} & Thm~\ref{ch21:thm:rational}\\
\lean{exists_monsterLift_eq_extraVnatAut} & extraspecial automorphisms lift, explicitly & none & Thm~\ref{ch21:thm:extra}\\
\lean{conwayNorton_monster_of_isVnatLiftable} & Conway--Norton for a liftable Monster element & $g$ liftable & Thm~\ref{ch21:thm:cnLiftable}\\
\lean{conwayNorton_autO_of_generators} & Conway--Norton for AutO from liftable generators & \lean{hvis}, \lean{hτ}, \lean{hGmat} & Thm~\ref{ch21:thm:cnLiftable}\\
\lean{griessTrX}, \lean{griessTrNat}, \lean{griessTrSigmaBlockOFull_nat} & $\zeta\leftrightarrow\zeta'$ transfer & displayed & Thm~\ref{ch21:thm:transfer}\\
\bottomrule
\end{longtable}

Every row is \texttt{\#print axioms}-clean (standard axioms) per the audit files of \S\ref{ch21:sec:audits};
the only hypotheses are those displayed, and none of them is an unestablished datum except
\lean{hlift}, \lean{hGmat} and \lean{hτ}/\lean{hvis} in the two general theorems of Part~C, which
Chapter~\ref{ch:22} reduces to the single named hypothesis \lean{GmatEqAutO}.

%======================================================================
\section{Notes and references}\label{ch21:sec:notes}

The Monster's Griess algebra was constructed by R.~L.~Griess Jr., \emph{The friendly giant}, Invent.\ Math.\
69 (1982) 1--102, building on his announcement in \emph{A construction of $F_1$ as automorphisms of a
196,883-dimensional algebra}, Proc.\ Natl.\ Acad.\ Sci.\ USA 78 (1981) 689--691.  The simplified construction
followed by AxiomMath (the decomposition $S^2\oplus\Q^{98280}\oplus\Q^{4096}\otimes\Q^{24}$, the extraspecial
group $2^{1+24}$ realised by signed permutations indexed by the Parker loop, and the triality element) is
J.~H.~Conway, \emph{A simple construction for the Fischer--Griess monster group}, Invent.\ Math.\ 79 (1985)
513--540.  The theorem that the full automorphism group of the Griess algebra is the Monster is J.~Tits,
\emph{On R.~Griess' ``friendly giant''}, Invent.\ Math.\ 78 (1984) 491--499; see also J.~Tits, \emph{Le Monstre},
S\'eminaire Bourbaki 1983/84, exp.\ 620, Ast\'erisque 121--122 (1985) 105--122.  The MOG, the hexacode and the
Golay code are in R.~T.~Curtis, \emph{A new combinatorial approach to $M_{24}$}, Math.\ Proc.\ Cambridge
Philos.\ Soc.\ 79 (1976) 25--42, and J.~H.~Conway and N.~J.~A.~Sloane, \emph{Sphere Packings, Lattices and
Groups}, 3rd ed., Springer, 1999, Chapters~10, 11 and~29.  The moonshine module, the identification of
$\Vnat_2$ with the Griess algebra, the triality and the action of the Monster on $\Vnat$ are in I.~Frenkel,
J.~Lepowsky and A.~Meurman, \emph{Vertex Operator Algebras and the Monster}, Pure and Applied Mathematics 134,
Academic Press, 1988 (Chapters~10--12 for the Griess algebra and $C=2^{1+24}.\mathrm{Co}_1$).  A modern
account of the Monster through its Griess algebra and its axes is A.~A.~Ivanov, \emph{The Monster Group and
Majorana Involutions}, Cambridge Tracts in Mathematics 176, Cambridge University Press, 2009.  The physics
of the orbifold is L.~Dixon, P.~Ginsparg and J.~Harvey, \emph{Beauty and the beast: superconformal symmetry in
a Monster module}, Comm.\ Math.\ Phys.\ 119 (1988) 221--241.  The conjecture is J.~H.~Conway and S.~P.~Norton,
\emph{Monstrous moonshine}, Bull.\ London Math.\ Soc.\ 11 (1979) 308--339, proved by R.~E.~Borcherds,
\emph{Monstrous moonshine and monstrous Lie superalgebras}, Invent.\ Math.\ 109 (1992) 405--444.

The Lean source of Part~A is the AxiomMath CFSG library: D.~Cureton, V.~Ilin, B.~Mehta, N.~Voss, R.~Cabral,
C.~Defant, E.~Gregory, A.~Hill, L.~Hong, Q.~Hu, D.~Jamner, A.~Jiang, T.~Jordan, A.~Kurghinyan, K.~Lau,
J.~Laurent, D.~Liang, J.~Limperg, S.~Mahns, B.~Nugent, C.~Olson, K.~Ono, E.~Paul, D.~Rajan, K.~Ram, C.~Ray,
H.~C.~Siu, A.~A.~Swaminathan, J.~Thorner, A.~Tsaplin and C.~Wang, \emph{Lean formalization of the
classification of finite simple groups}, 2026, \texttt{https://github.com/AxiomMath/CFSG}, commit
\texttt{69cb972f587a3ecc035e092253167ff2e303c962} (Lean v4.35.0-rc3), Apache License 2.0.  The port to
Lean~4.30 is by F.~K.~Popov (with Claude).  The record of the port is
\leanpath{Notes/\allowbreak monster\_aut/\allowbreak CFSG\_PORT\_FIDELITY.md} (the declaration-by-declaration comparison); the
remaining working notes are not part of the release.
